\documentclass[11pt]{ctexart}
\usepackage[a4paper,margin=1in]{geometry}
\usepackage{graphicx}
\usepackage{xcolor}
\usepackage{hyperref}
\definecolor{refgray}{gray}{0.45}
\title{\textbf{市场最新观点汇总}\\[0.4em]{\large 全球市场在AI结构性需求与政策再定价中分化，亚洲通胀与资本流动成为焦点}}
\author{KC桌面 自动生成}
\date{260611}
\begin{document}
\maketitle
\tableofcontents
\newpage
\section{一页摘要}
\begin{itemize}
  \item 中国5月出口超预期增长19.4\%，但驱动力从数量转向价格，AI相关产品贡献近一半增量，市场低估了供给优势的再定价而非需求弹性。
  \item 亚洲通胀反弹正在重塑央行利率路径，DB预计2026年亚洲CPI通胀率将是去年的三倍，市场尚未完全定价这一轮紧缩的力度与广度。
  \item 美国科技股遭遇机构客户创纪录抛售，BofA数据显示单周个股净流出142亿美元，资金正从科技+增长单极化配置转向分散化。
  \item 日本央行6月加息概率被市场低估，GS认为加息后的沟通路径比时点更重要，日本正进入加息但不紧缩的微妙阶段。
  \item 中国2万亿AI投资计划的核心在于80\%国产化率的硬约束，Citi认为这将系统性重塑国内半导体供应链的定价权。
  \item 欧洲半导体公司估值重估刚刚开始，GS上调英飞凌和意法半导体目标价，市场低估了AI对功率半导体的需求弹性。
  \item 亚洲动量交易已进入估值泡沫区，Bernstein建议削减台湾和韩国的动量敞口，转向高股息或低波动资产。
  \item 美元兑日元的真正支撑来自日本股市催生的对冲需求，Citi维持看空观点，认为160是顶部区间。
  \item 中国新能源车渗透率突破60\%，但竞争已从谁能造出电动车进入规模优势转化为定价权阶段，品牌分化加速。
  \item 全球钢铁市场拐点不在需求而在欧洲供给侧的重新定价，CBAM和保障措施正在重塑区域定价权。
\end{itemize}
\section{宏观与利率：亚洲通胀反弹与央行紧缩路径}
\textbf{亚洲通胀反弹正在重塑央行的利率路径，市场尚未完全定价这一轮紧缩的力度与广度。}

\begin{itemize}
  \item DB预计2026年亚洲CPI通胀将达到2.7\%，是去年的三倍，驱动力来自芯片价格飙升、能源成本冲击和中国再通胀。
  \item 印尼和菲律宾可能再加75个基点，韩国和印度可能加100个基点，马来西亚的加息也可能提前。
  \item 中国通胀呈现双轨制，PPI高企主要由能源和AI材料驱动，剔除后PPI可能仍是负值，CPI依赖能源基数效应而非消费需求回暖。
  \item NOM上调全年CPI预测从0.6\%至0.9\%，PPI从1.0\%至2.5\%，主要因油价高位和AI需求推升芯片价格。
  \item 日本央行6月加息可能性被市场低估，GS将加息预期从7月提前至6月，预计后续约半年一次加息，最终利率目标1.5\%。
  \item GS认为加息后的沟通路径比时点更重要，央行如何修改前瞻指引和管理接近中性利率的预期将决定资产定价走向。
  \item 日本经济呈现软数据与硬数据背离，消费者信心走弱但机械订单和消费活动指数稳健，为加息提供窗口。
  \item 中国金融体系正进入更可持续的高质量发展正循环，MS认为净息差可能提前见底，贷款增长更理性。
\end{itemize}
\begin{figure}[htbp]
\centering
\includegraphics[width=0.88\linewidth]{figures/fig_086.jpg}
\caption{Figure 1 - Figure 1: Slowing but above-trend growth and sharp rebound in inflation}
\end{figure}
\begin{figure}[htbp]
\centering
\includegraphics[width=0.88\linewidth]{figures/fig_087.jpg}
\caption{Figure 2 - Figure 2: Weaker growth momentum in April}
\end{figure}
\begin{figure}[htbp]
\centering
\includegraphics[width=0.88\linewidth]{figures/fig_088.jpg}
\caption{Figure 3 - Figure 3: Growth slows, inflation rebounds}
\end{figure}
\begin{figure}[htbp]
\centering
\includegraphics[width=0.88\linewidth]{figures/fig_189.jpg}
\caption{Exhibit 1 - Exhibit 1: Confidence is Deteriorating... Consumer Confidence and Economy Watchers Survey DI (Household Activity-Related)}
\end{figure}
\begin{figure}[htbp]
\centering
\includegraphics[width=0.88\linewidth]{figures/fig_190.jpg}
\caption{Exhibit 2 - Exhibit 2: ...But Hard Data Remains Solid Machinery Orders (Private Sector Excluding Volatile Orders) and BOJ Consumption Activity Index}
\end{figure}
{\small\color{refgray}\textbf{References:} [R001] NOM：中国通胀的结构性真相——市场低估了供给驱动的持续性，高估了需求复苏的弹性; [R008] DB：亚洲通胀反弹正在重塑央行的利率路径，市场尚未完全定价这一轮紧缩的力度与广度; [R011] GS：日本央行6月加息的可能性被市场低估，真正关键的是加息后的沟通路径; [R022] 摩根斯坦利：中国金融体系真正的转折点不是刺激，而是“高质量”的内生循环}

\section{AI/算力：资本开支上行风险与供应链结构性重置}
\textbf{市场真正低估的不是AI需求，而是2027年资本开支的上行风险以及供电模式的结构性重置。}

\begin{itemize}
  \item GS指出共识预测2027年超大规模云厂商资本开支仅9200亿美元、增速骤降至22\%，但过去三年分析师平均低估了45个百分点。
  \item 如果增量投资达到GDP的2\%至3\%，2027年资本开支将落在1.1万亿至1.25万亿美元区间，激进情景下可达1.43万亿美元。
  \item 中国计划未来五年投入2万亿元建设数据中心，要求至少80\%使用国产AI加速器，Citi认为这一目标可实现。
  \item Citi估算2万亿投资相当于新增约10GW的AI数据中心容量，年均2GW，将使AI级容量在现有基础上扩张约50\%。
  \item Bernstein数据显示北美数据中心项目管道总量增至324GW，自备发电模式在在建容量中占25\%，管道中高达40\%。
  \item 自备发电项目在电气分配设备上的支出比纯电网接入项目高出约30\%，Bernstein认为市场对此定价不充分。
  \item MS在Computex供应链核查显示800V DC电源机架按计划2026年Q3量产，多家超大规模客户已转向800V DC。
  \item CPO市场真正的分歧在量产节奏，MS预计2027年光学引擎出货量仅600-700万颗，远低于市场期待的2000-3000万颗。
\end{itemize}
\begin{figure}[htbp]
\centering
\includegraphics[width=0.88\linewidth]{figures/fig_436.jpg}
\caption{Exhibit 4 - Exhibit 4: Consensus expects the hyperscalers will spend \textbackslash{}\$757 billion in 2026 and \textbackslash{}\$920 billion in 2027}
\end{figure}
\begin{figure}[htbp]
\centering
\includegraphics[width=0.88\linewidth]{figures/fig_437.jpg}
\caption{Exhibit 5 - Exhibit 5: Hyperscalers expected to allocate 100\% of cash flows from operations to capex}
\end{figure}
\begin{figure}[htbp]
\centering
\includegraphics[width=0.88\linewidth]{figures/fig_440.jpg}
\caption{Exhibit 8 - Exhibit 8: AI hyperscaler 2027 capex scenarios 2027 AI hyperscaler capex scenarios (\$ billion)}
\end{figure}
\begin{figure}[htbp]
\centering
\includegraphics[width=0.88\linewidth]{figures/fig_339.jpg}
\caption{EXHIBIT 9 - EXHIBIT 9: Behind The Meter: Only \$4\textbackslash{}\%\$ of the Installed Base, But \$25\textbackslash{}\%\$ of What's Under Construction, and \$40\textbackslash{}\%\$ of The Project Pipeline Behind-the-Meter Mix of Project by Stage}
\end{figure}
\begin{figure}[htbp]
\centering
\includegraphics[width=0.88\linewidth]{figures/fig_388.jpg}
\caption{EXHIBIT 59 - EXHIBIT 59: Behind The Meter: Only 4\% of the Installed Base, But 25\% of What's Under Construction, and 40\% of The Project Pipeline Behind-the-Meter Mix of Project by Stage}
\end{figure}
\begin{figure}[htbp]
\centering
\includegraphics[width=0.88\linewidth]{figures/fig_266.jpg}
\caption{Exhibit 1 - Exhibit 1: 800 VDC equipment readiness 800 VDC Equipment Readiness Q3 2026 Option 1A – Power Rack (660 kW) Q2 2027 Option 1B – DC Power Center (1.6 MW) Q1 2028 Option 1C – Power Block}
\end{figure}
{\small\color{refgray}\textbf{References:} [R014] Citi：中国2万亿AI投资计划的核心不在于规模，而在于国产化率的硬约束; [R018] 摩根斯坦利：市场对800V DC延迟的担忧被夸大了，真正的信号是供应链加速而非推迟; [R021] Bernstein：数据中心真正的拐点不在建设量，而在供电模式的结构性重置; [R026] 摩根斯坦利：CPO市场真正的分歧不在技术，而在量产节奏的预期差; [R029] GS：市场真正低估的不是AI需求，而是2027年资本开支的上行风险; [R032] Citi：中国2万亿AI基建计划真正重塑的不是需求，而是国内供应链的定价权}

\section{能源与大宗：欧洲钢铁定价权重塑与储能政策重心转移}
\textbf{全球钢铁市场拐点不在需求而在欧洲供给侧的重新定价，储能正在成为全球能源政策的新重心。}

\begin{itemize}
  \item MS指出欧洲通过CBAM和保障措施系统性地重塑钢铁定价权，CBAM于2026年1月生效，保障措施7月收紧。
  \item 欧洲结构性短缺1000-1500万吨，欧盟热卷利润已升至404美元/吨，高于长期均值320美元/吨。
  \item 中国5月成品钢净出口989万吨，环比增长约10\%，但增量正被全球贸易区域化趋势抵消。
  \item HSBC在SNEC观展报告中指出地面光伏面临中国市场化定价改革和跨界竞争压力，利润率压缩可能是结构性的。
  \item 储能尤其是户用储能正在成为全球能源政策新重心，中国锂电池出口前十大市场前4个月同比增长29\%。
  \item 欧盟4月22日禁止用中国逆变器的项目拿欧盟资金，预计影响欧洲约20\%的公用事业级光伏+储能市场。
  \item HSBC观察到固态变压器作为连接光伏直流发电与数据中心配电的关键环节，成为多个企业竞相布局的焦点。
  \item 中国铝出口跳升至7个月新高，MS指出驱动力是套利窗口而非需求信号，可持续性取决于内外价差动态。
\end{itemize}
\begin{figure}[htbp]
\centering
\includegraphics[width=0.88\linewidth]{figures/fig_811.jpg}
\caption{Exhibit 1 - Exhibit 1: China's finished steel net exports increased 10\% MoM (day-adjusted) in May Steel net exports from China}
\end{figure}
\begin{figure}[htbp]
\centering
\includegraphics[width=0.88\linewidth]{figures/fig_814.jpg}
\caption{Exhibit5 - Exhibit5: Price spreads: EU HRC less China HRC}
\end{figure}
\begin{figure}[htbp]
\centering
\includegraphics[width=0.88\linewidth]{figures/fig_815.jpg}
\caption{Exhibit6 - Exhibit6: Spreads: EU HRC gross profit per tonne\textbackslash{}*}
\end{figure}
\begin{figure}[htbp]
\centering
\includegraphics[width=0.88\linewidth]{figures/fig_811.jpg}
\caption{Exhibit 1 - Exhibit 1: China's finished steel net exports increased 10\% MoM (day-adjusted) in May Steel net exports from China}
\end{figure}
{\small\color{refgray}\textbf{References:} [R010] 摩根斯坦利：中国大宗商品出口的“结构性分化”正在重塑全球定价权; [R037] HSBC：储能正在成为全球能源政策的新重心，光伏反而面临需求拐点; [R038] HSBC：SNEC 2026揭示的真正信号，不是光伏回暖，而是能源系统的权力交接; [R044] 摩根斯坦利：全球钢铁市场的真正拐点不在需求，而在欧洲供给侧的重新定价}

\section{地产与消费：中国餐饮品牌分化加速与家电估值逻辑切换}
\textbf{中国消费领域品牌分化加速，餐饮和家电龙头正从周期股转向复合增长引擎，市场定价仍停留在旧框架。}

\begin{itemize}
  \item GS五月餐饮追踪显示同店销售增长较四月走弱，但品牌分化加速，现制茶饮高基数效应开始显现。
  \item 奈雪的茶五月单店销售同比下滑从-11\%扩大至-24\%，而头部品牌展现出更强韧性，支撑来自品类扩张和新产品。
  \item 海底捞五月翻台率同比转负，对应恢复至2019年同期的80-85\%水平，管理团队认为端午假期移位是扰动因素。
  \item JPM提出中国家电三巨头不应再按国内家电周期估值，而是作为现金牛盈利+全球挑战者增长+工业科技期权价值的组合。
  \item 三巨头国内市占率约60\%，海外仅16\%，而海外可触达市场是国内的3倍，JPM预计2026-28年海外收入CAGR达7-11\%。
  \item B2B业务如商用暖通空调和数据中心液冷是估值提升的钥匙，美的被JPM选为最清晰的标的。
  \item GS医疗设备招标数据显示五月总招标金额同比微增0.1\%，结束连续五个月下滑，CT、MRI和直线加速器实现正增长。
  \item 联影医疗前5个月招标金额同比+9\%，是主要竞争对手中唯一正增长，MS认为下半年有望看到更明显复苏。
\end{itemize}
\begin{figure}[htbp]
\centering
\includegraphics[width=0.88\linewidth]{figures/fig_259.jpg}
\caption{Exhibit 3 - Exhibit 5: May growth recovery level remained relatively soft}
\end{figure}
\begin{figure}[htbp]
\centering
\includegraphics[width=0.88\linewidth]{figures/fig_260.jpg}
\caption{Exhibit 4 - Exhibit 4: DAU/user session of major delivery/food service brand APP As of May 2026. Exhibit 5: May growth recovery level remained relatively soft}
\end{figure}
\begin{figure}[htbp]
\centering
\includegraphics[width=0.88\linewidth]{figures/fig_538.jpg}
\caption{Exhibit 1 - Exhibit 1: Total bidding value of 9 main medical devices in China In Rmb, mn}
\end{figure}
\begin{figure}[htbp]
\centering
\includegraphics[width=0.88\linewidth]{figures/fig_539.jpg}
\caption{Exhibit 2 - Exhibit 2: The YoY change rate of bidding value}
\end{figure}
\begin{figure}[htbp]
\centering
\includegraphics[width=0.88\linewidth]{figures/fig_567.jpg}
\caption{Figure 1 - Figure 1: Home appliance sector stocks under JPM coverage - share price performance vs China/HK index Indexed to Jun 5, 2026 = 100}
\end{figure}
\begin{figure}[htbp]
\centering
\includegraphics[width=0.88\linewidth]{figures/fig_568.jpg}
\caption{Figure 3 - Figure 3: Global (ex-China) home appliance market size\textbackslash{}* and China Big 3's market share \textbackslash{}* Market size is based on ex-factory prices}
\end{figure}
{\small\color{refgray}\textbf{References:} [R017] GS：中国餐饮的五月数据揭示的不是需求疲软，而是品牌分化的加速; [R036] GS：中国医疗设备招标的拐点正在到来，但真正的分化才刚刚开始; [R039] JPM：中国家电的估值逻辑正在从“周期股”转向“复合增长引擎”}

\section{区域市场：印度制造业竞争力缺口与韩元贬值悖论}
\textbf{印度面临的是制造业竞争力的结构性缺口而非短期资本流入问题，韩元走弱源于资本市场繁荣催生的资本外流悖论。}

\begin{itemize}
  \item MS指出印度2024年12月至2026年3月国际收支累计赤字520亿美元，是2011-2013年期间的两倍以上。
  \item 印度BSE500指数成分股2026年Q1盈利同比增长13\%，远低于韩国171\%、台湾49\%、日本33\%和美国29\%。
  \item MS认为印度需要提升全价值链制造能力，瞄准高增长行业，整合供应链，加大研发投入。
  \item 卢比实际有效汇率已跌至低于10年均值3.7个标准差的位置，显示系统性失衡。
  \item NOM分析韩元走弱根源在于经常账户盈余与资本账户结构性失血并存的悖论。
  \item 韩国企业主动持有美元资产而非结汇，零售投资者对海外股票配置热情高涨，NPS的国内股票配置目标调整可能形成新抛压。
  \item 三星电子和SK海力士若再涨60\%，全球基金因单只股票持仓不得超过10\%的规则可能带来高达480亿美元的外资流出。
  \item Bernstein指出印度350cc摩托车市场真正的较量不是产品而是分销，Triumph尚未在Royal Enfield的文化领地上建立足够的分销网络。
\end{itemize}
\begin{figure}[htbp]
\centering
\includegraphics[width=0.88\linewidth]{figures/fig_005.jpg}
\caption{Exhibit 1 - Exhibit 1: INR on a REER basis has now fallen to more than 3.7 standard deviations away from the 10-year mean trend}
\end{figure}
\begin{figure}[htbp]
\centering
\includegraphics[width=0.88\linewidth]{figures/fig_006.jpg}
\caption{Exhibit 2 - Exhibit 2: The key driver has been weaker BOP dynamics}
\end{figure}
\begin{figure}[htbp]
\centering
\includegraphics[width=0.88\linewidth]{figures/fig_007.jpg}
\caption{Exhibit 3 - Exhibit 3: FII has recorded persistent outflows}
\end{figure}
\begin{figure}[htbp]
\centering
\includegraphics[width=0.88\linewidth]{figures/fig_012.jpg}
\caption{Exhibit 8 - Exhibit 8: India's 1Q26 earnings growth lagged other major indices more levered to the AI thematic EPS growth (\%Y, 1Q26)}
\end{figure}
\begin{figure}[htbp]
\centering
\includegraphics[width=0.88\linewidth]{figures/fig_188.jpg}
\caption{Figure 5 - \#\# Our key takeaways are as such: - Our estimates suggest fund allocations into Samsung Electronics and SK Hynix may have breached the 10\% limit around 13 May 2026 and 6 May 2026, respectively, broadly in line with when market discussions on the topic began. -}
\end{figure}
\begin{figure}[htbp]
\centering
\includegraphics[width=0.88\linewidth]{figures/fig_850.jpg}
\caption{EXHIBIT 2 - EXHIBIT 2: GST cuts announced in Aug'25}
\end{figure}
\begin{figure}[htbp]
\centering
\includegraphics[width=0.88\linewidth]{figures/fig_852.jpg}
\caption{EXHIBIT 7 - EXHIBIT 7: Distribution Gap by City tier}
\end{figure}
{\small\color{refgray}\textbf{References:} [R003] 摩根斯坦利：印度真正需要解决的，不是短期资本流入，而是制造业竞争力的结构性缺口; [R009] NOM：韩元走弱的真正推手不是经济基本面，而是资本市场成功的“悖论”; [R047] Bernstein：印度350cc摩托车市场真正的较量，不是产品，是分销}

\section{风险偏好：科技股信心结构性重置与亚洲动量泡沫}
\textbf{市场正经历科技股信心的结构性重置而非简单避险，亚洲动量交易已进入估值泡沫区。}

\begin{itemize}
  \item BofA报告显示机构客户以有数据记录以来最大规模抛售科技股，单周个股净流出142亿美元。
  \item 科技板块流出规模为BofA自2008年有数据以来最大，按市值占比算是2014年初以来最极端。
  \item 资金从大盘股完全流出，小盘和中盘股反而获得买入，客户买入价值型和平衡型ETF，卖出成长型ETF。
  \item Bernstein量化策略指出亚洲除日本外动量策略年内跑赢基准36\%-38\%，台湾和韩国动量超额收益超50\%。
  \item 台湾和韩国动量股估值创历史新高，盈利预期同样处于前所未有的高位，构成脆弱性极高的组合。
  \item 动量因子与成长因子的相关性已升至历史高位，低波动风格动量跌至25年最低，预示风格切换临界点。
  \item JPM零售投资者资金流显示散户在一周之内从极端买入科技股转向创纪录卖出，是亏损止损而非换仓。
  \item JPM指出日本股市最被低估的交易机会是同时具备价值特征和适度动量的股票，而非追高动量股。
\end{itemize}
\begin{figure}[htbp]
\centering
\includegraphics[width=0.88\linewidth]{figures/fig_267.jpg}
\caption{Exhibit 3 - Exhibit 3: Tech net flows as a \% of S\&P 500 Tech Market cap were the most negative since Feb. '14 Weekly Tech net flows as a \% of S\&P 500 Tech market cap, since 2008}
\end{figure}
\begin{figure}[htbp]
\centering
\includegraphics[width=0.88\linewidth]{figures/fig_268.jpg}
\caption{Exhibit 1 - Exhibit 1: Institutional clients are the biggest net sellers post-crisis Cumulative flows (\textbackslash{}\$ bn) by client type, February 2008-present}
\end{figure}
\begin{figure}[htbp]
\centering
\includegraphics[width=0.88\linewidth]{figures/fig_271.jpg}
\caption{Exhibit 8 - Exhibit 8: By sector, Tech stocks saw the largest outflows last week BofA client net buying (selling) by GICS sector (\textbackslash{}\$mn) – single stocks only}
\end{figure}
\begin{figure}[htbp]
\centering
\includegraphics[width=0.88\linewidth]{figures/fig_046.jpg}
\caption{Exhibit 1 - Exhibit 1-Exhibit 3) EXHIBIT 1: Asia ex Japan factor performance: YTD price momentum has led, up 36\%-38\% relative to market while low vol/high yield/value has underperformed the market the most Asia ex Japan 2026 YTD Factor Perfor}
\end{figure}
\begin{figure}[htbp]
\centering
\includegraphics[width=0.88\linewidth]{figures/fig_049.jpg}
\caption{Exhibit 4 - EXHIBIT 4: Asia Momentum 12m fwd. PE: Asia earnings momentum basket has been trading at record high valuations for a while. Price momentum in Asia is still not trading at record high valuations (at 1.3SD level), however in Asia ex}
\end{figure}
\begin{figure}[htbp]
\centering
\includegraphics[width=0.88\linewidth]{figures/fig_063.jpg}
\caption{Exhibit 17 - Exhibit 17-Exhibit 18) EXHIBIT 17: Momentum bubble is visible in both TW and KR where momentum stocks reached record high valuations in this cycle and already showing de-rating since the beginning of this month}
\end{figure}
\begin{figure}[htbp]
\centering
\includegraphics[width=0.88\linewidth]{figures/fig_617.jpg}
\caption{Figure 67 - Figure 1: Tech ETF Outflows at All-Time-Highs Weekly flows in \textbackslash{}\$B, as of Jun 10 \$\textasciicircum{}\{th\}\$}
\end{figure}
\begin{figure}[htbp]
\centering
\includegraphics[width=0.88\linewidth]{figures/fig_618.jpg}
\caption{Figure 2 - Figure 2: Retail P\&L in ETFs - Underperforming QQQ}
\end{figure}
{\small\color{refgray}\textbf{References:} [R006] Bernstein：亚洲动量交易已进入估值泡沫区，减仓窗口正在关闭; [R019] 美国银行：市场正在经历一场科技股信心的结构性重置，而非简单的避险; [R040] JPM：零售投资者在科技股上的踩踏式抛售，暴露的不仅是情绪，更是仓位结构的脆弱性; [R046] JPM：日本股市真正的博弈不在AI，而在价值与动量的交叉点}

\section{中国贸易：价格驱动的结构性重定价与AI主导的出口新格局}
\textbf{中国5月出口强劲但驱动力是价格而非数量，AI相关产品连续两个月贡献总出口增长近一半。}

\begin{itemize}
  \item 5月出口同比增长19.4\%，远超市场预期的15\%，但NOM和JPM均指出增长主要由价格效应驱动。
  \item NOM估算5月IC出口金额同比增长110.9\%，但出口数量仅增长2.1\%，价格贡献高达106.5个百分点。
  \item AI相关产品（IC+ADP）合计贡献9.3个百分点，占中国5月总出口增长的约48\%。
  \item BARC报告指出对美出口同比飙升35.4\%，部分得益于低基数，但区域贸易网络正在重新激活。
  \item JPM指出出口增长越来越窄、越来越依赖价格驱动、越来越与国内生产和就业脱钩。
  \item 存储芯片从量增驱动转变为价增驱动，新三样价格企稳同时伴随可观的量增。
  \item BARC认为中国出口驱动力正从传统成本优势转向AI资本开支与能源转型共同支撑的结构性需求。
  \item DB认为中国贸易结构性驱动力正从规模扩张转向以AI相关产品和能源密集型产品为代表的定价权优势。
\end{itemize}
\begin{figure}[htbp]
\centering
\includegraphics[width=0.88\linewidth]{figures/fig_001.jpg}
\caption{Figure 1 - Figure 1: China's trade momentum accelerated further}
\end{figure}
\begin{figure}[htbp]
\centering
\includegraphics[width=0.88\linewidth]{figures/fig_002.jpg}
\caption{Figure 2 - Figure 2: "AHEAD" factors remained the dominant drivers}
\end{figure}
\begin{figure}[htbp]
\centering
\includegraphics[width=0.88\linewidth]{figures/fig_003.jpg}
\caption{Figure 3 - Figure 3: Exports to major developed markets surged}
\end{figure}
\begin{figure}[htbp]
\centering
\includegraphics[width=0.88\linewidth]{figures/fig_426.jpg}
\caption{FIGURE 2 - FIGURE 2. ... with faster regional trade (Asean, Japan, Korea and Taiwan) and shipments to the US}
\end{figure}
\begin{figure}[htbp]
\centering
\includegraphics[width=0.88\linewidth]{figures/fig_427.jpg}
\caption{FIGURE 3 - FIGURE 3. Strong exports were led by AI-related and green-tech products}
\end{figure}
{\small\color{refgray}\textbf{References:} [R002] DB：中国贸易数据揭示的不是短期反弹，而是供给优势的再定价; [R020] NOM：中国5月出口强劲，但市场真正低估的不是需求，而是价格信号的重估; [R028] BARC：市场低估了中国出口的结构性韧性，AI与绿色技术正重塑贸易格局; [R030] JPM：市场低估的不是出口增速，而是贸易结构对增长含义的根本改变}

\section{中国新能源车：渗透率突破60\%后的竞争结构重构}
\textbf{中国新能源车渗透率突破60\%不是终点而是竞争烈度升级的起点，品牌分化加速且订单结构成为关键。}

\begin{itemize}
  \item DB数据显示五月新能源渗透率达到60.9\%，环比提升3个百分点，同比提升超过9个百分点。
  \item 五月新能源车批发134.7万辆，同比+12\%，环比+10.3\%，但乘用车批发总量同比下降4.9\%。
  \item DB周度订单追踪显示6月首周主要新能源车企合计订单环比下滑约5\%，同比下滑约15\%。
  \item 比亚迪周订单同比骤降36\%，订单中枢自3月峰值后持续下移，暴露规模领先者的结构性隐患。
  \item 蔚来集团周订单同比暴增419\%但环比下滑28\%，新品上市冲高后的正常回落。
  \item MS指出6月第一周订单回落是新品上市冲高后的正常化，真正的考验是订单能否在交付中得到兑现。
  \item NOM提出2026年国内汽车零售可能出现有史以来首次双位数同比下滑，前5个月累计零售同比下降19.3\%。
  \item 中国车企在欧洲市场4月销量突破16.3万辆，同比增长41.7\%，市场份额首次稳定在12\%。
\end{itemize}
\begin{figure}[htbp]
\centering
\includegraphics[width=0.88\linewidth]{figures/fig_470.jpg}
\caption{Figure 2 - Figure 2: Monthly total passenger-vehicle sales trend}
\end{figure}
\begin{figure}[htbp]
\centering
\includegraphics[width=0.88\linewidth]{figures/fig_471.jpg}
\caption{Figure 3 - Figure 3: Monthly new-energy vehicle (NEV) passenger-vehicle sales trend}
\end{figure}
\begin{figure}[htbp]
\centering
\includegraphics[width=0.88\linewidth]{figures/fig_473.jpg}
\caption{Figure 5 - Figure 5: Monthly NEV penetration rate trend}
\end{figure}
\begin{figure}[htbp]
\centering
\includegraphics[width=0.88\linewidth]{figures/fig_530.jpg}
\caption{Figure 2 - Figure 2: Weekly new orders trend of key automakers (Li, NIO, XPeng, Leap, Xiaomi, Tesla)}
\end{figure}
\begin{figure}[htbp]
\centering
\includegraphics[width=0.88\linewidth]{figures/fig_531.jpg}
\caption{Figure 2 - Figure 2: Weekly new orders trend of key automakers (Li, NIO, XPeng, Leap, Xiaomi, Tesla)}
\end{figure}
{\small\color{refgray}\textbf{References:} [R016] DB：中国车企在欧洲的渗透已从“试探”进入“阵地战”阶段; [R034] DB：中国车市正在经历的，不是周期波动，而是竞争结构的根本性重构; [R035] DB：中国新能源车市最值得关注的不是总量，而是订单结构的剧烈分化; [R043] 摩根斯坦利：中国电动车市场正进入“产品攻势间歇期”的定价博弈; [R045] NOM：中国汽车市场2026年可能首次出现国内需求双位数下滑}

\section{半导体：欧洲设备需求周期尚未见顶与存储供给纪律重塑}
\textbf{欧洲半导体设备需求周期尚未见顶，存储周期回调是延续的必要重置而非终结。}

\begin{itemize}
  \item GS上调ASML、ASMI、BESI和Nebius目标价，指出三大信号同时指向设备需求周期尚未见顶。
  \item SK海力士计划五年内翻倍晶圆产能，博通重申FY27 AI半导体收入将超1000亿美元。
  \item GS认为英飞凌和意法半导体正经历从周期性半导体公司向AI结构性增长公司的身份转换。
  \item 英飞凌当前两年远期收入CAGR约17\%，几乎是历史中位数9\%的两倍，支撑估值倍数扩张。
  \item Bernstein通过台湾设备进口数据分析，测试设备需求正在加速，光刻设备订单强劲。
  \item MS提出DRAM这轮回调不是周期的终结而是延续的必要重置，长期供应协议正在改变游戏规则。
  \item 如果LTA能覆盖未来三到五年70\%以上的总供给，DRAM股票估值中枢可从5倍PE向8-10倍PE迁移。
  \item JPM指出MLCC粉体行业已在2026年上半年进入强于预期的上升周期，高端产品线制造利润断层。
\end{itemize}
\begin{figure}[htbp]
\centering
\includegraphics[width=0.88\linewidth]{figures/fig_042.jpg}
\caption{Exhibit 3 - Exhibit 3: The expansion in Infineon's valuation multiple implied by our PT is congruent with the increase in 2Y FWD revenue CAGR versus history}
\end{figure}
\begin{figure}[htbp]
\centering
\includegraphics[width=0.88\linewidth]{figures/fig_044.jpg}
\caption{Exhibit 5 - Exhibit 5: The expansion in STMicro's valuation multiple in underpinned by an increase in 2Y FWD revenue CAGR versus history}
\end{figure}
\begin{figure}[htbp]
\centering
\includegraphics[width=0.88\linewidth]{figures/fig_456.jpg}
\caption{Exhibit 1 - EXHIBIT 1: May Taiwan import for SPE was \textbackslash{}\$4.4bn, -8\% MoM.}
\end{figure}
\begin{figure}[htbp]
\centering
\includegraphics[width=0.88\linewidth]{figures/fig_457.jpg}
\caption{EXHIBIT 2 - EXHIBIT 2: Likewise, May Taiwan SPE import from Japan was \textbackslash{}\$700mn, -17\% MoM.}
\end{figure}
\begin{figure}[htbp]
\centering
\includegraphics[width=0.88\linewidth]{figures/fig_458.jpg}
\caption{Exhibit 4 - EXHIBIT 3: May tester import from Japan and Malaysia collectively was \textbackslash{}\$613mn, +7\% MoM.}
\end{figure}
\begin{figure}[htbp]
\centering
\includegraphics[width=0.88\linewidth]{figures/fig_798.jpg}
\caption{Exhibit 1 - Exhibit 1: DRAM Stocks vs. 200-Day Moving Average}
\end{figure}
\begin{figure}[htbp]
\centering
\includegraphics[width=0.88\linewidth]{figures/fig_799.jpg}
\caption{Exhibit 2 - Exhibit 2: DRAM Stocks: Major sell-offs through the rally}
\end{figure}
{\small\color{refgray}\textbf{References:} [R005] GS：欧洲半导体公司的估值重估才刚刚开始，市场低估了AI对功率半导体的需求弹性; [R025] GS：欧洲半导体设备市场的下一个增长动力，不是周期，是结构; [R031] Bernstein：半导体设备进口数据揭示的不仅是周期性复苏，更是结构性分化; [R042] 摩根斯坦利：市场真正低估的不是需求，而是供给纪律的长期重塑; [R050] JPM：MLCC材料端正在经历一场被低估的定价权迁移}

\section{外汇与资本流动：日元对冲需求与韩元资本外流悖论}
\textbf{美元兑日元的真正支撑来自日本股市催生的对冲需求，韩元走弱源于资本市场繁荣催生的资本外流悖论。}

\begin{itemize}
  \item Citi双层模型显示日本股市和美元指数是USD/JPY最主要的驱动因素，利差收窄但汇率未跟随。
  \item Citi估算USD/JPY合理价位在161左右，实际汇率略低于此水平，可能因日本当局干预。
  \item Citi维持看空观点，认为160是顶部区间，年底前回落至155以下，但短期日元走弱风险仍存。
  \item NOM分析韩元走弱根源在于经常账户盈余与资本账户结构性失血并存的悖论。
  \item 韩国企业主动持有美元资产而非结汇，零售投资者对海外股票配置热情高涨。
  \item NPS的国内股票配置目标调整可能形成新抛压，三星电子和SK海力士的10\%持仓限制可能带来480亿美元外资流出。
  \item DB认为亚洲通胀反弹正在重塑央行的利率路径，市场尚未完全定价这一轮紧缩的力度与广度。
  \item GS对香港银行跨境监管风险的分析显示，即便极端假设下内地净新资金完全归零，对HSBC和渣打税前利润影响也仅约1\%和2\%。
\end{itemize}
\begin{figure}[htbp]
\centering
\includegraphics[width=0.88\linewidth]{figures/fig_072.jpg}
\caption{Figure 2 - Our two-tiered model (double decker model) analyzes the data from 2017 through 2025, and we add dummy variables for 2022 while also adding the cross terms (interaction terms) to the analysis (we estimated both the intercept dummy and the slope dummy for 2022).}
\end{figure}
\begin{figure}[htbp]
\centering
\includegraphics[width=0.88\linewidth]{figures/fig_073.jpg}
\caption{Figure 2 - © 2026 Citi Inc. No redistribution without Citi's written permission. Note: The deviation from the 200-day MA is based on the actual 200-day MA in the single deck model and on our 200-day MA estimate in the double decker model. Figure 2. USD/JPY: Double decker}
\end{figure}
\begin{figure}[htbp]
\centering
\includegraphics[width=0.88\linewidth]{figures/fig_078.jpg}
\caption{Figure 9 - © 2026 Citi Inc. No redistribution without Citi's written permission. Note: 200-day moving average (MA). Figure 9. USD/JPY and Japanese stock index (TOPIX)}
\end{figure}
\begin{figure}[htbp]
\centering
\includegraphics[width=0.88\linewidth]{figures/fig_188.jpg}
\caption{Figure 5 - \#\# Our key takeaways are as such: - Our estimates suggest fund allocations into Samsung Electronics and SK Hynix may have breached the 10\% limit around 13 May 2026 and 6 May 2026, respectively, broadly in line with when market discussions on the topic began. -}
\end{figure}
{\small\color{refgray}\textbf{References:} [R007] Citi：美元兑日元的真正支撑不是利差，而是日本股市催生的对冲需求; [R009] NOM：韩元走弱的真正推手不是经济基本面，而是资本市场成功的“悖论”; [R008] DB：亚洲通胀反弹正在重塑央行的利率路径，市场尚未完全定价这一轮紧缩的力度与广度; [R033] GS：市场对香港银行跨境监管风险的担忧，可能被放大了十倍}

\section{医疗健康：药价谈判结构性变化与设备招标拐点}
\textbf{美国医保体系正经历系统性财政责任再分配，中国医疗设备招标拐点到来但分化刚刚开始。}

\begin{itemize}
  \item GS指出CMS将最惠国机制视为更广泛的工具包，目标通过贸易动态在全球范围内施加定价压力。
  \item IRA谈判第一轮带来约20\%平均净节省，第二轮折扣幅度已接近40\%。
  \item GS认为医疗健康板块的结构性拐点比市场感知的更近，板块内部正在发生深刻分化。
  \item 2026年至今全球医疗健康并购交易额已接近1000亿美元，2025年全年仅为500亿美元。
  \item 中国医疗设备招标五月总金额同比微增0.1\%，结束连续五个月下滑，CT、MRI和直线加速器实现正增长。
  \item GS预计更为明显的复苏将在2026年下半年出现，但市场底部可能已经确认。
  \item Citi专家电话会指出BINSA法案通过概率低于市场隐含预期，核心障碍在立法流程而非政治意愿。
  \item Citi给出BINSA法案通过概率区间为30-40\%以下，财政部已明确表示现行法律下缺乏法定权限。
\end{itemize}
\begin{figure}[htbp]
\centering
\includegraphics[width=0.88\linewidth]{figures/fig_538.jpg}
\caption{Exhibit 1 - Exhibit 1: Total bidding value of 9 main medical devices in China In Rmb, mn}
\end{figure}
\begin{figure}[htbp]
\centering
\includegraphics[width=0.88\linewidth]{figures/fig_539.jpg}
\caption{Exhibit 2 - Exhibit 2: The YoY change rate of bidding value}
\end{figure}
\begin{figure}[htbp]
\centering
\includegraphics[width=0.88\linewidth]{figures/fig_547.jpg}
\caption{Exhibit 12 - Exhibit 12: CT scanner procurement value increased by +12\% yoy in May-26 vs. -31\% in Apr-26 CT}
\end{figure}
\begin{figure}[htbp]
\centering
\includegraphics[width=0.88\linewidth]{figures/fig_549.jpg}
\caption{Exhibit 14 - Exhibit 14: MRI procurement value increased by \$+20\textbackslash{}\%\$ yoy in May-26 vs. \$-50\textbackslash{}\%\$ in Apr-26 MRI}
\end{figure}
{\small\color{refgray}\textbf{References:} [R012] GS：CMS释放的信号比药价谈判本身更重要; [R013] GS：医疗健康板块的结构性拐点，比市场感知的更近; [R024] Citi：中国生物科技面临的地缘政治风险，市场可能高估了法案通过概率; [R036] GS：中国医疗设备招标的拐点正在到来，但真正的分化才刚刚开始}

\section{新兴科技：AI音乐变现与量子计算速度精度取舍}
\textbf{AI不是音乐的威胁而是超级粉丝变现的桥梁，量子计算真正变量是速度与保真度的取舍。}

\begin{itemize}
  \item Bernstein指出Spotify与环球音乐达成的AI翻唱和混音协议是增量收入池，而非重新分配现有收入。
  \item 保守估算2\%的用户愿意每月多付5美元，Spotify每年多赚3.5亿美元，版权方分到约2.1亿美元。
  \item 基准情景5\%渗透率×8美元/月，版权方年增收8.4亿美元，提升现有流媒体收入7.7\%。
  \item Bernstein量子计算专家研讨会揭示制造型量子比特（超导、硅自旋）速度快但一致性缺失。
  \item 天然型量子比特（离子阱、中性原子）天生一致性好但操作慢，不同应用场景对速度和保真度要求不同。
  \item Bernstein认为量子计算不会是赢家通吃，而是多模态共存的生态系统。
  \item Visa与OpenAI的合作标志着卡片支付正式进入代理时代，Bernstein认为卡片网络正成为代理经济最核心的信任基础设施。
  \item Bernstein指出亚马逊供应链的真实威胁被市场高估，前置库存模型在供应链结构上存在天然局限。
\end{itemize}
\begin{figure}[htbp]
\centering
\includegraphics[width=0.88\linewidth]{figures/fig_031.jpg}
\caption{Exhibit 6 - EXHIBIT 6: As Spotify's gross margin climbs toward the 35\% to 40\% target, the content and distribution share of every revenue dollar falls from 75 cents in 2022 to about 64 cents by 2030, with the marketplace, up roughly fourfold i}
\end{figure}
\begin{figure}[htbp]
\centering
\includegraphics[width=0.88\linewidth]{figures/fig_032.jpg}
\caption{EXHIBIT 8 - EXHIBIT 8: We see AI-driven pricing power and incremental engagement more than offsetting royalty leakage and engagement drag, adding \textbackslash{}\textasciitilde{}\textbackslash{}\$2.3Bn to 2030E streaming TAM Recorded music streaming TAM, AI headwinds \& tailwinds (\textbackslash{}\$B)}
\end{figure}
\begin{figure}[htbp]
\centering
\includegraphics[width=0.88\linewidth]{figures/fig_033.jpg}
\caption{EXHIBIT 9 - EXHIBIT 9: AI-driven upside: Streaming TAM CAGR expands to 9.3\% through 2030E Streaming TAM CAGR to 2030E: AI v. Non-AI scenarios}
\end{figure}
{\small\color{refgray}\textbf{References:} [R004] Bernstein：AI不是音乐的威胁，而是超级粉丝变现的桥梁; [R015] Bernstein：量子计算的真正变量不是技术路线，而是速度与保真度的取舍; [R023] Bernstein：亚马逊供应链的真实威胁，被市场高估了; [R048] Bernstein：Visa与OpenAI的合作标志着卡片支付正式进入代理时代}

\section{结语}
综合来看，当前市场主线围绕AI结构性需求、亚洲通胀再定价与全球资本流动重构展开。中国出口的价量背离、欧洲半导体设备的结构性增长、以及亚洲动量交易的估值泡沫，共同指向一个高度分化的投资环境。市场对2027年AI资本开支的上行风险、日本央行加息路径、以及欧洲钢铁定价权重塑的定价可能仍不充分。
\newpage
\section{报告覆盖清单}
\begin{itemize}
  \item [R001] 宏观与利率：亚洲通胀反弹与央行紧缩路径 - NOM：中国通胀的结构性真相——市场低估了供给驱动的持续性，高估了需求复苏的弹性
  \item [R002] 中国贸易：价格驱动的结构性重定价与AI主导的出口新格局 - DB：中国贸易数据揭示的不是短期反弹，而是供给优势的再定价
  \item [R003] 区域市场：印度制造业竞争力缺口与韩元贬值悖论 - 摩根斯坦利：印度真正需要解决的，不是短期资本流入，而是制造业竞争力的结构性缺口
  \item [R004] 新兴科技：AI音乐变现与量子计算速度精度取舍 - Bernstein：AI不是音乐的威胁，而是超级粉丝变现的桥梁
  \item [R005] 半导体：欧洲设备需求周期尚未见顶与存储供给纪律重塑 - GS：欧洲半导体公司的估值重估才刚刚开始，市场低估了AI对功率半导体的需求弹性
  \item [R006] 风险偏好：科技股信心结构性重置与亚洲动量泡沫 - Bernstein：亚洲动量交易已进入估值泡沫区，减仓窗口正在关闭
  \item [R007] 外汇与资本流动：日元对冲需求与韩元资本外流悖论 - Citi：美元兑日元的真正支撑不是利差，而是日本股市催生的对冲需求
  \item [R008] 宏观与利率：亚洲通胀反弹与央行紧缩路径 / 外汇与资本流动：日元对冲需求与韩元资本外流悖论 - DB：亚洲通胀反弹正在重塑央行的利率路径，市场尚未完全定价这一轮紧缩的力度与广度
  \item [R009] 区域市场：印度制造业竞争力缺口与韩元贬值悖论 / 外汇与资本流动：日元对冲需求与韩元资本外流悖论 - NOM：韩元走弱的真正推手不是经济基本面，而是资本市场成功的“悖论”
  \item [R010] 能源与大宗：欧洲钢铁定价权重塑与储能政策重心转移 - 摩根斯坦利：中国大宗商品出口的“结构性分化”正在重塑全球定价权
  \item [R011] 宏观与利率：亚洲通胀反弹与央行紧缩路径 - GS：日本央行6月加息的可能性被市场低估，真正关键的是加息后的沟通路径
  \item [R012] 医疗健康：药价谈判结构性变化与设备招标拐点 - GS：CMS释放的信号比药价谈判本身更重要
  \item [R013] 医疗健康：药价谈判结构性变化与设备招标拐点 - GS：医疗健康板块的结构性拐点，比市场感知的更近
  \item [R014] AI/算力：资本开支上行风险与供应链结构性重置 - Citi：中国2万亿AI投资计划的核心不在于规模，而在于国产化率的硬约束
  \item [R015] 新兴科技：AI音乐变现与量子计算速度精度取舍 - Bernstein：量子计算的真正变量不是技术路线，而是速度与保真度的取舍
  \item [R016] 中国新能源车：渗透率突破60\%后的竞争结构重构 - DB：中国车企在欧洲的渗透已从“试探”进入“阵地战”阶段
  \item [R017] 地产与消费：中国餐饮品牌分化加速与家电估值逻辑切换 - GS：中国餐饮的五月数据揭示的不是需求疲软，而是品牌分化的加速
  \item [R018] AI/算力：资本开支上行风险与供应链结构性重置 - 摩根斯坦利：市场对800V DC延迟的担忧被夸大了，真正的信号是供应链加速而非推迟
  \item [R019] 风险偏好：科技股信心结构性重置与亚洲动量泡沫 - 美国银行：市场正在经历一场科技股信心的结构性重置，而非简单的避险
  \item [R020] 中国贸易：价格驱动的结构性重定价与AI主导的出口新格局 - NOM：中国5月出口强劲，但市场真正低估的不是需求，而是价格信号的重估
  \item [R021] AI/算力：资本开支上行风险与供应链结构性重置 - Bernstein：数据中心真正的拐点不在建设量，而在供电模式的结构性重置
  \item [R022] 宏观与利率：亚洲通胀反弹与央行紧缩路径 - 摩根斯坦利：中国金融体系真正的转折点不是刺激，而是“高质量”的内生循环
  \item [R023] 新兴科技：AI音乐变现与量子计算速度精度取舍 - Bernstein：亚马逊供应链的真实威胁，被市场高估了
  \item [R024] 医疗健康：药价谈判结构性变化与设备招标拐点 - Citi：中国生物科技面临的地缘政治风险，市场可能高估了法案通过概率
  \item [R025] 半导体：欧洲设备需求周期尚未见顶与存储供给纪律重塑 - GS：欧洲半导体设备市场的下一个增长动力，不是周期，是结构
  \item [R026] AI/算力：资本开支上行风险与供应链结构性重置 - 摩根斯坦利：CPO市场真正的分歧不在技术，而在量产节奏的预期差
  \item [R027] 未被主题摘要引用 - 摩根斯坦利：中国保险业真正的分化不在规模，而在产品策略与渠道控制力
  \item [R028] 中国贸易：价格驱动的结构性重定价与AI主导的出口新格局 - BARC：市场低估了中国出口的结构性韧性，AI与绿色技术正重塑贸易格局
  \item [R029] AI/算力：资本开支上行风险与供应链结构性重置 - GS：市场真正低估的不是AI需求，而是2027年资本开支的上行风险
  \item [R030] 中国贸易：价格驱动的结构性重定价与AI主导的出口新格局 - JPM：市场低估的不是出口增速，而是贸易结构对增长含义的根本改变
  \item [R031] 半导体：欧洲设备需求周期尚未见顶与存储供给纪律重塑 - Bernstein：半导体设备进口数据揭示的不仅是周期性复苏，更是结构性分化
  \item [R032] AI/算力：资本开支上行风险与供应链结构性重置 - Citi：中国2万亿AI基建计划真正重塑的不是需求，而是国内供应链的定价权
  \item [R033] 外汇与资本流动：日元对冲需求与韩元资本外流悖论 - GS：市场对香港银行跨境监管风险的担忧，可能被放大了十倍
  \item [R034] 中国新能源车：渗透率突破60\%后的竞争结构重构 - DB：中国车市正在经历的，不是周期波动，而是竞争结构的根本性重构
  \item [R035] 中国新能源车：渗透率突破60\%后的竞争结构重构 - DB：中国新能源车市最值得关注的不是总量，而是订单结构的剧烈分化
  \item [R036] 地产与消费：中国餐饮品牌分化加速与家电估值逻辑切换 / 医疗健康：药价谈判结构性变化与设备招标拐点 - GS：中国医疗设备招标的拐点正在到来，但真正的分化才刚刚开始
  \item [R037] 能源与大宗：欧洲钢铁定价权重塑与储能政策重心转移 - HSBC：储能正在成为全球能源政策的新重心，光伏反而面临需求拐点
  \item [R038] 能源与大宗：欧洲钢铁定价权重塑与储能政策重心转移 - HSBC：SNEC 2026揭示的真正信号，不是光伏回暖，而是能源系统的权力交接
  \item [R039] 地产与消费：中国餐饮品牌分化加速与家电估值逻辑切换 - JPM：中国家电的估值逻辑正在从“周期股”转向“复合增长引擎”
  \item [R040] 风险偏好：科技股信心结构性重置与亚洲动量泡沫 - JPM：零售投资者在科技股上的踩踏式抛售，暴露的不仅是情绪，更是仓位结构的脆弱性
  \item [R041] 未被主题摘要引用 - JEF：泡泡玛特的韧性来自股东结构变化，而非基本面改善
  \item [R042] 半导体：欧洲设备需求周期尚未见顶与存储供给纪律重塑 - 摩根斯坦利：市场真正低估的不是需求，而是供给纪律的长期重塑
  \item [R043] 中国新能源车：渗透率突破60\%后的竞争结构重构 - 摩根斯坦利：中国电动车市场正进入“产品攻势间歇期”的定价博弈
  \item [R044] 能源与大宗：欧洲钢铁定价权重塑与储能政策重心转移 - 摩根斯坦利：全球钢铁市场的真正拐点不在需求，而在欧洲供给侧的重新定价
  \item [R045] 中国新能源车：渗透率突破60\%后的竞争结构重构 - NOM：中国汽车市场2026年可能首次出现国内需求双位数下滑
  \item [R046] 风险偏好：科技股信心结构性重置与亚洲动量泡沫 - JPM：日本股市真正的博弈不在AI，而在价值与动量的交叉点
  \item [R047] 区域市场：印度制造业竞争力缺口与韩元贬值悖论 - Bernstein：印度350cc摩托车市场真正的较量，不是产品，是分销
  \item [R048] 新兴科技：AI音乐变现与量子计算速度精度取舍 - Bernstein：Visa与OpenAI的合作标志着卡片支付正式进入代理时代
  \item [R049] 未被主题摘要引用 - GS：游戏行业真正的拐点不在2026年的游戏阵容，而在成本结构与IP运营能力的分化
  \item [R050] 半导体：欧洲设备需求周期尚未见顶与存储供给纪律重塑 - JPM：MLCC材料端正在经历一场被低估的定价权迁移
\end{itemize}
\section{逐篇报告摘录}
\subsection{[R001] NOM：中国通胀的结构性真相——市场低估了供给驱动的持续性，高估了需求复苏的弹性}
{\small\color{refgray}归类：宏观与利率：亚洲通胀反弹与央行紧缩路径}

NOM：中国通胀的结构性真相——市场低估了供给驱动的持续性，高估了需求复苏的弹性

中国5月通胀数据公布后，市场的主流叙事是“通胀温和，政策空间犹存”。CPI同比持平于1.2\%，PPI同比升至3.9\%，似乎验证了经济温和复苏、通胀压力可控的判断。但NOM这份最新研报揭示了一个更复杂的图景：当前通胀并非需求回暖的信号，而是一场由外部供给冲击主导的“成本推入式”再定价。市场真正需要关注的，不是通胀本身的高低，而是这种通胀的结构——它正在挤压企业利润、侵蚀居民购买力，并可能迫使政策在“稳增长”与“防输入性通胀”之间做出更艰难的选择。

这份报告的核心判断是：中国通胀已经走出了通缩，但走向的并非温和复苏，而是一种“类滞胀”的初期形态。PPI的高企主要由能源和AI相关材料驱动，剔除这两项后，PPI可能仍是负值。而CPI的稳定，则更多依赖能源价格的基数效应，而非消费需求的真正回暖。这意味着一件事：如果外部供给约束持续，中国将面临“成本上升但需求不足”的尴尬局面，这比单纯的通缩或通胀都更难应对。

1. 当前通胀的“双轨制”：上游热、下游冷，中间环节的利润正在被系统性压缩

NOM报告中最具穿透力的观察，在于对通胀结构性的拆解。5月PPI同比3.9\%的读数，看似强劲，但它的驱动力高度集中。NOM估算，能源和AI相关材料两项合计贡献了PPI通胀的3.96个百分点。换句话说，剔除这两个板块，PPI通胀将是负值。

具体来看，石油相关行业（占PPI篮子约14\%）贡献了1.90个百分点，有色矿采和加工业贡献了1.80个百分点，科技制造业贡献了0.26个百分点。而下游消费品部门的PPI通胀仍为负值（-0.8\%），尽管较4月的-1.0\%有所改善，但依然处于收缩区间。

这意味着什么？当前的通胀不是经济全面过热的结果，而是全球供应链瓶颈在特定环节的集中爆发。上游企业受益于价格上涨，但下游制造企业和最终消费者却在承担成本上升的压力。这种“上游热、下游冷”的结构，本质上是对中游制造业利润的系统性挤压。那些无法将成本转嫁给下游的企业，将面临毛利率的快速恶化。对于投资者而言，这提示了一个关键风险：在PPI高企但CPI低迷的背景下，制造业板块的整体盈利预期可能需要下修。

2. 能源冲击的“二次传导”：中国作为净进口国的贸易条件正在恶化，这比通胀本身更值得警惕

NOM报告特别强调了一个常被市场忽视的宏观视角：中国是全球最大的能源净进口国，也是芯片的净进口国。当能源和芯片价格因外部供给冲击而飙升时，中国的贸易条件正在显著恶化。

报告指出，贸易条件的恶化会在宏观层面收缩国际收支，并在微观层面挤压国内生产者和消费者。这并非理论推演。5月活动数据的广泛不及预期已经表明，由外部驱动的再通胀无法将中国经济从困境中解救出来。居民消费已被“以旧换新”政策的回补效应和持续的房地产下行所削弱，而供给驱动的再通胀将进一步抑制购买力。工业生产则可能因能源和化工原材料的供应中断而受损。

这个判断的深层含义在于：市场可能低估了能源价格对国内经济的“二次传导”效应。第一次传导是直接的成本推升，体现在PPI上。第二次传导则是通过恶化贸易条件、压缩企业利润、抑制居民消费，最终反噬经济增长。如果布伦特油价维持在同比30\%以上的高位，这种二次传导的影响将在未来两个季度内逐步显现。对于资产定价而言，这意味着对周期股的估值需要更谨慎，而对受益于成本下降的消费和服务板块，复苏的窗口可能比预期更窄。

3. AI超级周期正在成为PPI的新推手，但它的传导路径和持续性仍需验证

NOM报告的一个新亮点，是将“AI超级周期”纳入了PPI的分析框架。报告明确指出，除全球油价外，芯片价格的上涨已成为PPI通胀的新兴力量。芯片价格的飙升通过两个渠道传导至PPI：一是直接推高国内半导体产品（包括存储

[... middle omitted ...]

，AI驱动的价格上涨是否具有持续性？NOM报告没有完全展开这一部分。当前芯片价格的上涨，部分源于全球AI基础设施建设的军备竞赛，部分源于地缘政治导致的供应紧张。如果AI资本开支增速放缓，或者全球芯片产能扩张超预期，这一推力的边际效应可能递减。但至少在当前阶段，AI相关材料已经成为PPI通胀中不可忽视的结构性力量。对于投资者而言，这意味着与AI上游材料（如铜、先进封装材料、高带宽存储）相关的资产，其定价逻辑正在从“需求故事”切换到“通胀受益故事”。

4. 猪肉价格仍在拖累CPI，但这是否意味着消费需求真的“没救了”？

CPI通胀稳定在1.2\%，但核心CPI（剔除食品和能源）从1.2\%微降至1.1\%。其中，服务价格通胀从0.9\%放缓至0.8\%。最显著的拖累来自猪肉：猪肉价格同比下跌16.1\%，贡献了-0.31个百分点的CPI通胀。NOM报告指出，由于猪肉供应充裕且国内需求疲弱，猪肉价格短期内仍将承压。4月生猪屠宰量同比增速反弹至26.7\%，进一步印证了供给过剩的局面。

\subsection{[R002] DB：中国贸易数据揭示的不是短期反弹，而是供给优势的再定价}
{\small\color{refgray}归类：中国贸易：价格驱动的结构性重定价与AI主导的出口新格局}

这份由DB在2026年6月9日发布的宏观报告，核心判断不是“出口又变好了”这样的短期叙事。报告真正值得关注的信号是：中国贸易的结构性驱动力正在从单纯的规模扩张，转向以AI相关产品、能源密集型产品为代表的定价权优势。5月出口同比增长19.4\%，进口增长27.4\%，但这组数字背后的含义远比表面增速更深刻。

市场习惯将中国贸易数据解读为全球需求的晴雨表，或者简单归结为基数效应。DB的报告提供了一个不同的框架：它把增长拆解为“AHEAD”因素——AI相关产品、汽车、其他机械——并发现这些因素贡献了16个百分点的出口增长。这意味着，即便剔除全球周期的影响，中国自身的供给能力正在成为贸易增长的主引擎。

为什么现在这个判断重要？因为全球投资者正在重新评估中国资产的定价逻辑。如果贸易增长主要依赖的是中国在AI硬件、能源密集型产业的成本与技术优势，而不是单纯的外部需求拉动，那么这种增长的可持续性和质量都会不同。这对于汇率预期、产业投资方向、甚至全球供应链的重新布局都有直接的含义。

报告提出的一个关键观察是：AI相关出口在5月同比增速翻倍至81\%，其对整体出口的贡献从1-4月的5个百分点跃升至10个百分点。这不是一个边际变化，而是量级上的跃迁。与此同时，能源和石化产品的出口也显著走强——精炼油出口从-1\%反弹至27\%，塑料制品从7\%加速至12\%。这两条线索叠加在一起，指向一个更根本的叙事：中国正在将其在制造业和能源领域的规模优势，转化为全球贸易中的定价权。

以下是我们从这份报告中提炼出的五个层次判断，它们共同支撑一个核心主张：市场低估的不是中国贸易的短期韧性，而是供给侧优势正在被重新定价。

1. 这轮贸易加速中，AI相关出口的翻倍增长不是偶然，而是中国在硬件制造领域系统性优势的集中体现

DB报告中最引人注目的数据点是AI相关出口增速从42\%跃升至81\%。这不是一个孤立事件。报告定义的AI相关产品包括集成电路、手机、计算机和其他数据处理设备。这些产品恰恰是中国在过去十年持续投入、形成完整产业链的领域。

从供应链角度看，中国在AI硬件制造上的优势不仅是成本，更是生态完整度。从芯片封装测试到终端组装，从精密模具到物流网络，这种集群效应使得中国企业在面对AI需求爆发时，能够以更快的响应速度和更低的边际成本扩大产出。81\%的同比增速反映的不是一时的订单集中，而是这种系统能力的释放。

对于投资者而言，这意味着AI相关硬件出口的强势表现具有结构性支撑，不应被视为短期脉冲。报告没有展开的一个问题是：这种增长在多大程度上来自国内企业的自主创新，又在多大程度上来自全球科技巨头在中国的生产布局？这需要进一步的产业链调研来验证，但无论如何，中国作为AI硬件制造中心的地位正在被强化。

2. 能源密集型产品的出口走强，揭示了中国在能源成本端的全球竞争力正在转化为贸易份额

精炼油从-1\%到27\%、塑料从7\%到12\%的增速变化，在报告的叙事中可能被AI的光芒所掩盖，但这一条线索对长期格局的理解同样关键。DB明确提到，这“凸显了中国在能源密集型行业的定价优势”。

这个判断的逻辑是：在全球能源价格波动加剧的背景下，中国依托相对稳定的国内能源供应体系——包括煤炭、电力、以及日益增长的天然气和可再生能源——获得了生产成本上的比较优势。当欧洲、日韩等地区的制造业面临能源成本高企的压力时，中国企业有能力以更低价格出口能源密集型产品，从而抢占市场份额。

报告没有深入讨论的一点是：这种定价优势是否可持续？中国的能源结构转型——从煤炭向清洁能源的过渡——可能会在长期改变成本曲线。但至少在中期，中国在能源密集型产业上的竞争力是一个需要被纳入全球资产定价框架的因素。对于化工、塑料、炼油等行业的全球竞争格局而言，这意味着来自中国企业的竞争压力将持续存在，甚至可能加剧。

3. 对发达市场出口的突然飙升，背后是AI需求与关税扰动的双重作用，但结构性因素更值得关注

报告显示，中国对主要发达市场——美国、日本、韩国、台湾——的出口增速从1-4月的2\%飙升至5月的30\%。其中对美出口从-13\%反弹至35\%，对韩国出口从24\%升至42\%。DB将其归因于AI相关需求的拉动，同时也承认对美出口的反弹部分反映了低基数和关税扰动。

这里需要区分两个不同的逻辑。关税扰动带来的贸易抢跑是暂时性的——企业为了规避关税上调而提前发货，这种效应会在政策落地后消退。但AI需求驱动的增长则更具持续性。韩国和台湾的强劲进口数据尤其值得注意：这两个经济体是全球半导体产业链的核心节点，它们从中国进口的激增，很可能反映的是中国作为中间品供应国在AI产业链中的嵌入程度在加深。

对于观察者而言，这意味着评估中美贸易关系时，不能仅仅盯着关税谈判的短期进展。更深层的变量是：全球科技产业链对中国的依赖程度是否在上升？如果答案是肯定的，那么即便关税政策收紧，贸易流量的调整也会比预期更慢、成本更高。

4. 进口数据中的能源滞后效应，暗示中国正在成为全球能源贸易的“价格接受者”与“量级塑造者”

进口增长27.4\%中，能源进口的滞后反弹是一个值得单独分析的细节。报告指出，3-4月主要供应商对中国的能源出口已经增加，但直到5月才在进口数据中体现——原油从0\%加速至15\%，煤炭从-7\%至10\%，天然气从-16\%至10\%。

\subsection{[R003] 摩根斯坦利：印度真正需要解决的，不是短期资本流入，而是制造业竞争力的结构性缺口}
{\small\color{refgray}归类：区域市场：印度制造业竞争力缺口与韩元贬值悖论}

摩根斯坦利：印度真正需要解决的，不是短期资本流入，而是制造业竞争力的结构性缺口

市场对印度近期的政策反应普遍积极——取消资本利得税、扩大政府债券投资范围、提高海外个人投资限额，一系列措施密集出台。但摩根斯坦利最新发布的亚洲宏观报告提出了一个冷静的判断：这些措施能缓解短期压力，但无法解决根本问题。

这份由首席亚洲经济学家Chetan Ahya领衔的报告，核心论点是：印度面临的不是流动性问题，而是国际收支的结构性承压。而解决这一压力的唯一路径，是制造业竞争力的系统性提升。资本账户管理只是治标，制造业出口能力的建设才是治本。

报告的数据框架清晰地展示了问题的严重性：2024年12月至2026年3月，印度国际收支累计录得520亿美元的赤字。这个数字是 年期间250亿美元赤字的两倍以上。更值得关注的是，卢比的实际有效汇率已经跌至低于10年均值3.7个标准差的位置——这不是短期波动，而是系统性的失衡信号。

本文将从五个层次拆解这份报告的核心洞察，并在最后提出一个报告尚未完全回答的关键问题，供读者进一步思考。

1. 资本外流的根源不是市场情绪，而是相对盈利增速的结构性分化

许多市场观察者将印度的资本外流归因于地缘政治风险或全球利率环境。但摩根斯坦利的分析揭示了更深层的原因：印度企业的相对盈利增速正在失去吸引力。

报告提供了一组关键对比数据：2026年第一季度，印度BSE500指数成分股盈利同比增长13\%。这个数字本身并不差，但放在区域比较中就显得苍白——韩国KOSPI指数盈利增长171\%，台湾加权指数增长49\%，日本TOPIX增长33\%，美国标普500增长29\%。印度是主要经济体中盈利增速最低的几个之一。

这意味着什么？全球资本正在重新定价“增长溢价”。AI和工业供应链重构带来的盈利脉冲，正在将资金从传统的“人口红利”故事抽离，转向更直接受益于技术周期的市场。印度在这一轮资本再配置中，处于相对不利的位置。

报告进一步指出，印度的油气贸易逆差占GDP比例达到3.4\%，是亚洲主要经济体中最高之一。能源价格冲击通过贸易条件恶化，进一步削弱了印度对外资的吸引力。这不是一个临时性问题——只要能源结构没有根本性改变，这一结构性拖累就会持续存在。

面对国际收支压力，印度政府和央行已经打出了一套组合拳：取消政府证券投资的资本利得税和股息税、扩大“完全可进入路径”下的债券范围、提高海外个人投资限额、为国有企业外部商业借款提供汇率对冲支持。报告明确认为，这些措施在短期内是有效的。

但关键问题在于：这些措施解决的是“流量”问题，而不是“存量”问题。它们可以缓解卢比贬值压力、降低汇率波动，但无法提升印度经济自身的储蓄和创汇能力。

报告给出了一个清晰的预测路径：在能源价格温和回落的假设下，印度经常账户赤字将从2026财年的0.6\%扩大到2027财年的1.8\%，之后在2028财年收窄至1.2\%。但即使按照这个相对乐观的假设，经常账户赤字仍将持续存在。如果没有制造业出口的增长来弥补这一缺口，国际收支压力将反复出现。

这引出了一个核心判断：印度需要的是储蓄率的可持续提升，而储蓄率的提升必须依赖于就业增长，就业增长又必须依赖于制造业的扩张。这是一个环环相扣的逻辑链条。

报告提出了一个容易被忽视的洞察：制造业的扩张能够同时从两个方向改善国际收支。一方面，制造业出口增长可以直接缩小贸易逆差，改善经常账户；另一方面，制造业吸引的外国直接投资能够增强资本账户的稳定性。

这与印度过去依赖服务业出口的模式有本质区别。服务业出口虽然利润率较高，但创造就业的能力有限，且对资本账户的直接拉动作用较弱。制造业则不同——它既能够大规模吸纳就业（从而提升储蓄率），又能够通过供应链本地化吸引长期资本流入。

报告特别强调，印度制造业的目标不应只是低端组装

[... middle omitted ...]

“人口红利”之间找到自己的位置

报告虽然没有直接使用这个表述，但它的数据框架揭示了一个更深层的战略问题：全球资本正在按照“技术周期”而非“人口周期”重新配置。韩国和台湾的盈利增长主要受益于AI相关产业链的需求爆发，而印度在这一轮技术浪潮中的参与度相对有限。

这意味着，印度不能简单地复制中国的制造业崛起路径。中国制造业的腾飞恰逢全球贸易自由化的黄金时期和WTO的红利期。而当前的全球贸易环境更加碎片化，技术竞争更加激烈，供应链重构更多是出于安全考虑而非效率考虑。

印度需要找到自己的差异化定位。报告的观点是，制造业竞争力的提升可以帮助印度在“中国+1”的供应链多元化趋势中占据有利位置。但这要求印度在基础设施、营商环境、劳动力技能等方面做出实质性的改进，而不仅仅是政策层面的表态。

这里有一个报告没有完全展开的问题：印度制造业的竞争优势究竟在哪里？是劳动力成本、市场规模、还是政策激励？如果只是依赖关税壁垒和补贴，而不解决生产效率的根本问题，那么制造业的增长可能难以持续。

\subsection{[R004] Bernstein：AI不是音乐的威胁，而是超级粉丝变现的桥梁}
{\small\color{refgray}归类：新兴科技：AI音乐变现与量子计算速度精度取舍}

Bernstein：AI不是音乐的威胁，而是超级粉丝变现的桥梁

音乐产业正在经历一场被市场低估的结构性转变。当大多数讨论聚焦于AI是否将取代人类创作时，一份来自Bernstein的最新研报提出了相反的判断：AI恰恰是唱片公司实现“超级粉丝”变现的关键工具，而非生存威胁。

这份报告的触发点来自Spotify最近的投资者日活动。当天最重磅的公告不是用户增长或定价策略，而是一份史无前例的授权协议——Spotify与环球音乐集团达成合作，允许Premium用户基于授权曲库付费创作AI翻唱和混音版本。这是全球主流流媒体平台与主要唱片公司之间首次建立基于“同意、署名和补偿”原则的AI音乐框架。

Bernstein团队长期认为，AI是唱片公司实现超级粉丝变现的桥梁，而非行业颠覆者。这份协议就是最清晰的证据。从Suno和Udio最初闯入AI音乐领域至今不过短短数年，整个品类已经可以通过全球最大音乐平台的框架，演变为唱片公司明确的增长引擎。

理解这一变化的关键在于：这不是对现有流媒体收入的重新分配，而是在其之上叠加的全新收入层。

Spotify管理层在投资者日上明确了一个核心洞察——用户的付费意愿遵循幂律分布。头部一小部分高度活跃的用户愿意支付远超标准订阅费用的金额。公司不再试图仅通过提价和增加订阅人数来扩大音乐收入池，而是通过付费附加服务来捕获这部分“超级粉丝”的价值。

与环球音乐达成的翻唱和混音协议，正是这一策略在版权方层面的首次落地。它的结构是付费Premium附加服务，录音版权所有者和词曲作者共同分享粉丝创作带来的收益。对行业而言，最关键的信号是：这笔收入是增量。它叠加在流媒体收入之上，而非重新分配现有池子。而且，由于曲库本身就是输入素材，Spotify不需要增加用户的收听时长就能扩大对版权方的支付。

Bernstein用Spotify唯一公开披露过的附加服务“有声书+”作为锚点来估算规模。该服务年化收入约1亿美元，付费用户超过100万，意味着每位付费用户每月贡献约8美元，在总订阅用户中的渗透率约为0.3\%。这一成果在不到一年内、仅覆盖22个市场的情况下达成，为附加服务的定价区间（5至10美元）提供了可信的参照，也表明渗透率还有多年的增长空间。

基于当前约2.93亿订阅用户，Bernstein给出了三档情景测算。保守情景下，2\%的渗透率、5美元月费，Spotify每年可增加约3.5亿美元收入，其中约2.1亿美元流向版权方。基准情景下，5\%渗透率、8美元月费，Spotify年增量收入约14亿美元，版权方分得约8.45亿美元，相当于当前行业从Spotify获得的110亿美元年收入的7.7\%。乐观情景下，10\%渗透率、10美元月费，Spotify年增量达35亿美元，版权方分得约21亿美元，接近当前池子的19\%。

这些数字背后还有两个结构性优势。第一，翻唱和混音产品以词曲版权为核心，出版商和词曲作者获得的份额可能高于普通流媒体播放（后者以录音版权为主）。第二，购买附加服务的用户是整个用户群中最忠诚的部分，流失率更低、生命周期价值更高，意味着这笔增量收入比等额的普通订阅收入更具持久性。

Bernstein预期，华纳音乐将很快跟进签署自己的AI翻唱和混音授权协议。这一判断基于两个并行的增长引擎。

第一个引擎是Spotify。华纳音乐每年已从Spotify获得约12亿美元收入，付费创作附加服务可以干净地叠加在这一基础之上。第二个引擎是AI原生平台。华纳音乐是首家与Suno达成授权协议的主要唱片公司，也已与Udio达成和解，从这些快速扩张的池子中分得一杯羹。单是Suno一家，年化收入已达约2至3亿美元，且仍在快速扩张。

Bernstein预计，到2027年，授权AI将为唱片公司贡献低个位数的收入增长动力。这一判断

[... middle omitted ...]

控制权和收入份额。这解释了为什么Bernstein给予索尼“市场表现”评级，而非“跑赢大盘”。

这一判断的底层逻辑是Bernstein对Spotify战略演进的框架性分析。报告将Spotify的演进划分为三个阶段：策展时代（接入）、推荐时代（算法）和生成时代（AI）。在每个阶段，Spotify都为行业提供了新的变现方式——先是流媒体版税，然后是市场推广工具，现在是授权AI——但同时也在逐步掌握对发现机制和用户关系的控制权。唱片公司的角色从守门人逐步收窄为稀缺授权的供应商。

4. HYBE已经将超级粉丝变现工业化，Spotify正在向K-pop模式收敛

Bernstein报告中一个极具洞察力的判断是：HYBE已经将Spotify正在追求的方向工业化了。

HYBE的商业模式建立在“爱”而非“听”的基础上。通过Weverse这一封闭生态系统，HYBE控制着知识产权和粉丝接触渠道，在偶像的全生命周期内实现高频变现——内容、周边、互动，每一个触点都在叠加每用户平均收入。

\subsection{[R005] GS：欧洲半导体公司的估值重估才刚刚开始，市场低估了AI对功率半导体的需求弹性}
{\small\color{refgray}归类：半导体：欧洲设备需求周期尚未见顶与存储供给纪律重塑}

GS：欧洲半导体公司的估值重估才刚刚开始，市场低估了AI对功率半导体的需求弹性

这份GS最新发布的欧洲半导体研报，表面上只是上调了英飞凌和意法半导体的盈利预测与目标价，但真正值得关注的信号并非数字本身，而是其背后一套完整的估值逻辑重置。

报告传递了一个清晰判断：市场此前对AI半导体需求的认知过于集中在GPU和高端计算芯片上，而忽略了AI基础设施大规模部署对功率半导体、光通信等“使能器件”的结构性拉动。GS将英飞凌12个月目标价从75欧元上调至88欧元，意法半导体目标价从42欧元上调至58欧元，估值倍数分别从16倍和10倍的EV/EBITDA提升至18倍和13倍。但比目标价更重要的，是支撑这一倍数扩张的逻辑——这两家公司正在经历从周期性半导体公司向AI结构性增长公司的身份转换。

以下是我们从这份报告中提炼出的五个核心洞察，以及一个报告尚未完全回答的关键问题。

1. 英飞凌的估值扩张并非凭空而来，其背后是收入增速翻倍的支撑

GS对英飞凌的估值调整，最值得关注的是其推导逻辑。报告明确指出，英飞凌的历史中位两年远期收入复合增长率约为9\%，而当前预测值约为17\%，几乎翻倍。这一增速提升，与GS给予的估值倍数扩张幅度高度吻合——从历史中位约10倍的12个月远期EV/EBITDA，提升至目标价隐含的约21倍。

这不是简单的“给AI概念股溢价”，而是基于收入增速变化与估值倍数之间的历史对应关系。GS用同样的逻辑检验了盈利增速：英飞凌历史两年远期EPS复合增长率为19\%，当前预测为43\%，扩张幅度约2.3倍；而目标价隐含的12个月远期PE从历史约20倍提升至约39倍，扩张幅度同样约为2倍。

这意味着，GS并非在给英飞凌一个“AI故事溢价”，而是在系统性地论证：当一家公司的增长曲线发生结构性上移时，市场给予的估值倍数应当相应调整。这一逻辑如果成立，对欧洲半导体板块的定价体系将产生深远影响。

与英飞凌相比，意法半导体的估值调整幅度更为显著。GS将其12个月目标价从42欧元大幅上调至58欧元，上调幅度接近40\%。估值倍数从10倍提升至13倍的EV/EBITDA，目标价隐含的12个月远期PE更是从历史约18倍跃升至约38倍。

支撑这一调整的收入增速数据同样惊人。意法半导体的历史中位两年远期收入复合增长率为7\%，当前预测为14\%，同样实现翻倍。但EPS增速的扩张更为剧烈——历史两年远期EPS复合增长率为29\%，当前预测高达72\%，扩张幅度接近2.5倍。

然而，意法半导体的评级仍为中性，而非买入。这一反差本身就值得深思。报告给出的主要风险包括：消费类半导体库存调整速度、碳化硅业务竞争对手的进展速度，以及当前有利定价能否持续。换句话说，GS承认意法半导体在AI基础设施领域（尤其是光通信）的增长前景，但对其传统业务的周期性和竞争格局仍持保留态度。

3. 市场对AI半导体需求的认知存在结构性盲区：功率器件和光通信才是本轮上调的核心驱动

这份报告最值得提炼的洞察，不是目标价本身，而是GS为何上调这些预期。对于英飞凌，驱动因素来自AI相关功率半导体需求，包括GPU和CPU的电源管理芯片。报告特别提到“来自AI半导体供应链的近期积极数据点”，以及公司计划中的德累斯顿模块化产线，后者可在中期支持50亿欧元的收入运行率。

对于意法半导体，增长动力则来自AI基础设施中的光通信需求。报告指出“近期客户互动评论”和“AI基础设施中对其产品的持续需求”是上调预期的主要原因。

这两条线索指向同一个判断：AI基础设施的资本开支正在从核心计算芯片向外围器件扩散。当数据中心从建设阶段进入运营阶段，功率管理、光互联、热管理等配套器件的需求弹性可能被严重低估。这不仅是英飞凌和意法半导体的故事，也可能映射到整个欧洲工业半导体生态。

4. 盈利预测的

[... middle omitted ...]

S明确解释了这一差异的来源：更高的收入基数叠加更低的产能利用率费用（underutilization charges），以及运营杠杆的自然释放。这意味着，这些公司并非仅仅卖出了更多的产品，而是在产能利用率提升的过程中，固定成本被摊薄，利润率结构发生改善。

这一逻辑对于理解半导体公司的盈利周期至关重要。在需求复苏初期，收入增长对盈利的拉动往往被低估，因为固定成本占比高。而一旦产能利用率越过盈亏平衡点，利润增速将显著快于收入增速。GS正在押注这一拐点已经到来。

5. 报告尚未完全回答的关键问题：AI驱动的增长能否持续超越半导体周期的均值回归

尽管GS的估值逻辑自洽，但报告本身留下了一个关键悬而未决的问题：当前AI驱动的增长能否持续，以至于改变英飞凌和意法半导体的长期增长中枢？

GS使用了两年远期收入复合增长率作为估值锚定，但这一指标本质上是中短期预测。如果AI需求在两年后趋于平稳，或者面临竞争加剧导致的定价压力，那么当前的高增速可能只是周期性峰值，而非结构性跃迁。

\subsection{[R006] Bernstein：亚洲动量交易已进入估值泡沫区，减仓窗口正在关闭}
{\small\color{refgray}归类：风险偏好：科技股信心结构性重置与亚洲动量泡沫}

Bernstein：亚洲动量交易已进入估值泡沫区，减仓窗口正在关闭

这份报告的核心判断并不复杂，但它的警示信号值得认真对待：亚洲市场（除日本和中国外）的动量交易已经进入历史级别的估值与预期双泡沫区间，而韩国和台湾市场正是风险最集中的区域。Bernstein量化策略团队在6月初发布的这份报告中，明确建议削减这两个市场的动量敞口。

为什么现在重要？因为过去几个月，动量策略在亚洲取得了极为罕见的超额收益——在除日本外的亚洲市场，价格动量因子年内跑赢基准36\%至38\%；如果剔除中国，这一数字更是飙升至53\%至60\%。韩国和台湾的动量因子年内绝对超额收益甚至超过50\%。但极端收益往往伴随着极端风险。Bernstein团队指出，这些市场的动量交易不仅估值创下历史新高，盈利预期同样处于前所未有的高位，而这两者叠加在一起，构成了一个脆弱性极高的组合。

报告给出的建议很清晰：削减台湾和韩国的动量敞口，转向反动量篮子，即高股息或低波动资产。但真正值得深思的，不是这个建议本身，而是它背后反映的市场结构变化——动量因子与成长因子的相关性已升至历史高位，而动量在低波动和价值股中却在持续下降。这种极端的市场内部分化，往往预示着风格切换的临界点正在逼近。

1. 动量交易的估值泡沫并非均匀分布，台湾和韩国是真正的风险中心

Bernstein的量化分析显示，亚洲（除日本、除中国）的动量组合正处于历史最高估值水平，无论是12个月远期市盈率还是市净率，都已突破此前所有周期的高点。具体来看，台湾和韩国动量篮子的估值均达到历史峰值，而且从6月初开始已经出现估值收缩的迹象。相比之下，中国和印度的动量股距离泡沫估值区间还相去甚远，日本动量股的估值甚至仍低于长期平均水平。

这里需要特别关注一个细节：报告区分了价格动量和盈利动量。亚洲（除日本）的盈利动量组合已经处于泡沫区间相当一段时间，而价格动量组合虽然估值上升，但尚未达到2021年的峰值。然而，一旦剔除中国，情况就完全不同——亚洲（除日本、除中国）的价格动量组合市净率已升至历史均值上方4.2个标准差，这是报告数据覆盖期内从未见过的水平。

这意味着什么？台湾和韩国市场的动量交易，已经不仅仅是“贵”的问题，而是进入了历史极端区间。当估值达到这种水平时，任何负面催化剂都有可能触发大规模的去杠杆和仓位调整。对于持有这些市场动量仓位的投资者来说，当前的问题不是“是否应该减仓”，而是“减仓窗口还有多久”。

Bernstein团队指出，今年动量股上涨的主要驱动力并非估值扩张，而是盈利上调。AI技术带来的盈利预期提升，推动亚洲科技股的盈利预期突破了所有历史水平。报告数据显示，亚洲动量组合、亚洲（除中国）动量组合，甚至日本动量组合，都出现了创纪录的盈利上调幅度。

具体到国家层面，中国和台湾的动量组合盈利预期均处于历史最高水平，韩国动量组合虽然略低于2009年的峰值，但同样处于极高位置。问题在于，当估值已经极度拉伸时，这些极端乐观的盈利预期就构成了一个脆弱的基础——只要有任何迹象表明盈利增长可能放缓或不及预期，估值和预期的双重修正就会同时发生。

报告没有明确说，但这里隐含着一个关键判断：当前市场对AI相关公司的盈利预期，可能已经过度反映了技术变革的乐观前景。一旦出现盈利预警或指引下修，动量交易的反转速度可能会非常快。对于投资者而言，需要密切关注即将到来的财报季中，台湾和韩国科技龙头的盈利指引是否能够支撑当前的高预期。

3. 市场内部分化已达到极端水平，低波动策略的崩溃是风格切换的前兆

Bernstein的量化分析还揭示了一个重要的结构性变化：亚洲动量因子与成长因子的相关性已升至历史高位，而动量在价值股、质量股和低波动股中的表现却在持续下降。最值得注意的是，低波动策略的动量水平已降至过去25年来的最低点。

[... middle omitted ...]

基于这一逻辑。

对于资产配置者来说，当前需要思考的不是“动量交易是否已经见顶”，而是“当动量反转发生时，自己的组合是否做好了准备”。低波动策略的极端表现，实际上是在提醒投资者：市场可能已经极度拥挤在同一个方向上。

Bernstein的报告给出了清晰的判断和建议，但有两个关键问题仍然悬而未决。

第一，动量反转的具体触发因素是什么？报告提到了地缘政治风险的重新升级可能加速动量反转，但并没有给出更具体的催化剂分析。是AI盈利预期的下修？是美联储政策转向？还是台湾和韩国市场的特定事件？对于读者来说，理解可能的触发因素，比知道“应该减仓”这个结论更为重要。

第二，动量反转的幅度会有多大？报告指出估值和预期都处于历史极端水平，但并没有量化可能的下跌空间。2000年互联网泡沫破裂时的动量反转，或者2021年成长股估值收缩时的跌幅，可以作为参考，但每个周期都有其独特性。如果AI技术确实代表了生产力范式的转变，那么当前的估值是否可能被新的增长路径所证明？报告没有展开讨论这一点。

\subsection{[R007] Citi：美元兑日元的真正支撑不是利差，而是日本股市催生的对冲需求}
{\small\color{refgray}归类：外汇与资本流动：日元对冲需求与韩元资本外流悖论}

Citi：美元兑日元的真正支撑不是利差，而是日本股市催生的对冲需求

市场对日元走强的共识正在变厚。美日长期利差已经收窄到足以清晰支持日元升值的方向，短期利差的信号同样不容忽视。但美元兑日元汇率却依然顽固地维持在160附近。这种背离让许多宏观交易员困惑——是模型错了，还是市场定价里藏了一个被低估的结构性因素？

Citi这份题为《Japan FX: An analysis of the USD/JPY price formation mechanism》的报告，给出一个值得认真对待的判断：当前美元兑日元的强势，很大程度上来自日本股市处于历史高位所催生的日元卖出对冲需求。这不是一个短期噪音，而是一种与股市市值深度绑定的结构性买盘。Citi维持对美元兑日元的中长期看空观点，认为160日元附近就是顶部区间，年底前将回落至155以下。但要控制日元在短期内的阶段性走弱，仍需要日本央行货币政策正常化以及日本当局的卖出美元干预。

这份报告的核心贡献，不在于给出了一个具体的点位预测，而在于提供了一个分析框架——如何区分趋势、周期和噪音，以及为什么传统利差模型在当前环境下可能严重低估了股市对冲需求的力量。

1. 利差收窄已经明确指向日元升值，但汇率没有跟随，说明市场存在“第五因素”

Citi在报告中开宗明义地指出，美元兑日元的定价受到四个基本面因素的共同影响：美日货币政策差距、整体市场风险偏好、日本的国际收支状况，以及美元整体方向（尤其是对欧元的走势）。这四个因素构成了一个基本估值框架。

然而，报告明确承认，即便这四个因素全部纳入模型，依然存在一个“残差”——实际价格与模型估值的每日偏差。这个残差通常被归因于短期投机活动造成的错误定价。但问题在于，预测这个残差恰恰是市场参与者最困难也最重要的任务。

当前的情况正是如此。美日10年期国债利差已经显著收窄，按历史经验应该推动日元走强。但实际汇率却没有做出相应调整。Citi的多重回归模型估算出的当前公允价值大约在161日元附近，而实际汇率略低于这个水平，说明日元甚至还有轻微的低估——这很可能反映了日本当局的买入日元干预。

关键在于，这个“轻微低估”本身说明了一个反直觉的事实：不是利差失效了，而是有一个强大的对冲需求正在抵消利差的推动力。这个对冲需求就是Citi所说的“第五因素”——它不能被简单归为噪音，而是有结构性支撑的持续买盘。

2. 日本股市的历史高位正在产生一次“被动日元卖出”，规模可能远超市场预期

Citi的分析中，最具洞察力的部分是对日本股市（TOPIX）在模型中解释力的发现。在双层模型的上层结构中，TOPIX和美元指数（DXY）是解释力最强的两个变量，而长期利差和贸易条件指数的解释力反而相对较弱。

报告明确指出，日本股市对美元兑日元的影响具有高度的同时性——两者相互影响。但Citi的每日分析表明，美日利差方向仍然是美元兑日元高度敏感的变量，只是这种影响在模型中可能通过美元指数间接体现出来。

更值得关注的是，Citi认为日本股市处于历史高位这一事实，直接催生了大规模的日元卖出对冲需求。当日本股市上涨时，海外投资者持有的日本股票头寸增加，这些投资者为了对冲汇率风险，会卖出日元、买入美元。这种对冲需求是系统性的、持续性的，与日本股市的市值规模直接挂钩。

日本股市在过去几年的涨幅显著，海外资金流入规模巨大。这些资金的对冲需求，可能已经形成了一个自洽的循环：股市上涨吸引更多资金流入，资金流入需要对冲，对冲卖出日元，日元贬值又进一步利好日本出口企业，从而支撑股市。这个循环的强度，取决于日本股市能否维持在高位。

3. Citi的双层模型提供了一个实用的分析工具，但模型本身存在必须正视的局限

Citi构建的双层模型是一个值得关注的创新。它将美元兑日元的分析拆分为两

[... middle omitted ...]

量化模型都有其内在局限，尤其是在分析市场数据时。Citi在报告中坦承，独立变量之间存在显著的内生性和同时性问题，使得严格意义上的统计显著性识别变得困难。向量自回归（VAR）等计量方法可以在一定程度上缓解这个问题，但无法彻底解决。

对于读者而言，这份报告的价值不在于模型的精确预测能力，而在于它揭示了一个重要的分析视角：当传统利差模型失效时，可能需要引入股市对冲需求这个变量。这也意味着，未来美元兑日元的走势，将在很大程度上取决于日本股市的走向——如果日本股市出现显著回调，对冲需求的逆转可能会加速日元升值。

4. 报告尚未完全回答的关键问题：对冲需求的规模究竟有多大，以及它何时可能反转

Citi的报告在分析框架上做出了重要贡献，但有几个关键问题没有完全展开。

第一，海外投资者对日本股市的对冲比例是多少？不同的机构投资者有不同的对冲政策。一些养老金和保险资金可能选择完全不对冲汇率风险，而一些对冲基金可能选择动态对冲。对冲比例的不同，会直接影响对冲需求对汇率的边际影响。

\subsection{[R008] DB：亚洲通胀反弹正在重塑央行的利率路径，市场尚未完全定价这一轮紧缩的力度与广度}
{\small\color{refgray}归类：宏观与利率：亚洲通胀反弹与央行紧缩路径 / 外汇与资本流动：日元对冲需求与韩元资本外流悖论}

DB：亚洲通胀反弹正在重塑央行的利率路径，市场尚未完全定价这一轮紧缩的力度与广度

2026年二季度，亚洲经济正在经历一个看似矛盾但逻辑清晰的阶段：增长放缓但高于趋势，通胀加速但政策干预有效。DB在其最新发布的《Asia Macro Insight》报告中提出一个核心判断——亚洲央行将进一步加息，且加息的力度和广度可能超出当前市场定价。这不是一次常规的周期调整，而是一次由供给侧结构变化、地缘政治冲击和国内政策权衡共同推动的利率重置。

报告指出，2026年亚洲CPI通胀预计将达到2.7\%，是去年的三倍。这一数字本身并不惊人，但真正值得关注的是通胀的来源结构：芯片价格飙升带来的贸易条件改善、伊朗战争推高的能源成本、以及中国再通胀进程的超预期推进。三个因素叠加，使得亚洲央行的政策天平从“保增长”明确转向“稳价格与稳汇率”。

对于投资者而言，这意味着需要重新审视亚洲资产的定价逻辑。利率上行不再是尾部风险，而是基准情景。而哪些经济体、哪些行业能够在这一轮紧缩中保持韧性，将决定未来12个月的相对表现。

DB的报告清晰勾勒出当前亚洲通胀的三个来源，它们彼此独立但相互强化。

首先是半导体价格的飙升。报告指出，韩国和台湾正受益于强劲的半导体定价，这不仅仅是周期性的需求回暖，更是AI投资驱动的结构性变化。芯片价格的上涨幅度足以抵消进口价格上涨对贸易条件的侵蚀，使得韩国和台湾的名义增长可能分别达到1996年和1987年以来的最高水平。这是一种“好通胀”——由高附加值出口拉动，伴随企业盈利改善和名义GDP扩张。

其次是伊朗战争引发的能源成本冲击。报告数据显示，部分经济体在2026年5月的交通/燃料通胀环比大幅跳升，菲律宾达到50\%，泰国35\%，韩国和台湾各20\%。这一冲击具有外生性和不可控性，且短期内看不到缓解迹象——报告明确提到“美伊协议仍遥不可及”。

第三是中国再通胀的超预期推进。报告认为，伊朗战争实际上“陡峭化了中国的再通胀路径”。更重要的是，中国房地产市场数据暗示了“持久的底部形成”，这一底部并非由大规模宏观刺激驱动，而是由改善中的市场基本面支撑。如果这一判断成立，那么中国通胀的上行风险将不仅来自能源输入，更来自内需的自我修复。

三个因素叠加的结果是：亚洲通胀不再是暂时的、供给侧的扰动，而是正在内化为一个更持久、更广泛的价格上涨周期。

在通胀压力明确化之后，DB对各经济体的政策路径给出了具体的推演。这些推演值得逐条审视，因为它们共同指向一个结论：亚洲利率的顶部尚未到来。

报告预计，印尼央行和菲律宾央行将分别再加息75个基点，使本轮紧缩总量分别达到125和100个基点。韩国央行和印度央行预计将各加息100个基点，且韩国央行将率先行动。马来西亚央行可能将原定于2027年的25个基点加息提前至今年。新加坡金融管理局预计将进一步收紧。越南国家银行仍在观察名单上，但报告指出其依赖流动性管理来平衡增长与通胀/汇率风险的做法“在高企的美国利率面前将越来越难维持”。

这些预测的背后是一个清晰的逻辑：亚洲央行正在将汇率稳定纳入政策目标函数的核心位置。报告提到，亚洲当局已经通过行政和监管措施加强了对汇率稳定的干预，包括吸引更多外资流入和遏制投机性资本流动。这意味着，美联储的政策路径将通过汇率渠道直接约束亚洲央行的操作空间。

对于投资者而言，这意味着亚洲债券的久期风险正在上升，而货币政策的差异化将为套利交易提供机会，但风险也在同步放大。

DB的报告揭示了一个容易被忽视的宏观特征：亚洲内部正在经历显著的结构分化。北亚（中国、韩国、台湾）与东盟五国在增长动能上出现了明显的差距。

报告数据显示，北亚的出口在2026年预计增长18\%，而东盟五国仅为5\%。这一差距的核心驱动力是半导体。北亚经济体深度嵌入全球半导体供应链，而AI投资的

[... middle omitted ...]

题：中国房地产底部的可持续性与全球供应链风险

尽管DB的报告提供了丰富的分析和判断，但有几个关键问题仍未得到充分解答，而这些问题的答案将直接影响上述推演的可靠性。

第一个问题是中国房地产市场的底部是否真的“持久”。报告提到数据暗示底部形成，但同时也承认“广泛的复苏尚不确定”。这里的核心矛盾在于：如果中国房地产市场的改善是由市场基本面驱动而非政策刺激驱动，那么它应该更具可持续性。但当前的数据是否足以排除二次探底的风险？尤其是考虑到居民信心修复的脆弱性和地方政府财政压力，房地产的复苏路径仍存在显著的不确定性。

第二个问题是全球供应链中断的风险。报告将“真正的供应链中断”列为出口前景的主要下行风险，并指出交货时间的恶化是一个令人担忧的信号。但报告并未深入分析这一风险的概率和潜在影响。如果中东局势进一步升级，或者美国贸易政策（如USTR 301调查）导致更广泛的脱钩，亚洲的出口韧性将面临真正的考验。半导体定价的强劲能否抵消供应链中断的冲击，这是一个尚未被充分验证的命题。

\subsection{[R009] NOM：韩元走弱的真正推手不是经济基本面，而是资本市场成功的“悖论”}
{\small\color{refgray}归类：区域市场：印度制造业竞争力缺口与韩元贬值悖论 / 外汇与资本流动：日元对冲需求与韩元资本外流悖论}

NOM：韩元走弱的真正推手不是经济基本面，而是资本市场成功的“悖论”

韩国正在经历一个让许多宏观交易员感到困惑的时刻。经济账上，出口强劲推动经常账户盈余在2026年3月飙升至373亿美元，年化占GDP比例高达23\%；KOSPI指数年内涨幅达87\%，领跑全球主要股指。按照教科书逻辑，韩元应当成为亚洲表现最强的货币之一。然而现实是，韩元兑美元和名义有效汇率年内反而贬值约5.5\%，沦为亚洲表现最差的货币之一。

这份NOM研报的核心判断是：韩元当前面临的不是基本面问题，而是一个由资本市场繁荣自身催生的、结构性的资本外流悖论。理解这一悖论的运行机制，比猜测韩国央行何时干预或美元何时走弱，更能帮助投资者把握韩元的中期走向。

报告的价值在于，它系统性地拆解了韩国国际收支中几个相互作用的资本流动引擎，并给出了三种情景下的量化压力测试。这不仅仅是一份汇率分析，更是一份关于“当国内资产价格暴涨时，资本账户会如何反向吞噬经常账户盈余”的机制说明书。

1. 经常账户盈余与货币贬值并存的悖论，根源在于资本账户的结构性失血

韩国经济目前呈现出一个典型的“K型”图景：科技出口部门极度繁荣，但这一繁荣并未有效转化为韩元的买入力量。NOM的分析揭示了背后三个关键的资金漏损渠道。

第一个渠道是企业层面的美元囤积。尽管出口强劲，韩国企业并未将大量美元收入结汇为韩元。报告引用韩国当局的表态和当地媒体报道，指出企业正在主动持有美元资产，这直接压低了实际贸易结算的净流入规模。企业这种行为背后是对全球贸易环境的不确定性和对韩元潜在贬值空间的防御性布局。

第二个渠道是零售投资者的海外资产配置。韩国散户对海外股票，尤其是美国科技股的配置热情在过去一年极为高涨。在2026年4月和5月，韩国零售投资者出现了自2023年6月以来首次连续两个月净卖出美国组合资产，但这究竟是趋势反转还是暂时停顿，报告提出了明确的质疑。即将到来的SpaceX、Anthropic和OpenAI等IPO，很可能再次点燃韩国散户的海外配置热情。

第三个渠道是机构层面，即韩国国民年金公团的海外投资。NPS作为全球最大的养老基金之一，其资产配置决策对韩元有直接且巨大的影响。2026年5月，NPS宣布将国内股票目标配置比例从14.9\%大幅提升至20.8\%，这看似是在减少海外投资。但NOM指出，这恰恰是悖论的核心：NPS的国内股票配置比例提升，是在一个已经暴涨的市场背景下做出的。如果KOSPI继续上涨，NPS为了达成新的目标比例，反而可能需要在高位减持国内股票、增持海外资产，从而形成新的资本外流压力。

这三个渠道叠加在一起，就解释了为何韩国经常账户盈余创历史新高，韩元却持续走弱。资本市场繁荣本身，正在通过企业和居民的资产负债表调整，制造出与基本面反向的资本流动。

2. NPS的资产配置调整看似利好韩元，但市场上涨正在抵消其正向效应

NPS在2026年5月修订的资产配置目标，是近期韩国市场最受关注的变量之一。按照新目标，NPS对国内资产的配置比例将从41.3\%提升至45.3\%，这意味着到2026年底，NPS对海外资产的净流出将从原先的439亿美元骤降至仅12亿美元。这看起来是对韩元的重大利好。

但NOM的测算揭示了一个隐藏的风险：这个测算假设NPS的资产规模仅按计划线性增长。然而，自2026年3月底以来，KOSPI指数已经上涨了约60\%。如果NPS的国内股票持仓跟随市场上涨，其国内股票的实际占比已经远超新目标。报告估算，在60\%的涨幅假设下，NPS可能需要重新卖出约705亿美元的国内股票，并将其配置到海外资产中，才能满足新目标。

这意味着，NPS的新目标实际上创造了一个“动态再平衡陷阱”。市场涨得越多，NPS需要卖出的国内股票就越多，从而产生的资本外流压力就越大。报告还指出，N

[... middle omitted ...]

个被市场普遍低估的资本外流风险：全球基金对单只股票10\%的持仓上限规则。如果三星电子和SK海力士继续大幅上涨，那些受规则约束的全球基金将不得不强制卖出部分持仓，以将单只股票占比控制在10\%以内。

NOM的情景分析显示，如果三星电子和SK海力士股价上涨60\%，同时大盘上涨10\%，由此引发的海外资金流出规模可能高达480亿美元。如果涨幅达到100\%，流出规模将进一步扩大至953亿美元。这组数字远超市场此前的预期。

这一风险的特殊之处在于，它不是基于基本面恶化或风险偏好下降，而是基于规则驱动的被动卖出。这意味着，即使韩国科技股的基本面依然强劲，只要股价上涨过快，就会自动产生反向的资本流出压力。这种机制与NPS的再平衡类似，都是“成功带来的烦恼”。

对于持有韩国科技股的全球投资者而言，这一分析提供了一个重要的风险警示：当股价进入快速上涨阶段时，不能只关注基本面趋势，还需要评估持仓结构是否可能触发规则性的强制卖出。这种卖出可能不是一次性的，而是随着股价上涨而阶梯式发生的。

\subsection{[R010] 摩根斯坦利：中国大宗商品出口的“结构性分化”正在重塑全球定价权}
{\small\color{refgray}归类：能源与大宗：欧洲钢铁定价权重塑与储能政策重心转移}

摩根斯坦利：中国大宗商品出口的“结构性分化”正在重塑全球定价权

这份报告的价值不在于它记录了5月的贸易数据，而在于它揭示了一个正在发生的根本性变化：中国大宗商品的出口结构正在从“数量驱动”转向“套利驱动”，而这种转变对不同品种的定价逻辑和竞争格局有着截然不同的含义。

摩根斯坦利刚刚发布的5月贸易数据追踪报告，表面上是一组月度数字的更新。但如果我们把这些数字放在一起看，一个清晰的信号浮现出来：铝和铝制品出口创下7个月新高，钢铁出口在低位企稳，而铜的进口则呈现出一种“价格高企但需求不弱”的韧性。这些看似零散的现象，背后有一条共同的逻辑线——全球供应链的再调整正在改变中国作为“世界工厂”的出口角色。

这不是一个关于“需求复苏”的故事，而是一个关于“供给再定价”的故事。理解这一点，比记住任何一个单月数据都重要。

1. 铝出口的跳升不是需求信号，而是套利窗口被打开的必然结果

5月中国铝及铝制品出口达到63.2万吨，环比增长6\%，同比增长16\%，创下自2024年11月以来的最高水平。这个数字本身并不令人意外，真正值得关注的是它背后的驱动机制。

报告明确指出，中东地区的供给中断收紧了除中国以外的市场，推高了LME铝价，从而扩大了出口套利空间。换句话说，这轮出口增长的核心驱动力不是海外终端需求的突然爆发，而是中国生产商抓住了价格差扩大的窗口。

这带来了一个重要推论：这种出口增长的可持续性取决于套利窗口能开多久，而不是海外需求有多强。一旦LME价格回落或国内成本上升压缩了价差，出口量可能会迅速回调。对于投资者而言，这意味着需要跟踪的不是简单的出口量趋势，而是内外价差的动态变化。

更深一层的含义是，中国铝行业正在从“被动接受全球价格”转向“主动利用全球价差”。那些拥有低成本产能、能够灵活调节内外销售比例的企业，将在这轮结构性变化中获得更大的议价权。

钢铁出口在5月同比下降2\%，但环比增长9\%，达到1034万吨。这个环比改善被一些市场参与者解读为需求回暖的信号。但摩根斯坦利的分析揭示了一个更复杂的图景。

报告提到，5月CISA会员钢厂的日均产量同比下降4.4\%，较4月的-5.7\%有所收窄。但更关键的数字是：摩根斯坦利估算5月中国表观钢材消费量同比下降10.4\%，环比下降11.1\%。这背后是“更多的库存积压”。

换句话说，钢厂减产的速度没有赶上需求下滑的速度，导致库存被动累积。出口的环比改善更多是消化过剩产能的一种被动选择，而不是海外需求的主动拉动。

这引出一个值得深思的问题：如果国内需求持续疲软，而出口又面临越来越多的贸易壁垒，中国钢铁行业将如何重新平衡供需？报告没有给出答案，但暗示了方向——减产可能还需要继续，而出口的竞争将更加激烈，只有成本最低、产品差异化最强的企业才能生存。

3. 铜进口的韧性说明“高价”不等于“弱需求”，但供给瓶颈正在转移

铜是这份报告中唯一一个呈现出“进口增长”的品种。5月铜及铜制品进口44.6万吨，同比下降1\%但同比增长4\%。更值得注意的是，4月净精炼铜进口量达到了2025年9月以来的最高水平。

这打破了“铜价高企必然抑制需求”的简单逻辑。报告指出，尽管铜价上涨，但在4月和5月期间，进口套利窗口间歇性打开。洋山铜溢价维持在每吨60-75美元，表明实物需求持续强劲。

但真正值得关注的信号来自上游。铜矿砂及精矿进口量环比下降1\%，同比持平，为235万吨。与此同时，加工费持续下降，表明冶炼产能相对于矿山供给仍然过剩。报告提到，强劲的硫酸价格支撑了冶炼厂的利润，但这是一种脆弱的平衡——一旦硫酸价格回落，冶炼厂的利润空间将被迅速压缩。

这意味着铜产业链的瓶颈正在从冶炼端向上游矿山端转移。那些拥有自有矿山资源的企业，将在这轮结构性变化中获得更大的竞争优势。而对于纯粹依赖外购矿的冶炼企业，加工费的持续下降将是一个长期压力。

5月中国铁矿石进口9800万吨，环比下降6\%，同比持平。港口库存开始逐步下降，但仍处于高位。这个数据本身并不惊人，但它与钢铁产量的数据放在一起看，就揭示了一个重要的结构性特征。

报告显示，钢厂生铁产量同比下降2\%，而粗钢产量同比下降4\%。这意味着废钢使用比例在上升，或者炼钢工艺在发生变化。对于铁矿石需求而言，这是一个长期的结构性下行信号。

高位库存叠加需求结构变化，意味着铁矿石价格的上行空间可能被限制。这并不是说铁矿石价格不会反弹，而是说任何反弹都需要更强劲的需求驱动，而不是简单的供给收缩。对于钢铁企业而言，原料成本的压力可能正在边际缓解，但这不足以抵消终端需求疲软带来的利润压缩。

尽管报告提供了丰富的数据和洞察，但有一个关键问题没有得到充分展开：这些出口增长的质量如何？

铝出口的增长是套利驱动的，钢铁出口是过剩产能驱动的，铜进口是结构性需求驱动的。这三种不同的驱动机制意味着完全不同的可持续性。套利窗口可以随时关闭，过剩产能的消化是痛苦的，而结构性需求则相对稳定。

但报告没有深入分析的是：这些出口中有多少是真正具有定价权的产品，有多少是低附加值的初级加工品？中国企业在海外市场的竞争力，是建立在成本优势上，还是建立在技术或品牌优势上？这些问题的答案，将决定中国大宗商品行业的长期盈利能力和估值水平。

6. 投资者需要建立的新观察框架：从“量”到“价差”再到“结构”

\subsection{[R011] GS：日本央行6月加息的可能性被市场低估，真正关键的是加息后的沟通路径}
{\small\color{refgray}归类：宏观与利率：亚洲通胀反弹与央行紧缩路径}

GS：日本央行6月加息的可能性被市场低估，真正关键的是加息后的沟通路径

市场对日本央行加息的讨论，大多停留在“6月还是7月”的择时博弈上。但GS最新一期日本经济研报提出了一个更具结构性的判断：加息时点本身不是最重要的，真正决定资产定价走向的，是央行在加息后如何修改前瞻指引，以及如何管理市场对“接近中性利率”的预期。

这份报告由GS首席日本经济学家Akira Otani领衔，在植田和男6月3日讲话后，将加息预期从7月提前至6月。这个调整本身并不令人意外——市场已经部分定价。但报告真正有穿透力的部分，在于它对加息之后路径的推演，以及对央行沟通策略的精细化拆解。这些内容，恰恰是很多短期交易视角容易忽略的。

以下是我们从报告中提炼出的五个核心层次，它们共同指向一个结论：日本正在进入一个“加息但不紧缩”的微妙阶段，市场对央行沟通的敏感度将前所未有地上升。

1. 经济数据“软硬分化”为加息提供了窗口，但真正的推力来自政治沟通的铺垫

GS的分析起点，是一个看似矛盾但已被多次验证的现象：日本经济正在经历软数据与硬数据的背离。消费者信心和景气观察调查等软指标在走弱，但机械订单、消费活动指数等硬数据依然稳健。这种“软硬分化”本身并不构成加息的充分条件——事实上，过去几年类似的分化曾多次出现，央行都选择了按兵不动。

第一个信号是5月22日植田和男与首相高市早苗的会面。GS明确指出，这次会面为加息做了政治层面的铺垫。在日本的货币政策决策框架中，政府与央行的沟通一直是一个“房间里的大象”——市场知道它重要，但很难量化。GS这次提供了一个可操作的观察框架：植田在会面后两周的6月3日讲话中，非但没有像市场预期的那样“降温”，反而释放了更鹰派的信号。GS的推断是，植田在会面中获得了政府至少不会反对加息的确认。

第二个信号是B2B价格的加速上涨。报告显示，企业间交易价格正在显著上升，尤其是石油相关产品。GS认为，这种上游价格压力将逐步传导至下游的B2C交易，从而在未来推高消费者价格。这为加息提供了经济层面的支撑。

这两个信号叠加，意味着6月加息的决策基础不仅是经济数据，更是政治可行性的确认。这一点对理解日本央行的行为模式至关重要：它不是一个纯技术性机构，其决策节奏深受与政府沟通进展的影响。

2. 加息本身不是终点，前瞻指引的修改才是市场最需要解读的变量

如果6月加息，市场会立刻转向下一个问题：然后呢？GS对此给出了一个非常具体的判断：央行将在加息的同时修改前瞻指引，核心变化是加入“尤其关注物价上行风险”的表述。

这个修改看似微小，但GS指出其用意是双重的。一方面，央行需要维持鹰派立场，防止市场过早定价加息结束。另一方面，它需要暗示政策利率正在接近中性利率，从而避免市场产生“每次会议都会加息”的线性预期。

这是一种非常精细的预期管理。接近中性利率意味着未来的加息空间在收窄，但央行又不希望市场因此放松对通胀的警惕。GS对前瞻指引修改的预测，实际上是在告诉我们：日本央行正在试图走一条“加息但不紧缩”的钢丝——加息本身会推高短端利率，但引导市场预期利率接近中性，又在一定程度上压制长端利率的上行空间。

对于投资者而言，这意味着日本国债收益率曲线的形态变化，可能比绝对水平更重要。如果央行成功引导市场相信1.5\%是终端利率，那么长端利率的上行空间将被锁定，而短端利率的抬升将导致曲线平坦化。这对不同期限的债券、银行净息差、以及日元套利交易的影响，将是截然不同的。

3. 加息节奏将从“不定期”转为“半年一次”，但政治窗口的不确定性仍是最大变量

GS在6月加息后，预计下一次加息在2027年1月，然后7月达到1.5\%的终端利率。这意味着加息节奏将从过去的“不定期”转为大约半年一次。

这个节奏预测背后有一个重要的逻辑转变：随着政策利率升

[... middle omitted ...]

迟。

这里有一个值得深思的问题：如果日本政治格局发生变化，或者首相的立场出现调整，这个沟通模式是否还能持续？报告没有展开讨论，但这是所有关注日本资产的人都需要持续跟踪的变量。

4. 国债购买缩减计划可能比市场预期的更温和，但这本身就是一个信号

关于日本国债购买计划，GS的判断是：央行可能在3月前维持现有购买计划，从2027年4月起将月度购买量维持在约2万亿日元。这个路径比一些市场参与者预期的“更快缩减”要温和。

GS没有明确解释为什么央行会选择相对温和的缩减路径，但我们可以从报告的隐含逻辑中推导：既然央行希望引导市场相信政策利率将接近中性利率，那么它就不希望长端利率因为国债供给压力而过度上行。温和的缩减计划，实际上是在配合“接近中性利率”的预期管理。

对于日本国债投资者来说，这意味着供给压力的释放是渐进的，而非一次性的。短期内，长端利率的上行空间可能受到央行购买行为的压制；但中期来看，如果通胀持续高于预期，央行可能不得不调整购买计划，届时市场可能会重新定价。

\subsection{[R012] GS：CMS释放的信号比药价谈判本身更重要}
{\small\color{refgray}归类：医疗健康：药价谈判结构性变化与设备招标拐点}

市场对药品定价改革的关注点，大多集中在“砍价幅度有多大”这个数字游戏上。但GS在其第47届全球医疗保健年会上与CMS高层官员的对话，揭示了一个更深层的结构性变化：美国医保体系正在经历一场系统性的财政责任再分配。这场再分配的最终影响，可能比任何一轮药价谈判都要深远。

GS的研报显示，CMS领导层将其近期议程锚定在四个核心优先事项上：以患者为先的支付方、预防导向的战略、打击欺诈浪费和滥用、以及可负担性。这听上去是一份标准的政策清单。但真正值得关注的，是这些优先事项背后隐含的财政逻辑——联邦政府正在系统性收紧其在医疗支出中的敞口，并将更多成本压力向下传导至州政府、制药企业和医疗服务提供方。

这不是一个孤立的事件。这是一轮医疗体系供给侧的结构性再定价。

1. 药价谈判只是表象，真正重要的是MFN机制正在重塑全球定价体系

CMS在会议上披露的数据值得反复推敲：第一轮IRA谈判带来了约20\%的平均净节省，而当前政府主导的第二轮谈判，折扣幅度已接近40\%。这个数字本身并不意外——市场对此已有预期。但GS分析师敏锐地捕捉到了一个更关键的信号：CMS将最惠国（MFN）机制视为一个“更广泛的工具包”，其目标不仅仅是美国市场，而是通过贸易动态在全球范围内施加定价压力。

这意味着什么？MFN的逻辑是，如果制药企业在美国市场给出了一个折扣价格，那么其他联邦医保项目也应该享受同样的价格。但CMS的表述暗示，他们正在将这个逻辑延伸至全球——如果一家制药公司在某个国家接受了较低的价格，那么这个价格应该成为美国市场的参照基准。

这实际上是在构建一个全球性的药品定价锚定体系。对于跨国制药企业而言，这不再是“美国市场单独承压”的问题，而是全球定价体系正在被重新定义。GS的研报没有完全展开这一点的二阶影响，但一个合理的推演是：如果MFN机制从当前的17家生物制药企业进一步扩展，那么全球药品定价的离散度将被大幅压缩，差异化定价策略的空间将被显著收窄。

这里有一个关键问题值得继续追问：CMS没有对MFN是否扩展至17家以外的企业做出表态。这个沉默本身可能就是一个信号——政策制定者正在观察第一轮扩展的效果，同时也在评估制药行业的反应。对于投资者而言，这个不确定性本身就是需要定价的风险因子。

2. 联邦医保资金流向正在发生结构性转变，州政府将成为新的成本承担者

如果说药价谈判是“看得见的改革”，那么CMS在Medicaid资金分配上的调整，则是“看不见但影响更大的变化”。

GS研报披露了一个关键数据：联邦在Medicaid资金中的占比，已从项目初期的约60\%上升至目前的约70\%。CMS官员明确表示，这被视为“联邦支出过剩”。他们的应对策略是：收紧项目诚信、重新平衡成本分摊，并逐步将Medicaid报销标准向Medicare基准靠拢，引入支付上限。

这句话翻译过来就是：联邦政府正在减少对Medicaid的补贴力度，州政府需要承担更多责任。考虑到Medicaid是许多州预算中最大的支出项目之一，这一调整将直接挤压州政府的财政空间。

更值得关注的是CMS对“提供者税机制”和“Medicaid扩展人群”的态度。这两个领域被明确标记为“吸引更多联邦资金”的来源，而CMS认为这导致了“超额联邦支出”。这意味着，过去州政府通过一些财务工具来最大化联邦匹配资金的做法，正在被系统性地收紧。

对于医疗服务提供方而言，这意味着两重压力：一方面，联邦层面的报销增速可能放缓；另一方面，州政府为了平衡预算，可能会进一步压缩对医疗服务提供方的支付。GS的研报没有直接点明，但一个清晰的推论是：那些高度依赖Medicaid收入的医疗企业，将面临更大的现金流压力。

CMS对健康保险交易所（ACA Exchanges）的关注，表面上是一个“清理门户”的动作——移除不合

[... middle omitted ...]

际效果之间，可能存在偏差。

另一个被低估的变化是CMS推动的“改善计划间竞争”。这意味着，未来ACA市场的结构可能会从“少数大型保险公司主导”转向“更多元化的竞争格局”。对于大型管理式医疗公司来说，这既是挑战也是机遇——挑战在于市场份额可能被侵蚀，机遇在于那些能够通过数据和技术实现更精准风险管理的公司，可能在新格局中占据优势。

4. 报告尚未完全回答的关键问题：IRA谈判的长期效果与创新激励的权衡

GS的研报在讨论药价谈判时，提供了一个重要的观察框架，但同时也留下了一个尚未完全回答的问题：IRA谈判的长期效果，究竟会如何影响制药行业的研发投入和创新产出？

当前的数据显示，第一轮谈判带来了约20\%的净节省，第二轮接近40\%。但这两个数字背后有一个关键变量没有被充分讨论：这些节省在多大程度上来自于“价格调整”，在多大程度上来自于“产品组合变化”？如果制药企业将资源从受IRA影响的产品转向不受影响的产品，那么所谓的“节省”可能只是账面数字的转移，而非真实的成本降低。

\subsection{[R013] GS：医疗健康板块的结构性拐点，比市场感知的更近}
{\small\color{refgray}归类：医疗健康：药价谈判结构性变化与设备招标拐点}

过去三年，全球医疗健康投资界一直在等待一个明确的信号：基本面改善何时能转化为估值修复。GS在第47届全球医疗健康年会上组织的一场投行家圆桌讨论，给出了一个值得认真对待的答案——这个拐点可能比多数人预期的更近，但它不会以整齐划一的方式到来。

这份讨论纪要的核心信息不是“医疗健康板块要涨了”，而是一个更微妙的判断：板块内部正在发生深刻的分化，决定胜负的关键不再是赛道选择，而是企业在资本配置、资产组合和运营效率上的执行力。那些能够将规模转化为议价权、将数据积累转化为竞争壁垒的企业，将在下一轮周期中占据显著优势。

这场讨论的价值不在于提供了多少新数据，而在于它揭示了市场定价与基本面之间正在形成的裂口——以及这个裂口如何为有准备的投资者创造机会。

2026年至今，全球医疗健康并购交易额已接近1000亿美元，而2025年全年仅为500亿美元。这个数字本身已经足够引人注目，但真正值得关注的是背后的驱动力。

投行家们观察到，本轮并购活跃度并非简单来自低利率或充裕资本，而是源于产业内部的结构性压力。大型生物科技公司正在主动追逐中型标的，且交易模式发生了显著变化：买方更早介入谈判、更频繁使用或有价值权（CVR）来管理风险、更专注于大型药企通常忽略的细分治疗领域。

这意味着什么？传统上，大型药企的并购逻辑是“补品种”——用现金换收入。而当前这轮并购的底层逻辑是“补能力”——用交易换效率。大型药企意识到，内部研发的回报率正在下降，而通过并购获取经过验证的、执行风险较低的商业化阶段资产，是提升资本回报率的最可靠路径。

对投资者而言，这意味着需要重新评估哪些公司最可能成为并购目标。不再是单纯看管线深度，而是看资产的商业化成熟度、与买方现有组合的协同效应、以及管理团队在交易后的整合能力。GS的M\&A排名框架中，那些被归类为“高概率”（30\%-50\%）的标的，其价值重估可能才刚刚开始。

2. 中国生物科技的“双刃剑效应”：既是威胁，也是效率催化剂

中国生物科技产业是这场讨论中最具争议性的话题之一。投行家们直言，这是一个“高度两极分化”的议题——一部分人将其视为生存威胁，另一部分人则看到高质量的创新、低成本资本和临床试验的高吞吐能力。

但真正值得注意的洞察是，中国生物科技的崛起正在迫使美国创新企业提升运营效率。这不是一个零和博弈，而是一个竞争性催化过程。中国公司以更低的成本、更快的速度推进临床，这直接挑战了美国企业的成本结构和时间表。那些无法在效率上做出回应的美国公司，将面临持续的估值压力。

同时，双方存在真实的利益交集：中国实体需要进入美国市场，而美国战略投资者看到了以合理价格获取高质量资产的机会。这意味着跨境交易将继续活跃，但交易结构将更加复杂，监管审查也会更严格。

对产业决策者而言，这意味着不能再用“中国威胁论”来回避效率问题。真正需要回答的是：你的研发运营模式能否在更低的成本基础上保持创新质量？如果不能，你的资产在未来两年内的吸引力将显著下降。

3. AI在医疗健康中的落地：效率提升真实存在，但商业转化仍需验证

AI是这场讨论中反复出现的主题，但投行家们的表述比市场叙事更为谨慎。他们承认，AI已经在行政性工作中产生了可衡量的效率提升——法律文件处理、监管申报、实验室流程优化等领域都有明显改善。但同时，AI在核心业务环节的落地面临结构性约束。

在生物制药领域，AI的采用速度受到“责任速度”的限制——在涉及患者安全的产品开发中，技术的推进必须与风险管控同步，这天然限制了AI的应用节奏。在医疗设备领域，AI集成面临医生信任建立和产品审批的双重挑战，这“都需要时间”。

在生命科学工具和诊断领域，AI的角色更为明确：优化实验室效率、加速数据驱动的吞吐、在诊断中捕获纵向数据以构建数据壁垒。这里的关键洞察是，AI的真正价

[... middle omitted ...]

原因是什么？核心在于定价权的持续悬而未决。投行家们提到了“存在性的定价压力”，这使得市场难以对终局价值形成稳定预期。投资者无法确定五年后的定价环境，因此不愿意为当前的收入增长支付溢价。

但这里有一个值得注意的信号：投行家们预计，随着大型企业重置市场预期并建立起“超越并上调”的业绩模式，板块动能可能回归。这意味着，一旦有一两家龙头企业通过连续的超预期表现打破了市场对定价的悲观预期，整个板块的估值框架可能迅速重估。

对于投资者而言，医疗设备板块的核心问题不是“基本面好不好”，而是“市场何时开始重新定价”。这需要关注两个先行指标：一是大型企业是否开始建立稳定的业绩超预期模式，二是关于定价政策的关键监管信号是否出现。报告没有给出具体时间表，但结构性的估值错位本身就是一个值得关注的信号。

欧洲医疗健康板块的讨论揭示了一个被多数投资者忽视的结构性问题：规模瓶颈。投行家们指出，欧盟虽然是可靠的创新来源，但许多公司面临因资本获取受限导致的规模化障碍——尤其是相对于美国市场而言。

\subsection{[R014] Citi：中国2万亿AI投资计划的核心不在于规模，而在于国产化率的硬约束}
{\small\color{refgray}归类：AI/算力：资本开支上行风险与供应链结构性重置}

Citi：中国2万亿AI投资计划的核心不在于规模，而在于国产化率的硬约束

当市场还在讨论中国AI算力缺口有多大时，一份来自Citi研报的最新信号值得决策者认真对待：中国计划在未来五年投入2万亿元人民币建设数据中心和网络基础设施，以支撑国内AI产业发展。这个数字本身并不令人意外——市场早已预期中国会在AI基础设施上持续加码。真正被低估的，是这份计划中一个看似技术性的硬性约束：电信运营商运营的数据中心，必须使用至少80\%的国产AI加速器。

这个80\%的国产化率目标，才是理解未来五年中国半导体产业链投资逻辑的关键。它不仅仅是一个采购指引，更是一个产业格局的再分配机制。Citi的研报正是在这个判断基础上，给出了对代工厂、封测、设备以及AI芯片供应商的正面看法。

这份报告发布的时间点也值得玩味。当前，中国AI芯片市场已经发生了结构性变化：英伟达H200的进口基本停滞，国产AI芯片实际上已经占据了市场的绝大多数份额。这意味着，80\%的国产化目标并非遥不可及的远景，而是对已经发生的市场现实的确认和制度化。但问题在于，这种确认对不同环节、不同体量的企业意味着完全不同的东西。

Citi研报中最具洞察力的判断之一，在于它对国产化率目标实现路径的评估。报告明确指出，80\%以上的AI芯片国产化目标是“可实现的”，因为国产AI芯片厂商今年实际上已经有效占据了市场的绝大多数份额。但这并不意味着所有国产厂商都能均分这块蛋糕。

关键在于，国有资本投资的数据中心在采购时，天然倾向于更广泛的供应商选择。这与过去市场集中采购少数头部芯片的逻辑有本质区别。Citi认为，这种变化对中小型AI芯片厂商更为有利——国有数据中心更愿意从更广泛的供应商处采购。

这意味着，竞争格局将从“赢家通吃”转向“多极共存”。过去，市场关注的焦点是华为昇腾能否追上英伟达的算力水平。但在80\%国产化率的框架下，竞争维度发生了变化。不再是谁的算力绝对最强，而是谁能够稳定供货、满足合规要求、并在特定场景下提供足够的性能。这为寒武纪、壁仞、摩尔线程、MetaX等二线厂商打开了结构性机会窗口。

从Citi提供的AI加速器对比表中，我们可以清晰地看到性能梯队的分布。华为的910C和920系列在FP16性能上已经达到0.8-0.9 PFLOPS，接近英伟达H20的0.3 PFLOPS的2-3倍。但更值得注意的是，壁仞的BR100和华为950系列的FP16性能已经达到1.0 PFLOPS，与英伟达H20的0.3相比，差距在缩小。虽然与H200的2.0 PFLOPS仍有明显差距，但在80\%国产化率的保障下，这种差距不再构成市场准入的障碍。

2. 代工厂和封测环节将受益于更确定的需求预期，但估值已经部分反映

Citi研报将代工厂和封测列为直接受益者，逻辑清晰：国产AI芯片出货量的增加，必然拉动对先进制程代工和先进封装的需求。中芯国际和华虹半导体作为国内主要的代工平台，将获得更稳定的订单预期。封测环节的长电科技和通富微电同样受益。

但这里需要做一个冷静的区分。从Citi提供的板块股价表现数据来看，过去12个月，代工板块上涨158.9\%，封测板块上涨103.9\%，前端设备上涨127.5\%，后端设备更是暴涨306.8\%。这些数字说明，市场已经对国产替代的逻辑进行了相当程度的定价。

真正的增量在哪里？Citi的框架指向的是“更确定的需求预期”和“更清晰的路线图”。过去，国产AI芯片的出货量受制于政策不确定性、客户接受度和供应链瓶颈。现在，2万亿投资计划和80\%国产化率目标，相当于为未来五年的需求画出了一条底线。对于代工厂和封测厂而言，这意味着产能利用率的下行风险被显著降低，资本开支的回报预期更加清晰。

但投资者需要警惕的是，当估值已经处于高位时，确定性改善的边际价值在下降。真正

[... middle omitted ...]

.8\%，前端设备上涨127.5\%，差距超过一倍。

这种分化背后有深刻的产业逻辑。后端设备（封测设备）的国产化难度相对较低，技术壁垒和客户验证周期较短，因此在国产替代浪潮中率先受益。而前端设备（晶圆制造设备）涉及光刻、刻蚀、薄膜沉积等核心工艺，技术门槛极高，客户验证周期长达数年。

Citi将设备供应商列为受益方，但并未详细展开不同环节设备国产化的时间差。这里需要补充一个判断：2万亿投资计划对设备的拉动，将呈现明显的“后端先行、前端后至”的特征。短期内，封测设备的订单弹性最大；中长期看，前端设备的国产化突破才是决定中国半导体产业链自主可控深度的关键。报告对此着墨不多，这恰恰是需要投资者自行深度挖掘的方向。

4. 报告尚未完全回答的关键问题：算力性能差距如何在不影响AI应用落地的情况下被消化

Citi的研报提供了一个清晰的受益框架，但它也留下了一个核心问题没有完全回答：当国产AI芯片与英伟达最先进产品的性能差距仍然显著时，2万亿投资能否真正转化为有竞争力的AI应用？

\subsection{[R015] Bernstein：量子计算的真正变量不是技术路线，而是速度与保真度的取舍}
{\small\color{refgray}归类：新兴科技：AI音乐变现与量子计算速度精度取舍}

Bernstein：量子计算的真正变量不是技术路线，而是速度与保真度的取舍

量子计算行业正站在一个关键的十字路口。市场上充斥着“量子霸权”的叙事和各类技术路线的争论，但真正决定产业走向的核心问题，并非哪种量子比特最终胜出，而是一个更根本的取舍：你是愿意在毫秒级完成计算但容忍一定误差，还是愿意等待一秒以上但确保结果的精确性？

Bernstein最新发布的量子计算行业专家研讨会纪要，揭示了这一被多数投资者忽视的核心矛盾。该机构邀请了一位拥有15年行业经验、横跨硬件、系统和软件应用全栈的量子技术管理者，系统拆解了当前量子计算的技术格局与竞争态势。这份纪要的价值不在于罗列技术参数，而在于给出了一个可操作的判断框架：不同应用场景对计算速度和保真度的要求截然不同，这意味着未来的量子计算不会是“赢家通吃”，而是一个多模态共存的生态系统。

对于产业决策者和投资者而言，理解这一框架，远比押注某一技术路线更为重要。

1. 量子计算的核心张力不在于技术路线，而在于“制造”与“天然”的物理本质差异

大多数市场讨论将量子计算技术路线简单归类为“谁更快”或“谁更准”，但这种比较忽略了最根本的物理学差异。Bernstein专家提出的分析框架极具洞察力：所有量子比特应被分为“制造型”和“天然型”两大类。

制造型量子比特，包括超导、硅自旋和拓扑等路线，本质上是人类工程设计的产物。它们的优势在于速度快——操作频率可达千赫兹甚至兆赫兹级别，计算重复速率远超其他方案。但代价是“一致性”缺失：即使是精心制造的两个超导量子比特，其特性参数仍存在约10\%的差异。这导致控制系统的复杂度大幅上升，同时噪声更高、相干时间更短。

天然型量子比特，包括离子阱和中性原子，则利用了自然界中本就存在的完美一致性。每一个原子在物理上完全相同，这意味着控制架构的设计可以更加可靠。它们天然具有长相干时间和强抗噪声能力。但代价是速度缓慢——这类系统的计算重复速率通常不超过1赫兹，即每秒只能完成一次计算。

这一区分绝非学术概念。它直接决定了不同技术路线在商业应用中的竞争力边界。制造型量子比特擅长需要大量重复采样以估算连续数值的问题，如化学模拟和材料科学；天然型量子比特则更适合需要一次性给出离散正确答案的问题，如密码学和组合优化。

2. 当前格局是“三强并立”，但中性原子的崛起揭示了技术格局的动态性

Bernstein专家明确指出，目前量子计算领域存在三个领先的技术路线：超导、离子阱和中性原子。但这一格局并非静态。

超导目前综合领先，可视为“1A”梯队。它在速度、规模和工业成熟度上占据绝对优势，背后有IBM、谷歌等大型科技公司持续投入。其量子比特数量已达数百个，门操作的保真度和可重复性在过去五年显著提升。但超导的长期前景存在隐忧：由于制造型量子比特存在于平面结构上，实现非近邻的长程交互极为困难，而这恰恰是构建高效纠错逻辑量子比特的关键。

离子阱和中性原子则天然具备实现长程交互的能力。中性原子的崛起尤其值得关注：就在三年前，这一路线还处于行业末流，但哈佛和MIT的合作团队及其衍生公司QuEra在近三年内取得了革命性进展。专家特别强调，目前只有QuEra团队能在中性原子上实现如此优异的表现，而超导和离子阱领域则有多个学术和工业团队同时展现良好性能。

这一动态格局对投资者的启示是：技术路线的领先地位并非一成不变。中性原子的快速崛起证明，量子计算领域的“黑马”可能来自任何方向。押注单一技术路线具有显著风险。

3. 应用场景决定技术选择：速度与保真度的取舍将塑造市场分层

量子计算并非一个统一的市场。不同的应用场景对计算特性的要求截然不同，这将直接决定哪种技术路线在哪些领域占据优势。

对于化学模拟和材料科学这类需要估算实数值的应用，关键在于能否在合理时间内完成大量

[... middle omitted ...]

投资者而言，关键不是判断“哪种技术最好”，而是“哪种技术最适合哪些应用”。

4. 大型科技公司在资本密集型赛道上占据结构性优势，但初创企业仍有生存空间

量子计算是极其资本密集型的行业。Bernstein专家判断，云服务模式将成为主流，而非企业自建量子计算机。这一判断暗示了行业竞争格局的关键特征。

IBM和谷歌等大型科技公司凭借资本、基础设施和人才储备，在整体竞争中占据结构性优势。IBM尤其值得关注：其Qiskit软件生态已成为行业事实标准，在软件和生态系统层面建立了强大护城河。但专家也指出，IBM的硬件领导地位正受到天然型量子比特系统（尤其是离子阱和中性原子）的挑战，因为后者在量子比特质量上已展现出显著优势。

初创企业的生存策略则需要更加聚焦。与大型科技公司正面竞争并非明智选择。专家认为，初创企业的成功更可能来自三个路径：技术专精（在某一特定技术路线上做到极致）、战略合作（与大型科技公司或行业客户形成互补）、以及政府资助项目（利用国防和科研需求获得稳定资金）。

\subsection{[R016] DB：中国车企在欧洲的渗透已从“试探”进入“阵地战”阶段}
{\small\color{refgray}归类：中国新能源车：渗透率突破60\%后的竞争结构重构}

2026年4月，中国车企在欧洲市场的单月销量突破16.3万辆，同比增长41.7\%，市场份额首次稳定在12\%的关口。这个数字本身已经足够引起关注，但DB这份最新月度监测报告真正值得深读的信号，不是总量增长，而是增长结构的质变——中国车企正在从过去依赖俄罗斯市场的单一支点，转向对西欧核心市场的全面渗透。

这份报告以月度追踪数据为基础，覆盖了欧洲主要国家的中国品牌销量。数据本身是冰冷的，但当我们将这些数字放在一起看时，一个清晰的判断浮现出来：中国车企的欧洲战略已经完成了“市场测试”阶段，进入了需要重资产投入、本地化运营和品牌建设的“阵地战”阶段。对于产业决策者和投资者而言，这意味着评估框架需要从“能不能卖出去”切换到“能不能留下来”。

1. 俄罗斯市场的收缩正在被西欧市场的爆发所对冲，中国车企的欧洲版图正在重新绘制

这是这份报告最核心的结构性变化。2025年之前，中国车企在欧洲的销量高度依赖俄罗斯市场。但2026年前四个月的数据显示，中国车企在俄罗斯的销量同比下降13.8\%，而同期在英国的销量飙升127.7\%，意大利增长129.5\%，德国增长139.0\%。

这不是简单的“东方不亮西方亮”。俄罗斯市场的收缩有地缘政治和制裁的叠加因素，而西欧市场的增长背后是中国车企产品力、渠道布局和品牌认知的综合提升。DB的数据清晰显示，4月份中国车企在西欧核心市场的销量占比已经超过了俄罗斯——英国、意大利、西班牙、德国、法国五国合计销量达到7.76万辆，而俄罗斯为5.33万辆。

这意味着什么？中国车企正在完成一次市场重心的战略转移。过去几年，俄罗斯是中国品牌出海的“练兵场”和“利润池”，但这一池子正在变浅。真正决定中国车企海外命运的市场，现在转移到了西欧。这不是一个渐进的变化，而是一个需要重新评估企业海外资产配置和供应链布局的拐点。

2. 头部车企的竞争格局正在分化：规模优势开始转化为市场地位，而非所有玩家都能跟上

从车企层面看，DB的数据揭示了明显的梯队分化。奇瑞以4.08万辆的单月销量位居第一，但它的增长主要依赖俄罗斯市场——在俄罗斯销量占比仍然较高。真正值得关注的是比亚迪和SAIC（上汽）的表现。

比亚迪4月在欧洲销量2.84万辆，同比增长57\%，前四个月累计10.76万辆，同比增长94.3\%。这个增速意味着比亚迪正在欧洲市场复制其在中国市场的渗透路径，而且速度更快。SAIC以2.98万辆的月销量紧随其后，但增速（35.7\%）明显低于比亚迪。更值得注意的是，SAIC前四个月累计增速仅为9.8\%，远低于行业平均的40.4\%，显示其增长动能正在放缓。

另一个值得关注的变量是零跑（Leapmotor）。虽然绝对销量不大，但415.8\%的同比增速和590\%的前四个月累计增速，使其成为增速最快的中国品牌。零跑与Stellantis的合资公司正在发挥作用，这种“借船出海”的模式是否可持续，将是未来几个季度的观察重点。

从竞争格局看，中国车企在欧洲市场的“马太效应”正在显现：头部企业通过规模效应降低物流、认证和渠道成本，而小体量玩家（如蔚来、赛力斯、北汽）的销量仍在低位徘徊，甚至出现同比下滑。这提示了一个残酷的现实：欧洲市场不会给所有中国品牌机会，最终能留下来的可能只有3-5家。

3. 英国和意大利正在成为中国品牌渗透率最高的西欧市场，德国市场的突破仍需要时间

分国家来看，中国车企的市场渗透率差异巨大。英国是中国品牌渗透率最高的西欧市场，4月份达到17\%，前四个月平均水平在13\%以上。意大利紧随其后，4月份渗透率达到13\%。这两个市场的共同特点是：消费者对价格敏感度较高，对新品牌的接受度相对开放，且当地传统车企的竞争壁垒不像德国那样坚固。

德国市场则呈现另一番景象。尽管中国车企在德国的销量增速很快（4月同比增长130

[... middle omitted ...]

问题，但数据本身已经暗示：德国市场可能是中国车企在欧洲最需要耐心的地方。

4. 报告尚未完全回答的关键问题：关税和本地化生产的博弈将如何重塑竞争格局

DB的这份报告是数据驱动的，它忠实地记录了“发生了什么”，但对于“接下来会发生什么”保持了克制。这正是我们需要追问的地方。

欧盟对中国电动车的反补贴关税已经落地，但报告中的销量数据尚未充分反映这一政策的影响。关税对价格敏感型消费者的冲击是直接的，但中国车企的反应速度也很快：比亚迪已经在匈牙利建厂，奇瑞和吉利也在评估欧洲本地化生产方案。问题在于，本地化生产需要18-24个月的周期，在这个窗口期内，关税如何影响销量增速，报告没有给出答案。

另一个未解问题是：中国车企在欧洲的盈利状况如何？销量增长是确定的，但为了抢占市场份额，中国车企在欧洲的定价策略是否侵蚀了利润？DB的报告没有提供单车利润或毛利率数据，而这些恰恰是判断“增长质量”的关键。如果销量增长是靠大幅降价和巨额营销费用换来的，那么这种增长的可持续性需要打折扣。

\subsection{[R017] GS：中国餐饮的五月数据揭示的不是需求疲软，而是品牌分化的加速}
{\small\color{refgray}归类：地产与消费：中国餐饮品牌分化加速与家电估值逻辑切换}

GS：中国餐饮的五月数据揭示的不是需求疲软，而是品牌分化的加速

GS最新发布的餐饮月度追踪报告显示，五月同店销售增长（SSSG）较四月进一步走弱。五一假期需求温和，消费情绪未见明显提振，部分地区还受到天气扰动。如果你只看这个开头，很容易得出一个笼统的判断：消费复苏又踩了一脚刹车。

但这份报告的真正价值，不在于确认“五月数据变差”这个已知事实，而在于它揭示了一个更为关键的结构性变化：品牌之间的分化正在加速，且这种分化并非简单的“大品牌抗跌、小品牌承压”，而是由商业模式、品类属性和成本结构共同决定的。那些仅仅用“消费降级”或“需求疲软”来解释一切的人，正在错过真正的投资信号。

报告中最值得关注的判断是：五月数据确认了高基数效应开始显现，但真正决定企业未来走势的，不是需求总量的波动，而是品牌能否在分化中构建起新的定价权和成本优势。 这一点，在现制茶饮（FMD）和火锅两个赛道上表现得尤为突出。

1. 现制茶饮的高基数压力只是表象，真正的变量是品牌能否用新品和品类扩张对冲下滑

五月，现制茶饮品牌普遍面临同店销售增长的压力，这在GS的预期之内。去年同期的强劲表现形成了一个较高的比较基数，高基数效应开始显现。但如果我们把目光从“同比”移开，观察环比和绝对量，会发现更值得关注的信号。

奈雪的茶五月单店销售额同比下滑幅度从四月的-11\%扩大至-24\%。这个数据本身令人担忧，但更关键的是，它与其他头部品牌形成了鲜明对比。报告指出，头部品牌展现出更强的韧性，支撑力来自品类扩张和新产品。这意味着，在同样的宏观环境下，能否持续推出有吸引力的新品、能否在品类上找到新的增长点，正在成为区分品牌强弱的核心能力。

这里需要警惕一个常见的误读：把高基数压力等同于“行业不行了”。实际上，高基数压力是周期性的、可预期的，而品牌的产品创新能力和品类扩张能力是结构性的、可积累的。后者才是决定企业在下一个周期中位置的关键。对于投资者而言，五月的压力测试恰恰提供了一个窗口，去观察哪些品牌真正具备穿越周期的产品力。

2. 火锅赛道的复苏并非整体性的，海底捞的“低基数红利”正在消退

海底捞五月平均翻台率同比转为小幅负增长，而四月还是中单位数增长。GS估算，这对应的翻台率约为3.8倍，恢复至2019年水平的80\%-85\%，略高于四月，但仍低于一季度水平。管理层将部分原因归咎于端午节日历效应（2025年端午节在五月，而2024年不在），但即使剔除这一因素，节后翻台率也出现了小幅减速。

GS的判断很克制：“市场预期也已经有所下调。”这句话背后的含义是：市场此前对海底捞的期待，部分建立在对“低基数红利”的线性外推之上。但现实是，低基数的红利正在消退，而真正的需求恢复并没有跟上。海底捞一季度至今的翻台率同比增幅仅为低单位数，略低于GS全年4\%单店销售增长的假设。

这引出一个更本质的问题：当低基数效应完全消失后，海底捞靠什么驱动增长？是继续下沉开店，还是通过产品创新和场景拓展提升客单价？报告没有给出明确的答案，但数据本身已经提出了这个问题。对于投资者来说，五月的翻台率数据可能不是终点，而是一个需要重新审视增长假设的起点。

3. 太二的韧性值得关注，但它的成功模式能否复制到其他品牌仍存疑

在多数品牌承压的五月，太二的表现是一个亮点。太二在中国大陆市场录得低双位数同店销售增长，虽然略低于四月，但仍高于一季度10.9\%的水平。更值得注意的是，节后增速出现了环比改善，五一期间受南方阴雨天气影响仅为中单位数至高单位数增长。

GS将太二的韧性归因于新模型门店的持续驱动和老模型门店的正增长。这里的“新模型”是指太二近年推出的面积更小、模型更轻的新店型，在成本结构和运营效率上都有优化。这实际上指向了一个重要的战略判断：在消费环境偏弱时，品牌扩张的核心竞争力正在从

[... middle omitted ...]

周9.9元优惠券，蜜雪冰城推出了咖啡机门店的买一送一活动，瑞幸也推出了每日9.9元精选单品。乍看之下，这似乎是价格战升级的信号。

但GS的判断是：“促销力度和规模仍然是有纪律的。”这是一个非常重要的定性判断。如果促销是无序的、恶性的，那么行业利润将被迅速侵蚀，所有参与者都会受损。而“有纪律”的促销，意味着品牌在通过促销教育消费者、拉动堂食的同时，仍然在控制成本底线和品牌定位。

这里的洞察是：真正需要警惕的不是促销本身，而是促销是否导致了定价体系的崩塌。从现有数据看，九毛九和海底捞的定价趋势仍然稳定。如果后续数据出现定价松动，那将是一个比同店销售下滑更危险的信号。目前，这个风险尚未兑现，但需要持续跟踪。

5. 报告尚未完全回答的关键问题：当高基数效应和促销常态化叠加，行业利润率的底部在哪里？

GS的这份报告提供了丰富的高频数据，但有一个问题它没有完全展开：在可预见的未来，当高基数效应持续、促销活动常态化、而消费者价格敏感度仍在上升时，餐饮行业的利润率底部在哪里？

\subsection{[R018] 摩根斯坦利：市场对800V DC延迟的担忧被夸大了，真正的信号是供应链加速而非推迟}
{\small\color{refgray}归类：AI/算力：资本开支上行风险与供应链结构性重置}

摩根斯坦利：市场对800V DC延迟的担忧被夸大了，真正的信号是供应链加速而非推迟

2026年6月9日，一则来自SemiAnalysis的报道在市场引发震动：NVIDIA原生800V DC电源架构的大规模出货被推迟至2028年，远低于市场预期。消息传出后，相关电源与散热供应链个股出现波动。但摩根斯坦利在Computex期间的供应链核查给出了截然不同的结论。

这份6月10日发布的研报，虽然篇幅不长，却直接回应了市场中最大的一个争议点——AI基础设施的电源架构升级节奏是否已经放缓。我们的判断是：市场对这条消息的解读出现了方向性错误。真正的信号不是延迟，而是加速。800V DC的产业化路径不仅没有中断，反而在多个维度上比预期走得更快。

1. 800V DC的延迟消息与摩根斯坦利在Computex的一手供应链信息完全相悖

SemiAnalysis的报道声称NVIDIA已将原生800V DC电源架构的大规模出货推迟至2028年。但摩根斯坦利分析师Sharon Shih在Computex期间的直接沟通中获取的信息截然不同：NVIDIA在GTC Taipei明确表示，800V DC电源机架（Power Rack）将在2026年第三季度准备就绪并进入量产。这一时间表与“推迟至2028年”的说法存在两年的差距。

更关键的是，摩根斯坦利的供应链核查还发现了一个被市场忽视的趋势：原本计划走+-400V DC路线的多家超大规模云服务商，已经在将研发方向转向800V DC。这意味着，800V DC不仅没有延迟，反而正在成为行业共识的下一代标准。+-400V DC原本被视为过渡方案，如今其资源正在被重新分配至800V DC。

这里需要指出一个市场认知的盲区。许多人将“800V DC”等同于“NVIDIA的800V DC”，但电源架构的产业生态远比单一芯片厂商复杂。即使NVIDIA内部某个项目的时间表有所调整，也不代表整个800V DC产业链的节奏放缓。摩根斯坦利核查的是整个供应链——从电源机架制造商到保护器件供应商，再到标准认证机构——而不仅仅是NVIDIA一家。

2. 台达电将在2026年第四季度成为首家量产800V DC独立电源机架的厂商，客户是美国超大规模云商

这份研报中最具冲击力的信息点在于：台达电（Delta Electronics）将在2026年第四季度向一家美国超大规模云商批量交付800V DC独立电源机架。这意味着，在NVIDIA的800V DC方案大规模出货之前，产业链中已经有一家核心供应商率先实现了从研发到量产的跨越。

摩根斯坦利明确指出，台达电将是“first vendor to mass produce 800 VDC standalone power racks”。这一判断的意义在于：电源架构的升级并非完全绑定于GPU芯片的迭代周期。超大规模云商在建设AI集群时，可以独立于NVIDIA的芯片路线图，先行部署800V DC的电源基础设施。这类似于数据中心在GPU到位之前先完成液冷系统的部署——基础设施的前置投资是可行的，也是正在发生的。

当然，摩根斯坦利也坦承，初期的出货量不会很大。800V DC的保护器件开发和工业安全标准仍需要时间，UL等主要监管机构尚未完成相关标准的最终定义。但“量产”这个动作本身，已经打破了市场关于“800V DC遥遥无期”的叙事。

3. 800V DC的产业化路径已有清晰的路线图，从660kW到4.8MW+分三个阶段推进

摩根斯坦利在研报中附了一张极为关键的表格，展示了800V DC设备就绪度的完整路线图。这份路线图分为三个阶段：

第一阶段（2026年第三季度）：Option 1A——Power Rack（660kW），部署就绪。这是最基础的电源机架形态，也是台达

[... middle omitted ...]

。

此外，表格中还列出了几个关键支撑领域的进展：接地与保护架构正在与UL和OEM厂商协同收敛；设备认证路径已经建立并正在推进；集电设备方面，MCCB已经可用，SSCB生态系统正在成熟；高压母线、连接器和线缆方面，标准化125A架构正在开发中。这些细节表明，800V DC的产业生态正在全方位推进，而非某个单一环节的突破。

4. 报告尚未完全回答的关键问题：800V DC的保护器件和工业安全标准何时完成定义？

尽管摩根斯坦利的供应链核查给出了积极的信号，但研报本身也坦承了当前的不确定性。最核心的未解问题集中在两个领域：

第一，保护器件的成熟度。800V DC系统需要全新的故障隔离和保护机制。传统的交流保护器件无法直接复用，而固态断路器（SSCB）虽然生态系统正在成熟，但距离大规模商用仍有距离。摩根斯坦利提到“MCCB available; SSCB ecosystem maturing”，这意味着目前可用的保护器件仍以传统断路器为主，下一代固态断路器尚未完全就绪。

\subsection{[R019] 美国银行：市场正在经历一场科技股信心的结构性重置，而非简单的避险}
{\small\color{refgray}归类：风险偏好：科技股信心结构性重置与亚洲动量泡沫}

美国银行：市场正在经历一场科技股信心的结构性重置，而非简单的避险

上周，标普500指数录得自2025年4月以来最大单周跌幅，跌幅达2.6\%。在这样一次看似寻常的市场回调背后，美国银行（BofA）的最新客户资金流报告揭示了一个远比“避险”二字更为深刻的结构性信号：机构客户正以有数据记录以来最大的规模抛售科技股，单周个股净流出高达142亿美元。这并非一次战术性减仓，而是一场正在发生的、对科技资产定价逻辑的集体重估。

这份报告的核心价值不在于告诉我们“谁在卖”，而在于揭示“卖的方式”和“卖的同时在买什么”。它指向一个正在成型的市场新叙事：资金正从过去几年主导市场的“科技+增长”单极化配置，转向一种更为分散、更注重估值保护与现金回报的多元格局。对于产业决策者和高净值投资者而言，理解这一资金流向的底层逻辑，远比预测下周指数涨跌更为重要。

1. 机构客户创纪录的科技股抛售，本质是对“规模溢价”的重新定价

美国银行报告中最引人注目的数据点，是科技板块出现了自2008年有数据记录以来的最大规模资金流出。若以占板块市值百分比衡量，这一流出规模也是2014年初以来最极端的。卖盘的主导力量是机构客户，但对冲基金和私人客户同样在抛售。这不是某一类资金的孤立行为，而是整个专业投资者群体的共识性撤退。

过去两年，以Magnificent Seven为代表的科技巨头，其股价上涨的核心驱动力并非盈利增长，而是投资者愿意为“确定性”和“规模”支付的溢价。当宏观环境从通胀转向增长放缓，当AI商业化路径从“讲故事”进入“验证收入”阶段，这种溢价的基础开始松动。机构客户此刻的大规模抛售，可以理解为对“科技股是否仍应享有相对于其他板块的估值特权”这一问题的集体否定回答。

这轮抛售的启示在于：市场并非不看好科技行业的长期前景，而是不再接受用未来三年的预期利润来为今天的股价买单。科技公司的估值，正在从“增长故事”回归到“现金流折现”的基本定价框架。

2. 资金并非逃离股市，而是进行一场“大轮动”：从小盘到价值，从科技到医疗

如果只看科技股流出，很容易得出“市场悲观”的结论。但美国银行的数据显示，客户的整体行为远比表面复杂。他们在抛售科技个股的同时，却在连续第11周净买入股票ETF。卖出集中在大型股，而小型股和中型股则获得了净买入。ETF层面，客户买入价值型和混合型ETF，同时卖出成长型ETF。

这构成了一个清晰的轮动图谱：资金正从“大盘科技成长”流向“小盘价值”和“小盘成长”。更值得关注的是行业层面的分化。医疗保健ETF录得了自2021年10月以来最大的单周资金流入。在11个行业板块中，只有工业、房地产和公用事业三个板块获得了个股层面的净买入。

这些信号叠加在一起，指向一个判断：市场正在进行一场预期中的“防御性轮动”，但防御的方向并非简单的“买债券”或“持有现金”，而是转向那些估值更合理、现金流更稳定、且具备独立于科技周期的增长逻辑的资产。医疗保健、公用事业和工业，恰好符合这些特征。这轮轮动的深度和持续性，将取决于宏观数据的进一步演变。

3. 回购放缓但未消失，金融和可选消费正在接力科技成为回购主力

报告另一个值得深挖的维度是公司回购行为的变化。美国银行的客户回购数据在最近一周连续第二周放缓。然而，滚动四周平均值却升至3月底以来的最高水平。这组看似矛盾的数据，揭示了一个季节性规律：回购活动在财报季前后通常会放缓，但在非财报窗口期会回升。

更关键的结构性变化在于回购的行业分布。报告明确指出，科技板块的回购减速最为显著，而可选消费、金融和能源板块的回购则有所提速。金融板块已成为近期回购的最大贡献者，可选消费紧随其后。

这对于投资者意味着什么？回购是公司管理层对自身股价信心的直接表达。当科技公司减少回购，而金融和消费公司加大回购，这暗

[... middle omitted ...]

最核心的问题保持了沉默：这轮机构抛售何时会停止？抛售的“目标价格”或“估值水平”是什么？

报告显示，机构客户在此前连续五周净买入后，突然转为创纪录的卖盘。这种行为的急剧转向，通常意味着触发因素是宏观层面的预期变化（比如对经济增长或利率路径的重新评估），而非个股基本面恶化。如果这一判断成立，那么抛售的“终点”将取决于宏观叙事何时企稳，而非科技公司自身的财报表现。

另一个未解之谜是：私人客户（Private Clients）在历史数据中一直是净卖出者，但过去12个月却成为净买入者。报告指出，私人客户上周的净卖出规模是2024年11月以来最大的。这是否意味着散户层面的情绪也开始转向？如果私人客户从净买入转为持续净卖出，市场的边际买盘将进一步减弱。

这些悬而未决的问题，恰恰是这份报告最有价值的部分——它告诉我们市场正在发生什么，但没有告诉我们接下来会发生什么。这正是深度研究与碎片化信息的区别所在。

5. 给决策者的观察框架：关注“资金流向的边际变化”，而非“绝对水平”

\subsection{[R020] NOM：中国5月出口强劲，但市场真正低估的不是需求，而是价格信号的重估}
{\small\color{refgray}归类：中国贸易：价格驱动的结构性重定价与AI主导的出口新格局}

NOM：中国5月出口强劲，但市场真正低估的不是需求，而是价格信号的重估

中国5月贸易数据再次超出预期。以美元计，出口同比增长19.4\%，远高于市场一致预期的15.0\%和NOM自己的13.7\%预测；进口增长27.4\%，也高于预期的26.0\%。贸易顺差进一步扩大至1054亿美元。

表面上看，这是一份“需求强劲”的报告。但NOM这份研报的核心洞察恰恰相反：5月贸易数据的驱动力并非真实需求的扩张，而是一场由价格主导的结构性重定价。 芯片价格飙升、能源价格冲击、以及关税政策变动带来的基数效应，共同制造了一个“名义增长强劲、实际量价严重背离”的贸易图景。

理解这一点，比知道出口增长了19.4\%重要得多。因为名义数据的强劲，正在掩盖一个关键事实：中国贸易的定价权正在发生转移，而这种转移将深刻影响产业链利润分配、汇率定价逻辑以及资产配置的方向。

1. AI相关出口已连续两个月贡献中国总出口增长的近一半，但驱动力是价格而非数量

这是整份报告最值得关注的结构性变化。NOM估算，5月集成电路（IC）单独贡献了出口增长的5.9个百分点，自动数据处理设备（ADP）贡献了3.4个百分点，两者合计9.3个百分点，占中国5月总出口增长的约48\%。这已是连续第二个月AI相关出口贡献率接近50\%。

但真正的看点不在于比例高，而在于驱动方式。5月IC出口金额同比增长110.9\%，但出口数量仅增长2.1\%，这意味着价格贡献高达106.5个百分点。进口端同样如此：IC进口金额增长68.0\%，但数量转为负增长-1.0\%，价格贡献69.8个百分点。DRAM和NAND二季度合约价格上涨的传导效应，正在制造一场“名义繁荣、实际平淡”的半导体贸易。

这意味着什么？对于投资者而言，需要区分两类受益者：一类是真正受益于AI需求放量的产能持有者，另一类只是被动享受价格上涨带来的库存重估。当前市场可能过度定价了前者的持续性，而低估了后者对利润率的脉冲式影响。一旦芯片价格周期转向，名义增长数字将迅速回落。

2. 对美出口飙升至37.3\%背后，关税政策的“减法”比“加法”更重要

5月中国对美出口同比增长37.3\%，创2021年3月以来新高。NOM指出了三个驱动因素：低基数、关税净减少、以及美国AI需求。

其中关税因素最值得深究。2026年2月底，美国最高法院推翻了基于IEEPA的关税（包括10\%的互惠基线关税和10\%的芬太尼关税），随后特朗普政府依据Section 122实施了新的10\%临时关税。NOM估算，这一调整使美国对华有效关税率从IEEPA时期的较高水平降至3月的约27.7\%。5月中旬，美国财长贝森特也公开确认了中国近期面临的关税税率因最高法院裁决而降低。

这是一个被市场低估的变量。市场往往只关注“加关税”的负面冲击，却忽略了“减关税”带来的短期利好。而且，6月2日美国宣布的新一轮301关税（对中国12.5\%）将在7月24日Section 122关税到期后生效，这意味着6-7月可能出现明显的“抢出口”行为。

所以，5月对美出口的高增长，部分是一次性的政策窗口红利，而非竞争力的根本性改善。投资者需要警惕三季度数据中可能出现的“基数效应反转”和“抢出口透支”。

3. 能源进口的“量价背离”比想象中更极端，中国正在主动减少采购

进口端同样存在严重的量价背离。5月原油进口金额同比增长15.3\%，但进口数量同比暴跌29.0\%，较4月的-20.0\%进一步恶化。NOM指出，这反映了霍尔木兹海峡供应中断与北京在高油价下主动减少采购的双重影响。

但更值得关注的是，原油对进口增长的贡献仅为1.9个百分点，而IC的贡献高达10.8个百分点。这意味着当前中国进口的“高增长”本质上是一个半导体价格故事，而非能源故事。市场如果简单用“进口高增长=内需强劲”的逻辑推

[... middle omitted ...]

对比的是，电子产品出口全面飙升：IC增长111.0\%，ADP设备增长66.1\%，手机增长43.9\%。NOM特别指出，手机出口的增长同样由价格效应驱动，因为投入成本飙升。

这组对比揭示了一个趋势：中国出口结构正在被AI超级周期强行重塑。传统劳动密集型产品的竞争力在关税与成本压力下持续流失，而技术密集型产品虽然名义增长亮眼，但利润空间受制于上游芯片定价权。对于企业而言，核心问题不再是“要不要转型”，而是“转型后的利润归属权在哪里”——如果芯片价格持续高企，下游组装和制造环节的利润空间将被持续挤压。

5. 报告尚未完全回答的关键问题：价格驱动的增长何时见顶，利润分配如何变化

NOM的这份报告数据扎实、逻辑清晰，但它也留下了一些尚未完全展开的问题。

第一个问题是芯片价格的持续性。当前IC价格的飙升主要来自DRAM和NAND合约价格上涨的传导，但这一轮涨价周期能持续多久？NOM没有给出明确判断。如果AI需求增速放缓或供给端产能释放，价格下行将对名义贸易数据产生剧烈冲击。

\subsection{[R021] Bernstein：数据中心真正的拐点不在建设量，而在供电模式的结构性重置}
{\small\color{refgray}归类：AI/算力：资本开支上行风险与供应链结构性重置}

Bernstein：数据中心真正的拐点不在建设量，而在供电模式的结构性重置

市场对AI基础设施的讨论，大多集中在GPU算力需求、资本开支节奏、以及“泡沫何时破裂”的宏观叙事上。这些讨论并非没有价值，但它们往往忽略了一个正在发生的、更为根本的结构性变化：数据中心从“电网依赖型”向“自备发电型”的加速迁移。Bernstein最新发布的北美数据中心追踪报告（2026年5月）提供了一个关键信号：这一迁移的速度和规模，正在重塑整个产业链的价值分配，而市场对此的定价可能仍不充分。

这份报告的核心价值，不在于它更新了324GW的总管道容量，也不在于它记录了4.4GW的新增在建工程。真正值得关注的，是“自备发电”（Behind-the-meter, BTM）管道容量在一个月内增长了14GW，同比年初至今暴增42GW，达到129GW。更关键的是，目前在建设的容量中已有25\%采用BTM模式，而管道中的项目这一比例已高达40\%。这意味着，未来几年内，将有近一半的新增数据中心不再完全依赖公共电网供电。

为什么这很重要？因为供电模式的选择，直接决定了哪些供应商能够从中受益。当数据中心自备发电时，其对电气设备的需求结构会发生根本性变化。Bernstein估算，自备发电项目在电气分配设备上的支出比纯电网接入项目高出约30\%。这才是这份报告最值得仔细拆解的核心洞察：市场真正低估的不是数据中心的需求总量，而是供电模式重置带来的供应链价值再分配。

1. 管道总量增长10\%并非重点，真正值得追问的是“谁在推动增长”

Bernstein的报告显示，5月份北美数据中心项目管道总量增长了29GW，达到324GW，环比增长10\%。这个数字本身并不令人意外——市场已经习惯了AI基础设施的指数级增长叙事。但报告揭示的增长驱动力分布，却与传统认知存在显著差异。

当月新增的29GW中，开发商（Developers）贡献了22GW，占比75\%；加密货币矿工（Crypto Miners）贡献了6GW，占比21\%。而超大规模云服务商（Hyperscalers）的管道容量反而出现了小幅下降。这意味着，当前数据中心建设热潮的“增量部分”，并非由亚马逊、微软、谷歌这些市场最熟悉的巨头直接驱动，而是由一批新兴的第三方开发商和转型中的矿工在推动。

O'Leary Ventures单月新增15GW管道容量，凯雷集团新增3GW，这些名字对大多数市场参与者而言仍相对陌生。加密货币矿工方面，Core Scientific、Riot、TeraWulf分别新增2.3GW、1.4GW和1.2GW。这些矿工正在从单纯的比特币挖矿转向AI基础设施服务商，其现有的大规模电力基础设施和运营能力，使其成为“时间就是算力”竞争中的关键参与者。

这一变化意味着什么？它暗示数据中心建设的决策权正在发生微妙转移。超大规模云服务商仍然是最重要的最终租户，但建设节奏和选址决策的主动权，正在越来越多地向拥有电力资源和土地储备的第三方实体倾斜。对于投资者而言，单纯跟踪云服务商的资本开支指引，已经不足以判断实际建设进度。

2. 在建工程稳步爬坡，但真正值得关注的信号来自“搁浅容量”的稳定

5月份的数据中心在建容量增加了4.4GW，达到63GW。其中超大规模云服务商贡献了1.4GW，开发商贡献了1.3GW，新云服务商（Neoclouds）贡献了726MW。Meta和亚马逊分别贡献了638MW和536MW，占超大规模云服务商新增量的约90\%。

这些数字本身是积极的，但报告中最耐人寻味的数据点，是“搁浅容量”（Stranded Capacity）的变化。所谓搁浅容量，是指因当地社区反对（NIMBY）、监管审批受阻或其他原因而无法推进的项目。报告显示，搁浅容量在5月份增加了1.9GW，达到34GW，占管道总

[... middle omitted ...]

降，将意味着供给侧的制约正在缓解，建设加速的确定性将进一步增强。

3. 自备发电模式的加速推进，正在从根本上改变产业链的价值分配

这是Bernstein这份报告最核心的洞察，也是市场讨论中最容易被忽视的部分。BTM管道容量在5月份增长了14GW，环比增长12\%，年初至今累计增长42GW，达到129GW。更值得关注的是结构性的变化：

目前只有4\%的已安装数据中心采用自备发电模式。 但在建容量中，这一比例已经跃升至25\%。 在管道项目中，这一比例进一步攀升至40\%。 5月份新增的管道容量中，近50\%采用了BTM模式。

这些数字的递进关系非常清晰：自备发电不是一种边缘尝试，而是正在成为新建数据中心的主流选择。开发商是这一趋势最积极的推动者，5月份新增了18GW的BTM管道容量，年初至今累计贡献了35GW，占BTM总增量的80\%以上。从运营商类型来看，新云服务商对BTM的依赖度最高（管道中64\%采用BTM），其次是开发商（50\%），超大规模云服务商相对最低（36\%）。

\subsection{[R022] 摩根斯坦利：中国金融体系真正的转折点不是刺激，而是“高质量”的内生循环}
{\small\color{refgray}归类：宏观与利率：亚洲通胀反弹与央行紧缩路径}

摩根斯坦利：中国金融体系真正的转折点不是刺激，而是“高质量”的内生循环

市场对中国金融板块的讨论，长期陷入一个二元对立框架：要么依赖大规模财政刺激拉动信贷，要么承受经济减速带来的资产质量恶化。但摩根斯坦利在近期发布的一份系列研报中，提出了一个更具穿透力的判断——中国金融体系正在进入一个“更可持续的高质量发展”正循环，而这个循环的驱动力并非来自政策强心针，而是来自产业结构自身调整带来的利润修复、信贷理性化以及居民资产配置的结构性迁移。

这份报告的核心主张值得每一位关注中国资产定价的决策者认真对待：市场真正低估的，不是短期需求刺激的力度，而是供给侧出清与金融体系内生质量改善叠加后，银行板块盈利能力和估值体系可能发生的系统性重定价。

1. 出口与工业利润的“超预期正循环”正在取代旧有的刺激依赖

报告开篇即点明一个关键背景：尽管宏观数据在近几个月出现一些分化与波动，但一个没有大规模刺激的“高质量发展”路径，正在为金融体系提供更可持续的支撑。这个判断并非空谈，而是建立在三个相互印证的数据链条之上。

第一，出口增长在2026年一季度表现强劲，YTD出口增速在4月仍维持在14.5\%的同比高位。更重要的是，摩根斯坦利指出，这次出口增长并非简单的价格竞争或全球补库，而是由“科技与产业升级”驱动。这意味着出口增长的可持续性和含金量，与过去几轮周期有本质区别。

第二，工业利润的修复正在加速。2026年4月，制造业利润YTD同比增速进一步加快至20.4\%。与此同时，PPI在4月环比上涨1.7\%，同比上涨2.8\%，正式回到正值区间。PPI回正与利润增长同步发生，说明企业端不仅是在“量”上恢复，更是在“价”上获得了改善，这是利润修复质量提升的关键信号。

第三，最值得关注的信号来自制造业产能扩张的理性化。报告数据显示，制造业固定资产投资增速从3月的4.1\%降至4月的1.2\%，而制造业利润增速却在加速。这意味着企业正在用更少的资本投入换取更高的利润回报，产能过剩的风险正在系统性降低。

这三个信号合在一起，构成了一个清晰的逻辑链条：产业升级驱动的出口为经济提供了外需支撑，工业利润修复改善了企业资产负债表，而资本开支的理性化则降低了未来产能过剩和坏账积累的风险。这个正循环一旦确立，银行体系面临的信用风险将显著下降，而无需依赖额外的政策刺激来“兜底”。

如果说工业端的数据更多反映供给侧，那么支付数据和居民金融资产的变化，则从需求侧为金融体系的稳健性提供了微观证据。

报告披露了一组关键数据：2026年一季度，全系统支付总额同比增长29\%，而2025年四季度这个数字仅为5\%。这是一个巨大的跃升。其中，银行卡消费增速在1Q26同比转正至3.2\%，而2025年全年仍处于负增长区间。网联支付增速虽然从2025年四季度的16.5\%小幅回落至7.2\%，但依然保持在正增长区间。

支付数据的全面回暖，至少传递了两个重要信息。其一，居民消费信心正在实质性修复，而非停留在政策口号层面。其二，消费支付从第三方支付平台向银行卡回流的趋势在延续，这对银行的中间业务收入（特别是银行卡手续费）是直接利好。

更值得关注的是居民金融资产配置的结构性变化。报告提到，新的居民储蓄正在持续向投资迁移。这个趋势如果延续，意味着银行的负债成本有望进一步下降（因为定期存款占比下降），同时财富管理业务将获得持续的资金流入。这对于银行的净息差和手续费收入都是双重利好。

综合来看，支付数据回暖和居民资产迁移，共同构成了银行资产负债表两端同时改善的微观基础：资产端，消费信贷质量受益于居民收入预期改善；负债端，资金成本受益于储蓄向投资的迁移。这个微观基础，比任何宏观政策信号都更可靠。

对于长期关注中国金融的投资者而言，信贷增速放缓往往被视为经济疲弱的信号。但摩根斯坦利在这份报告

[... middle omitted ...]

体经济对信贷的依赖正在降低，而财政政策在发挥托底作用的同时，并未引发信贷泡沫。

对于银行股估值而言，信贷增速放缓最直接的正面含义是：银行不再需要通过牺牲定价来争夺市场份额。当银行能够以更合理的价格、更审慎的标准投放贷款时，净息差的底部将更早到来，资产质量的确定性也将更高。报告明确提出，摩根斯坦利预计中国银行的净息差将“比预期更早地反弹”，这正是基于信贷供给端理性化这一核心假设。

4. 报告尚未完全回答的关键问题：利润修复的可持续性与行业分化

尽管报告勾勒了一个令人信服的正循环图景，但作为严谨的读者，我们必须承认其中仍存在几个尚未完全回答的关键问题。

第一个问题是制造业利润修复的可持续性。报告显示，4月制造业利润YTD增速高达20.4\%，但这是建立在2025年同期较低基数之上的。随着基数效应消退，利润增速能否维持在高位，取决于出口需求是否持续强劲以及PPI能否稳定在正值区间。而出口需求本身又面临全球贸易环境的不确定性，报告对此着墨不多，这需要投资者自行跟踪判断。

\subsection{[R023] Bernstein：亚马逊供应链的真实威胁，被市场高估了}
{\small\color{refgray}归类：新兴科技：AI音乐变现与量子计算速度精度取舍}

过去几周，市场对亚马逊推出“Supply Chain by Amazon”的讨论持续升温。投资者最关心的问题是：当亚马逊向第三方卖家开放其强大的履约网络，UPS和FedEx的包裹业务会不会遭遇结构性冲击？

直觉上，答案似乎是肯定的。亚马逊的履约能力有目共睹——Prime两日达甚至当日达已经成为行业标准，其物流基础设施的投入和效率也让传统承运商相形见绌。如果更多商家能够接入这套体系，UPS和FedEx的份额流失似乎是必然。

但Bernstein这份研报给出了一个反直觉的核心判断：亚马逊的前置库存模型在供应链结构上存在天然局限，真正适合采用该模式的零售市场上限约为20\%，而考虑到数据隐私和服务容错等现实约束，实际可渗透的市场可能只有低个位数百分比。

这不是一个关于“亚马逊能不能做”的问题，而是一个关于“哪些供应链结构天然不适合前置库存”的问题。Bernstein的分析框架揭示了决定这一问题的三个核心变量：持有成本、货位成本和运输节省。当这三者的数学关系摆到台面上，亚马逊供应链的真实威胁边界就清晰了。

要理解Bernstein的核心论点，首先要接受一个基本前提：供应链不是“一刀切”的解决方案。不同品类的产品属性、客户需求和成本结构，决定了完全不同的网络拓扑设计。

Bernstein在报告中明确指出，供应链设计需要权衡三大类成本：运输成本、库存处理成本（包括仓储和地产），以及软性成本（如融资成本、过时折价和降价损失）。高周转、低价值密度的商品，适合采用密集的前置网络以最小化最后一公里成本和周期时间；而高价值、慢周转的商品，则更适合集中化库存以保持合并效应并降低持有成本。

亚马逊的模型本质上是一种“前置库存”模型——将库存分散部署在离消费者更近的位置，以此缩短交付时间并降低运输成本，代价是更高的库存投资和设施投入。这种模型在“便宜运输、快速销售、低价值密度”的商品类别中具有优势，但在“昂贵运输、慢速销售、高价值密度、品类异质性高”的商品中，结构性不经济。

这一判断的市场含义是明确的：并非所有零售品类都能从亚马逊的供应链服务中受益，大多数品类甚至可能因为前置库存而承担更高的总成本。 对于UPS和FedEx而言，它们的高价值客户群体——比如工业品、医疗设备、高端消费品等——天然不适合亚马逊的模型，这些客户的长距离运输需求不会消失。

Bernstein构建了一个透明的数学模型来量化前置库存与传统集中库存的成本差异。模型的核心公式很简单：Delta（本地-集中）= 持有惩罚 + 货位惩罚 - 运输节省。当Delta为负数时，前置库存模型才更优。

这个公式揭示了一个反直觉的结论：决定一个品类是否适合前置库存的关键变量，不是运输成本，而是库存持有成本和货位成本。

持有惩罚是最具决定性的因素。当库存从集中节点分散到20个本地节点时，安全库存量会按照平方根法则放大。对于单价超过200美元的商品，这种库存碎片化带来的持有成本增长，几乎不可能被运输节省所抵消。

货位惩罚是一个结构性变量。它随SKU广度增长，但与商品单价无关。对于品类繁多的零售商，将数千个SKU分散部署到每个本地节点，意味着每个节点都需要承担固定的货位成本。即使机器人技术能够降低单位货位需求，也无法从根本上改变这一成本结构。

运输节省是唯一有利于前置库存的变量，但其规模有限。Bernstein的模型假设每单位运输可节省约9-10美元，即使将节省幅度放大到1.2倍，对于单价超过400美元的商品，前置库存仍然不经济。

Bernstein的敏感性分析清晰地展示了这一点：在单价50美元、运输节省乘数0.4倍的情况下，Delta为正28.9万美元，集中模型胜出；只有在单价50美元以下且运输节省乘数超过0.8倍时，前置库存才可能占优。这意味着亚马逊供应链服务的主

[... middle omitted ...]

家需要与亚马逊共享详细的销售数据、库存信息和客户行为数据。对于许多品牌商和零售商而言，这是不可接受的。它们不愿意将核心商业数据暴露给一个既是平台又是竞争对手的实体。Bernstein在报告中提到，对数据保护的容忍度会显著限制实际采用率。

第二，服务容错。亚马逊的履约网络虽然高效，但并非没有缺陷。对于高价值商品或对交付时间有严格要求的品类，商家可能不愿意将全部履约控制权交给第三方。一旦出现服务失败，品牌声誉的损失可能远超运输成本节省带来的收益。

这两个约束条件叠加，意味着即使理论上适合前置库存的品类中，也只有一小部分商家会真正选择采用亚马逊的供应链服务。Bernstein的结论是：对UPS和FedEx的包裹量而言，来自亚马逊供应链服务的风险是可控的，远没有市场担忧的那么严重。

4. 报告尚未完全回答的关键问题：亚马逊的定价策略和长期竞争动态

虽然Bernstein的分析框架具有很强的说服力，但报告对几个关键问题没有给出充分答案，这恰恰是投资者需要持续跟踪的方向。

\subsection{[R024] Citi：中国生物科技面临的地缘政治风险，市场可能高估了法案通过概率}
{\small\color{refgray}归类：医疗健康：药价谈判结构性变化与设备招标拐点}

Citi：中国生物科技面临的地缘政治风险，市场可能高估了法案通过概率

当前市场对中美生物科技脱钩的担忧，可能正在定价一个低概率事件。Citi在最新一份专家电话会纪要中释放了一个清晰但容易被忽视的信号：旨在限制美国对华生物科技投资的BINSA法案，在本届国会会期内成为法律的可能性非常低。这个判断不是基于乐观情绪，而是基于对美国立法流程、行政机构态度和产业游说力量的结构性分析。

这份报告的价值不在于它给出了一个“风险解除”的结论，而在于它拆解了为什么市场习惯性高估这类法案的落地概率。对于持有中国CXO、生物科技资产的投资者，以及正在评估供应链风险的产业决策者，理解这个判断背后的逻辑，比直接相信结论更重要。

1. 法案通过概率低于市场隐含预期，核心障碍在立法流程而非政治意愿

Citi专家给出的BINSA法案通过概率区间为30-40\%以下。这个数字本身已经低于许多市场参与者的隐含预期。但更值得关注的是概率的结构：三种可能的立法路径各有不同的障碍，而每一条路径都面临实质性的制度阻力。

路径一，作为FY2027 NDAA修正案附加上市，概率30-35\%。关键障碍在于众议院军事委员会主席明确反对将不相关条款附加到NDAA上。FY2027 NDAA的众议院版本已经在委员会层面清除了BINSA相关条款。路径二，由财政部采取行政行动，概率35-40\%。但财政部已明确表示，现行法律下它没有监管生物科技对外投资的法定权限。路径三，作为独立法案通过，概率低于10\%，因为目前没有参议院配套法案。

这些障碍不是“有可能被克服”的，而是制度性的。NDAA历来有“不相关条款”的惯例约束，财政部的法定权限边界清晰，而独立法案需要完整的立法周期——专家估计至少1至1.5年。在一个国会会期内完成所有这些步骤，难度远大于市场直觉。

这意味着，市场对BINSA的定价，可能包含了一个“快速通过”的尾部风险溢价，而这个风险的实质概率被高估了。

2. 生物科技被排除在COINS法案之外并非偶然，而是产业特征决定的制度选择

为什么生物科技没有被纳入已有的COINS法案？这个问题比BINSA本身更值得深思。Citi专家的解释揭示了一个被市场忽略的结构性逻辑：半导体、人工智能和量子计算被优先纳入，是因为它们“明显具有军事赋能属性”且已被高度安全化。生物科技则不同——它拥有庞大的民用和商业应用场景，对普通立法者而言，不那么容易被归类为“硬技术”。

这个区别决定了监管路径的根本不同。财政部明确表示没有法定权限监管生物科技对外投资，不是行政不作为，而是法律框架的边界使然。要改变这个边界，需要新的立法授权，而新的立法授权需要面对比修改已有法案更高的立法门槛。

更深层的含义是：生物科技领域的“国家安全化”进程，天然慢于半导体和AI。这不是因为生物科技不重要，而是因为它的产业特征——高度依赖全球临床合作、跨国供应链和许可协议——使得任何仓促的监管都可能造成超出预期的商业损害。这种损害的不确定性，反过来又成为立法者的顾虑。

3. 制药行业的游说力量是法案推进的结构性阻力，且正在积极运作

Citi专家明确指出，美国大型制药公司和生物科技公司正在积极游说，目标是完全阻止BINSA或争取大幅修改。这个游说力量不是边缘性的，而是产业核心利益驱动的。

大型制药公司依赖中国的临床试验、制造和合作伙伴关系。对外投资限制可能直接打乱研发管线、延迟药物审批、提高成本。如果生物科技被纳入一个粗放的对外投资监管框架，常规的许可和合作交易都将面临风险。商业联盟还指出，生物科技已经受到BIOSECURE式采购规则和供应链审查的压力，再加对外投资资本管制将增加额外的不确定性。

这些不是抽象的理由，而是具体的商业损失测算。制药行业在美国政治中的影响力——尤其是在马萨诸塞州等生物科

[... middle omitted ...]

康德纳入，但Citi认为基本面影响极小，因为只有联邦资助项目被禁止与药明康德合作。2026年2月撤回的版本已经包含药明康德，市场对此已有预期。

但这里有一个值得深入追问的问题：1260H清单的纳入虽然不构成制裁，但提高了美国实体与被列入公司交易的合规负担。对于药明康德这样的CXO企业，其客户结构中美国大型制药公司的占比是关键变量。如果客户因为合规成本上升而主动减少合作，即使法律上不禁止，商业影响也可能逐步显现。

Citi专家指出，被列入公司可以先向国防部申请行政复议，如果复议不成功，可以在联邦法院提起诉讼，质疑清单缺乏合理依据。药明康德已表示反对被纳入，专家预期其将迅速启动上诉程序。但法律程序的周期和结果都不确定。这里需要持续跟踪。

5. 报告尚未完全回答的关键问题：如果法案通过概率被高估，市场何时开始重新定价

Citi这份报告的核心贡献是提供了一个概率框架，但概率本身不是交易信号。真正的问题是：如果市场正在高估BINSA的通过概率，这个错误定价何时会被纠正？

\subsection{[R025] GS：欧洲半导体设备市场的下一个增长动力，不是周期，是结构}
{\small\color{refgray}归类：半导体：欧洲设备需求周期尚未见顶与存储供给纪律重塑}

这份GS最新发布的欧洲科技硬件研报，在标题中使用了“更新估值与目标价”这样克制的措辞。但当我们穿透数字调整的表层，会发现这远不止是一次季度性的财务模型微调。它实际上在传递一个更为根本的信号：欧洲半导体设备与AI基础设施领域，正在经历一轮由多重结构性力量共同驱动的价值重估。

GS分析师Alexander Duval及其团队，同时上调了ASML、ASMI、BESI和Nebius四家公司的目标价，并同步提升了远期盈利预测。这四家公司分别覆盖光刻、沉积、封装和AI算力基础设施——几乎构成了从芯片制造到AI部署的完整价值链。而驱动这轮集体上调的，并非单一事件，而是三个彼此独立、却同时指向同一方向的信号。

这份报告最值得关注的判断是：市场此前对欧洲AI使能者的定价，低估了其长期增长空间的持续扩展能力。当前的上修，只是这一过程的开始。

GS在报告中列举了近期三个关键数据点，分别来自存储、AI基础设施和中国市场。单独看，每个信号都有其自身的产业逻辑；放在一起，它们共同回答了一个关键问题：这轮半导体设备投资的强度，是否已经接近尾声？

第一个信号来自SK海力士。据媒体报道，这家全球领先的HBM供应商计划在未来五年内将晶圆产能翻倍。GS明确指出，这一计划进一步增强了市场对“多年存储器投资周期”的信心。HBM（高带宽存储器）是AI芯片的关键配套，其产能扩张意味着对前道设备的需求将持续释放。这不是一个短期补库存行为，而是基于AI推理和训练需求的结构性扩产。

第二个信号来自博通。博通重申其预期：到2027财年，AI半导体收入将超过1000亿美元，对应10GW的数据中心容量建设。更重要的是，博通预计在2028财年，即使基数已经很高，仍能实现强劲增长。这直接挑战了“AI基建投资将在2026年见顶”的叙事。

第三个信号来自设备行业的中国业务。LAM Research和KLA近期的评论显示，中国WFE（晶圆厂设备）支出在2026年将保持“大致持平或微幅增长”。考虑到此前市场普遍预期中国设备采购将在2025年后明显回落，这一信号意味着全球设备需求的下行风险被显著降低。

三个信号合并，GS得出的结论是：全球设备需求环境比此前预期的更为健康。这不是一个“见顶回落”的故事，而是一个“高位持续”的叙事。

2. “主权AI”与GPU定价权：AI基础设施的定价逻辑正在被重写

如果说半导体设备端的信号更多是“量的扩张”，那么AI基础设施端的信号则是“价的强化”。

GS引用了Nvidia的数据：在2027财年第一季度，其“主权AI”收入同比增长超过80\%。主权AI指的是各国政府主导的AI基础设施建设，这一增速远超市场此前对政府支出的预期。这意味着，除了大型云厂商（Hyperscalers）之外，一个由国家级需求驱动的新增长极正在形成。对于Nebius这类AI云服务商而言，主权AI意味着客户基础的多元化和合同期限的延长。

更值得关注的是GPU定价。Nvidia指出，H100和A100的定价年初至今分别上涨了20\%和15\%。在AI芯片供应持续紧张、需求仍在加速的背景下，定价权牢牢掌握在供给端。这对整个AI基础设施生态的影响是深远的：更高的GPU定价意味着更高的算力租赁价格，也意味着Nebius这类AI Neocloud公司的单位收入能力在提升。

GS正是基于这一逻辑，上调了Nebius的远期收入预测和目标价。Nebius在6月8日宣布将在英国投资约17亿英镑，部署三个由Nvidia GPU驱动的新设施，总容量预计到2027年达到65MW。这是供给端积极响应需求信号的具体案例。

这一小节的核心含义是：AI基础设施的定价逻辑正在从“规模决定价格”转向“稀缺性决定价格”。对于投资者而言，这意味着在评估AI基础设施公司时，不能再用传统的云服务估值框架

[... middle omitted ...]

从895欧元上调至955欧元，估值倍数从24倍提升至25倍。其逻辑核心是AI驱动的先进逻辑芯片需求，ASMI在这一领域有显著敞口。这是一个典型的“AI应用拉动设备需求”的故事。

BESI的目标价从286欧元上调至315欧元，估值倍数从31倍提升至33倍。驱动因素来自两个方面：主流封装业务的强劲需求，以及混合键合（Hybrid Bonding）业务在先进逻辑芯片应用中的顺风。BESI的估值提升，反映的是封装环节在AI价值链中地位的上升。

Nebius的目标价从234美元上调至267美元，估值倍数从8倍提升至9倍。其逻辑最为直接：强劲的AI基础设施需求、稳定的电力扩容进展、以及优于预期的定价环境。

四家公司，四种不同的估值驱动因素。这提示我们：欧洲AI使能者并非同质化板块，每一家公司的价值重估都有其独特的产业叙事。投资者需要区分“行业贝塔”和“公司阿尔法”。

尽管GS对中国WFE支出的判断比市场更乐观，但报告并未完全展开讨论一个关键问题：中国业务的高水平能否持续？

\subsection{[R026] 摩根斯坦利：CPO市场真正的分歧不在技术，而在量产节奏的预期差}
{\small\color{refgray}归类：AI/算力：资本开支上行风险与供应链结构性重置}

摩根斯坦利：CPO市场真正的分歧不在技术，而在量产节奏的预期差

市场对CPO（共封装光学）的担忧正在加剧。过去几周，投资者反馈的核心焦虑是：CPO导入是否在推迟？相关标的的抛售是否意味着长期逻辑动摇了？

摩根斯坦利在最新发布的CPO专题报告中给出了一个反直觉的判断：长期故事没有改变，但市场短期预期确实需要重置。这份报告的价值不在于告诉你CPO会来，而在于精确地告诉你它将以什么样的节奏、在什么样的约束条件下到来——以及当前股价中隐含的假设与这个节奏之间有多大的裂口。

报告的核心结论可以浓缩为一句话：2027年光学引擎出货量预计仅为600-700万颗，远低于市场期待的 万颗。这不是技术失败，而是量产爬坡的物理约束。而真正的爆发，将从2028年开始。

1. 市场对2027年的出货量预期存在3-5倍的过度乐观，这是短期情绪修正的核心驱动力

摩根斯坦利对2027年光学引擎（Optical Engine, OE）出货量的基准预测是600-700万颗，涵盖scale-up和scale-out两种方案。而投资者普遍期待的出货量落在 万颗区间。

这个差距不是技术路线之争，而是量产节奏的根本性误判。市场往往倾向于线性外推——一旦看到台积电宣布将PIC产能扩张至10kwpm，就默认出货量会同步放大。但摩根斯坦利的分析指出，产能扩张只是必要条件，远非充分条件。

这里的关键变量是良率。报告明确指出，SolC（Silicon on Carrier）良率目前仅为50-60\%，下游组装良率更是低至20-50\%。这两个数字才是决定最终出货量的真正瓶颈。产能可以快速扩张，但良率的提升需要时间、经验曲线和工艺迭代。

这意味着什么？2027年的CPO相关收入可能远低于当前股价所隐含的水平。对于持有CPO概念股的投资者来说，未来几个季度的财报季很可能成为预期修正的触发点。这不是基本面恶化，而是市场定价从“完美场景”向“现实场景”回归的过程。

摩根斯坦利提出的一个关键时间框架是：2026到2028年，收发器、光学方案（CPO/NPO）和铜互联将共存。主流解决方案仍然停留在1.6T和3.2T，CPO还需要时间爬坡。

这个判断的重要性在于，它打破了市场对“CPO一旦导入就全面替代”的二元叙事。实际上，数据中心内部的互联架构正在变得更加复杂，而非更加简单。不同的传输距离、带宽密度和功耗要求，将催生不同技术方案的分层应用。

对于产业链上的公司而言，这意味着收入结构在过渡期内不会是单一驱动。那些同时布局多种技术路线的供应商，可能比押注单一路线的公司更具韧性。而那些在CPO组件领域有独特优势的公司，其收入贡献在2028年之前可能仍然只是整体营收的一小部分。

报告特别提到，FOCI的CPO相关收入要到2026年三季度末才开始爬坡，在此之前收入将保持温和。这实际上给出了一个可验证的时间节点——投资者可以用三季度财报来检验CPO业务的真实进展。

3. 良率不是短期噪音，而是决定CPO供应链价值分配的结构性变量

SolC良率50-60\%、下游组装良率20-50\%——这两个数字值得反复咀嚼。它们不仅仅是技术参数，更是供应链利润分配的前置条件。

在半导体制造的历史上，低良率阶段往往是设备商和材料商享受超额利润的时期。因为下游厂商愿意为提升良率支付溢价。但在CPO的语境下，情况可能有所不同。CPO的封装和测试环节涉及大量新工艺，良率的提升依赖于设备精度、工艺稳定性和操作经验的积累。

这引出一个关键问题：谁将从良率提升的过程中受益最大？摩根斯坦利维持对台积电、日月光、FOCI、AllRing、MPI、鸿海精密和WinWay等核心CPO供应商的增持评级，这暗示该机构认为这些公司有能力在良率爬坡过程中建立竞争优势。但对于那些仅提供标准化组件的公司，低良率

[... middle omitted ...]

大价值

尽管摩根斯坦利的分析框架非常清晰，但有两个问题报告尚未完全展开。

第一，规模效应的时间点。报告指出CPO从2028年开始爆发，但没有量化“爆发”的具体幅度。如果2027年出货量仅为600-700万颗，2028年能否跃升至 万颗？这取决于良率提升的速度、下游客户（云厂商）的采购节奏以及标准化的推进程度。目前来看，这些因素都还存在较大的不确定性。

第二，价值分配的结构。CPO产业链涉及设计、晶圆制造、封装、测试、光引擎组装、光纤连接等多个环节。每一环节的竞争格局和技术壁垒差异很大。报告列出了多家增持标的，但并未给出一个清晰的“谁在什么时间点受益最多”的排序。对于机构投资者而言，这可能是决定超额收益来源的核心问题。

这两个未解问题恰恰是当前市场分歧的根源。乐观者看到的是2028年后的爆发潜力，悲观者看到的是2027年出货量远不及预期。而真正的投资机会，可能存在于两者之间的预期差中。

5. 投资者的观察框架应从“CPO何时到来”转向“哪些约束条件正在被打破”

\subsection{[R027] 摩根斯坦利：中国保险业真正的分化不在规模，而在产品策略与渠道控制力}
{\small\color{refgray}归类：未被主题摘要引用}

摩根斯坦利：中国保险业真正的分化不在规模，而在产品策略与渠道控制力

中国保险业正在经历一场由监管驱动、龙头主导的深度结构性转型。摩根斯坦利在最新发布的2026年行业观察报告中指出，过去两年间，上市保险公司已基本完成从传统险向分红险的产品切换，但这一轮转型的真正含义并非简单的产品替代，而是行业竞争格局的重新洗牌——大型保险公司与中小型保险公司正在走向截然不同的发展路径。

这份报告最值得关注的核心判断是：市场对保险板块的定价可能低估了龙头公司在负债端转型与资产端优化中积累的结构性优势，同时也高估了中小保险公司在低利率环境下维持增长的能力。行业整体风险可控，但分化正在加速，而2026年一季度开始执行的更严格的偿付能力报告披露规则，将第一次让外界清晰地看到这种分化。

摩根斯坦利的数据显示，到2026年一季度，上市保险公司的新业务中分红险占比已超过80\%，部分公司甚至接近90\%。这一数字在2025年上半年还不到50\%。转型速度之快，远超市场预期。

但报告真正有价值的洞察在于：大型保险公司与中小型保险公司对这一趋势的态度正在发生根本性分歧。大型保险公司将继续深耕分红险，依托代理人渠道的惯性维持产品延续性；而许多中小保险公司自2026年二季度起已经开始重新评估策略，部分公司甚至重新加大了传统险的销售力度。

第一，实际成本更高。虽然分红险的保证负债成本低于传统险，但计入非保证部分后的实际负债成本更高，降低分红实现率又可能削弱产品竞争力。

第二，利润分享机制侵蚀股东收益。分红险的客户利润分享特性意味着，当保险公司获得超额投资收益时，更多收益将流向客户而非股东，这实际上会减慢行业去风险化的进程。

第三，资本占用可能更高。分红险在更积极的资产配置策略下，资本密集度可能高于传统险，这对偿付能力本就紧张的中小公司构成额外压力。

第四，现金流管理更为复杂。分红险的资金管理涉及分红特别储备和退保风险，对于管理能力有限的中小保险公司而言，这种复杂性可能超出其运营能力边界。

这一判断的关键含义在于：市场不应简单地将所有保险公司视为“受益于分红险转型”的标的。产品策略的分化意味着，未来几个季度，大型保险公司与中小型保险公司在盈利能力、资本效率和增长可持续性上的差距将显著拉大。

2. 银保渠道的繁荣掩盖了一个结构性弱点：渠道依赖度决定了利润质量

银保渠道正在经历一轮强劲反弹。2025年新业务保费增长超过15\%，这一势头延续至2026年一季度，FYP同比增长约17\%，RP FYP同比增长约19\%。支撑这一增长的因素包括渠道利润率提升以及2023年大规模定期存款到期带来的再配置需求。

但摩根斯坦利的分析揭示了一个容易被忽视的结构性问题：银保渠道的依赖度与保险公司规模之间存在明显的负相关关系。

数据显示，资产规模在100亿至1万亿元之间的保险公司，其总保费中来自银保渠道的比例高达50\%至100\%。而资产规模超过1万亿元的大型保险公司，包括大多数上市险企，银保渠道的占比则显著更低，其代理人渠道贡献了总保费的60\%至80\%。

这意味着什么？大型保险公司拥有更强的渠道控制力和更可持续的价值创造能力。代理人渠道虽然增长较慢，但利润率更高、客户粘性更强、产品销售的复杂性管理能力更强。而中小保险公司高度依赖银保渠道，本质上是在用渠道控制权的让渡换取规模增长，这种增长的质量和可持续性值得怀疑。

对于投资者而言，这一观察框架意味着：评估一家保险公司的价值，不应只看保费增速，更要看渠道结构的质量。银保渠道占比过高的公司，在监管收紧或银行合作条件变化时，可能面临更大的增长波动和利润率压力。

3. 投资策略正在出现账户层面的分化，大型保险公司在另类投资中保持优势

在低利率和波动的市场环境下，保险公司的投资策略正在从“一刀切”转向账户层面的精细

[... middle omitted ...]

详细展开另类投资的具体类别和收益表现，但这一判断隐含了一个重要推论：在利率持续走低的背景下，另类资产的配置能力将成为保险公司投资收益率差异化的关键变量。

这里需要指出的是，摩根斯坦利并未提供另类投资的具体收益率数据或与公开市场投资的直接对比，因此这一判断的量化验证仍有待后续报告的补充。但基于行业逻辑，这一推论是合理的。

4. 偿付能力约束是硬天花板，但大型保险公司正在将其转化为竞争壁垒

偿付能力压力是贯穿整份报告的一条暗线。摩根斯坦利明确表示，尽管监管有所放松，但偿付能力压力仍将继续制约部分保险公司的资产配置和负债增长，且2026年偿付能力规则大幅修订的可能性有限。

这一判断的行业含义是深远的。偿付能力约束不仅限制了中小保险公司的增长空间，更重要的是，它正在成为大型保险公司维持竞争优势的制度性壁垒。当中小保险公司因资本约束而无法灵活调整资产配置或承接更多新业务时，大型保险公司凭借更充裕的资本缓冲和更优化的偿付能力管理，可以在同样的监管框架下获得更大的业务弹性。

\subsection{[R028] BARC：市场低估了中国出口的结构性韧性，AI与绿色技术正重塑贸易格局}
{\small\color{refgray}归类：中国贸易：价格驱动的结构性重定价与AI主导的出口新格局}

BARC：市场低估了中国出口的结构性韧性，AI与绿色技术正重塑贸易格局

这份由BARC在5月发布的中国贸易数据研报，其核心判断并非简单的“出口超预期”。真正值得关注的信号是：中国出口的驱动力正在从传统的成本优势与全球周期顺风，转向由AI资本开支与能源转型共同支撑的结构性需求。5月出口同比增长19.4\%，远超市场预期的15\%，这背后不是一次性的基数效应，而是全球制造业活动持续处于四年高位（PMI 52.6）与高附加值产品需求爆发共同作用的结果。

报告揭示了一个重要的产业现实：市场此前担心中东能源冲击会抑制外部需求，但事实证明，AI相关产品与绿色技术出口恰恰成为了对冲这一风险的关键力量。半导体出口同比暴增111\%，自动数据处理设备增长66\%，这些数字背后是持续的全球AI资本开支周期。中国作为AI制造组件的关键供应商，正在享受这一轮技术投资红利。与此同时，能源冲击反而加速了全球绿色转型，中国在电动车、锂电池、风电机组和太阳能电池等领域保持低成本高质量的供给优势，这些品类的年初至今出口均维持两位数增长。

这份报告的价值不仅在于数据本身，更在于它提供了一个观察中国出口新逻辑的分析框架。传统的贸易分析聚焦于目的地市场的景气度，而BARC的分析指向了一个更深刻的判断：中国出口的结构正在发生根本性转变，高附加值产品正在成为新的增长锚点。以下是我们基于报告逻辑的五个核心洞察。

5月对美出口同比飙升35.4\%，远高于4月的11.3\%，这确实部分得益于去年同期“解放日”关税后的低基数。但BARC的月度环比数据提供了更有说服力的证据：5月对美出口环比增长6.2\%，好于 年约5.6\%的平均水平（剔除2025年关税扭曲）。这意味着对美贸易流正在经历真正的正常化，而非仅仅基数效应的反弹。对美出口对整体增长的贡献从4月的1.2个百分点跳升至3.2个百分点。

然而，更值得关注的是区域贸易的全面激活。对东盟出口增长24.3\%（4月为15.2\%），对日韩台出口增长27.5\%（4月为16.8\%），对非洲出口维持18.6\%的双位数增长。与之形成对比的是，对欧盟出口增速从13.4\%放缓至7.6\%，对英国出口从9.6\%降至1.7\%。这种分化表明，中国出口的韧性并非依赖单一市场，而是建立在多元化的区域贸易网络之上。东盟和东亚经济体与中国供应链的深度绑定，正在形成一种独立于欧美需求的增长引擎。对于产业链上的企业而言，这意味着未来需要同时关注多个区域市场的景气度，而非仅仅盯住中美贸易关系。

报告中最具冲击力的数字来自高科技产品出口。5月高科技产品占总出口比重达到29.8\%，同比增长51\%（4月为39\%）。其中半导体出口价值同比翻倍（+111\%），自动数据处理设备增长66\%。这些数据不是孤立的月度波动，而是全球AI资本开支周期持续发酵的结果。

BARC明确指出了这一逻辑链条：全球AI投资热潮支撑了对中国AI相关出口的需求，因为中国是AI制造组件的关键供应国。这一判断意味着，中国高科技出口的驱动力已经从“替代效应”转向“需求创造效应”。过去，市场对中国高科技出口的理解更多是“美国限制什么，中国就加速生产什么”。但现在，真正推动出口的是全球AI基础设施建设的真实需求。半导体、服务器、通信设备等品类的高增长，反映的是产业链下游的刚性采购。

这对投资者的含义是：中国高科技出口的景气度，现在与全球科技巨头的资本开支计划高度相关。如果主要云服务提供商和AI芯片企业的资本开支增速放缓，中国相关出口将面临压力。但反过来，只要AI投资周期没有出现系统性逆转，中国作为制造环节的关键参与者，其出口增长就有结构性支撑。报告没有完全回答的问题是：中国在AI产业链中的价值获取能力究竟有多强？组件出口的高增长是否伴随着利润率的改善？这需要更细致的微观数据来验证。

与AI相关的

[... middle omitted ...]

口份额是否会面临天花板？但至少在中短期内，能源转型的刚性需求为中国绿色技术出口提供了坚实的支撑。

4. 进口数据揭示了中国经济的双重面孔：需求疲软与价格驱动并存

5月进口同比增长27.4\%，同样超出预期。但BARC的分析指出，这一增长主要由价格效应驱动，而非需求全面回暖。机电产品进口增长38.2\%，半导体进口价值飙升，这与中国高科技出口的扩张逻辑一致——出口导向型企业需要进口中间品来完成生产。真正反映内需强度的商品表现分化明显：农产品进口增速从20\%降至4\%，汽车进口继续萎缩25\%，铁矿石进口量转为下降。

最值得警惕的是原油进口量同比下降29\%，降幅较4月的20\%进一步扩大。尽管天然气进口在连续两个月双位数下降后趋于稳定，但整体能源进口的疲软表明，中国工业活动的实际强度可能低于市场预期。进口数据的“量价背离”是一个重要的观察框架：名义高增长背后，实际需求可能比看起来更弱。对于依赖中国需求的全球大宗商品市场而言，这意味着不能简单地将进口金额增长解读为需求复苏信号。

\subsection{[R029] GS：市场真正低估的不是AI需求，而是2027年资本开支的上行风险}
{\small\color{refgray}归类：AI/算力：资本开支上行风险与供应链结构性重置}

GS：市场真正低估的不是AI需求，而是2027年资本开支的上行风险

当大多数投资者还在为2026年84\%的资本开支增速感到兴奋时，一个更重要的信号正在被系统性低估：2027年的资本开支曲线可能远比共识预测陡峭。这不是一个简单的“多花几百亿”的问题，而是意味着AI基础设施的投资强度正在逼近甚至超越历史上铁路、汽车和电力网络建设时期的GDP占比。

GS这份最新研报的核心判断是：共识预测中2027年超大规模云厂商资本开支仅为9200亿美元、增速骤降至22\%，但这一数字可能被大幅低估。如果增量投资达到GDP的2\%至3\%，2027年资本开支将落在1.1万亿至1.25万亿美元区间；在更为激进的情景下，基于现金流和投资级债券市场容量，这一数字甚至可能达到1.43万亿美元。

这不是一个关于“AI有没有泡沫”的争论，而是一个关于“资本开支曲线形状”的定价问题。市场当前对AI基础设施股票的定价，隐含的是资本开支增速从84\%断崖式下滑至22\%的预期。如果真实曲线更为陡峭，那么整个AI基础设施受益股的盈利和估值框架都需要重估。

理解这份报告的关键，不在于2026年的数字，而在于2027年的假设。

当前共识预测显示，超大规模云厂商2026年资本开支为7570亿美元，同比增长84\%；但到了2027年，这一数字仅微增至9200亿美元，增速骤降至22\%。这种急剧的增速下滑，构成了当前市场对AI基础设施股票定价的核心假设。

问题在于，这个假设的历史记录并不好。GS研报指出，在过去三年中，分析师对超大规模云厂商资本开支的预测平均低估了45个百分点。2024年初，共识预测增速为19\%，实际执行结果为54\%；2025年初预测22\%，实际为73\%；2026年初预测36\%，目前已经上修至84\%。

每一次共识预测都被现实大幅超越，但每一次新的预测又再次回到保守区间。这种模式背后反映的是一种系统性偏差：分析师倾向于用线性外推来预测指数级增长，而AI基础设施的资本开支恰恰处于一个非线性扩张的阶段。

2027年的关键变量在于，超大规模云厂商的营收积压正在以前所未有的速度膨胀。谷歌云和AWS的合并积压订单从六个月前的3580亿美元飙升至8320亿美元。GS股权分析师判断，供需平衡最早要到2027年下半年才能实现。谷歌在最近的财报电话会上明确表示，预计2027年资本开支将“显著高于”2026年，并已通过股权融资筹集了860亿美元。

这些信号合在一起指向一个结论：2027年资本开支增速大概率不会自然减速至22\%，共识预测很可能再次被上修。

2. 1.1万亿到1.4万亿美元的资本开支并非空想，历史和技术逻辑都支持

第一个情景基于历史类比。当前AI增量投资占GDP的比例在2025年为0.9\%，2026年预计达到1.5\%。这个水平已经与1990年代电信和ICT硬件建设的高峰期相当，但远低于铁路建设时期的3.4\%和电力网络建设时期的2.2\%。如果AI投资强度达到GDP的2\%至3\%——这在历史上并非极端值——2027年资本开支将落在9500亿至1.25万亿美元区间。

第二个情景基于股权分析师对token需求的预测。GS预计到2030年，token消费量将增长24倍，主要由企业级AI代理驱动。如果token需求、能源强度和投入成本这三个变量同时上行，名义资本开支的弹性会非常大。

第三个情景基于资产负债表容量。当前超大规模云厂商的净债务与EBITDA之比仅为0.4倍，处于极低水平。理论上，它们可以将净债务增加9500亿美元，同时仍将净杠杆率维持在1倍以下——这对应AA评级企业的上沿。虽然投资级债券市场短期内可能面临单发行人集中度的约束，但GS信贷策略团队认为，离岸债券发行、合资项目融资和私人基础设施市场将填补这一缺口。

这三个情景的共同指向是：2027年

[... middle omitted ...]

注意到，近期AI基础设施股票中出现了杠杆和期权交易的显著增加。在一个已经高度拥挤的交易中，杠杆资金的涌入会放大任何方向的价格波动。上周AVGO财报不及预期和超大规模云厂商股权融资的新闻，已经引发了AI交易组合的剧烈震荡。

一个更值得关注的信号是：AI相关股票的回报离散度在2026年第二季度达到了53个百分点，创下历史新高。这种离散度并非来自多空分歧，而是来自基础设施板块内部收益率的分化。当板块内部分化如此之大时，意味着市场对“谁真正受益”的定价正在经历剧烈调整。

4. 企业AI采用仍处于“谈得多、做得少”的阶段，盈利传导尚未验证

报告中最具洞察力的部分，可能不是资本开支预测，而是对企业AI采用现状的刻画。

2026年第一季度财报季中，约54\%的公司在其电话会上讨论了AI与生产力的话题。但只有11\%的公司量化了AI在具体用例上的生产力提升，仅有2\%的公司量化了AI对盈利的影响。更关键的是，那些谈论AI的公司和没有谈论AI的公司，在利润率和股价表现上几乎没有差异。

\subsection{[R030] JPM：市场低估的不是出口增速，而是贸易结构对增长含义的根本改变}
{\small\color{refgray}归类：中国贸易：价格驱动的结构性重定价与AI主导的出口新格局}

JPM：市场低估的不是出口增速，而是贸易结构对增长含义的根本改变

五月出口数据再次超预期，年增长率达到19.4\%，远超市场共识的15\%和JPM自己预测的12.1\%。但这份JPM研报的核心判断，并非出口有多强劲，而是出口的增长模式正在发生一次被市场低估的质变——增长越来越窄、越来越依赖价格驱动、越来越与国内生产和就业脱钩。这意味着，仅仅盯着出口增速来判断经济走向，已经不够了。

过去几年，市场习惯将出口增长等同于制造业景气，将贸易顺差扩大等同于对GDP的净贡献增加。但这份报告揭示了一个不那么直观的现实：出口增长的主要驱动力，已经从“量”转向了“价”，而且集中在极少数产品类别。当增长主要来自价格效应而非产量扩张时，它对上游生产、就业和资本开支的拉动效应，会显著弱化。

更值得关注的是，这种结构变化发生在全球贸易政策环境高度不确定的背景下。美国正在将关税工具从IEEPA转向301和232条款，欧盟对华贸易立场也在持续硬化。这意味着，中国出口面临的风险已经从“需求端波动”转向了“政策端冲击”。对于产业决策者和投资者而言，理解这种变化，远比预测下个月出口数字更重要。

五月出口环比增长3.1\%，延续了四月6.7\%的强劲势头。但更值得拆解的是增长来源。这份研报明确指出，出口增长主要集中在三个领域：存储芯片与模块、AI数据中心建设设备、以及新能源产品（电动汽车、太阳能、电池）。这三类产品合计贡献了约60\%的出口增量。

关键变化在于，半导体出口价格已经大幅上涨。存储芯片从2025年的“量增驱动”转变为了近几个月的“价增驱动”。“新三样”的价格也在企稳，同时伴随可观的量增。而其他出口产品的平均价格，在四月已经转为小幅正增长，且可能因能源成本上升而面临进一步上行压力。

这意味着什么？出口增长的名义数字虽然好看，但剔除价格因素后的实际量增可能远低于表面水平。对于依赖出口量的制造业企业来说，这意味着收入增长并不一定转化为产量增长和产能利用率提升。对于投资者而言，这意味着需要重新审视出口相关板块的盈利预期——价格上涨带来的收入增长，其可持续性和传导效应与产量增长完全不同。

2. 进口结构揭示了一个更重要的信号：AI供应链拉动和商品囤货，而非内需复苏

进口连续第六个月扩张，五月环比增长2.9\%，年增长率维持在27.5\%的高位。但这份报告做出了一个关键区分：进口增长更多来自AI相关高科技进口和商品囤货，而非广泛的国内需求复苏。

从来源地看，五月进口增长主要由亚洲拉动，韩国环比增长10.2\%，日本增长2.9\%。从产品看，高科技进口环比增长5.9\%，同样受到价格效应推动，尤其是从韩国进口的存储产品。大宗商品进口则出现分化：煤炭、天然气、大豆、铜的进口量增加，而原油和成品油继续下滑。

原油进口下降的部分原因是去年的储备积累以及企业倾向于消耗库存。但报告提出一个值得关注的判断：随着中东冲突持续，如果进口恢复至更“正常”的水平，可能会收紧区域石油市场。

对于投资者而言，进口结构的变化意味着：不能将进口增长简单解读为内需回暖的信号。AI供应链的资本开支和商品囤货行为，与消费驱动的进口增长，对经济和市场的含义截然不同。前者更多反映特定产业的资本密集度提升，后者才意味着广泛的消费复苏。

这是整份报告最核心的洞察。当出口增长越来越窄、越来越依赖价格驱动时，贸易顺差对GDP的净出口贡献，其计算方式和实际含义都变得复杂。

报告指出，出口增长集中在存储芯片、AI存储模块、电动车、太阳能和电池这几类产品，同时半导体出口价格大幅上涨。进口方面，高科技进口同样受价格效应推动。这种“价格对价格”的结构，使得贸易顺差的名义扩大并不一定对应实际产出和就业的增加。

更关键的是，净出口贡献对相对价格变动（芯片价格、大宗商品价格、汇率）和产品组合变得高度敏感。

[... middle omitted ...]

增有限、且高度依赖价格波动的行业。

报告对贸易政策风险的分析，提供了比市场共识更细致的视角。美国法院否决了基于IEEPA和Section 122的关税后，特朗普政府正在转向301和232条款。新的调查涵盖16个经济体的产能过剩问题，以及60个经济体的强迫劳动执法。这些措施的覆盖范围和推进速度，比2017-18年的对华关税周期更广、更快。

虽然特朗普访华带来了短暂的稳定——建立了贸易委员会和投资委员会渠道，并在关键矿产、波音和农业领域达成了初步承诺——但报告指出，这些协议缺乏实质性内容。中国方面将协议称为“初步的”，法院也在审查关税权力。除非新的委员会能在9月习近平主席访美前将对话转化为正式协议，否则贸易风险仍将是长期的不确定性来源。

对于企业而言，这意味着供应链规划面临的不只是关税本身，更是政策路径的不确定性。2017-18年的关税周期有相对清晰的升级路径，而当前的政策工具切换和法院裁决，使得风险变得难以预测。企业需要为多种政策情景做好准备，而非仅盯着关税税率。

\subsection{[R031] Bernstein：半导体设备进口数据揭示的不仅是周期性复苏，更是结构性分化}
{\small\color{refgray}归类：半导体：欧洲设备需求周期尚未见顶与存储供给纪律重塑}

Bernstein：半导体设备进口数据揭示的不仅是周期性复苏，更是结构性分化

当市场还在争论半导体周期何时见顶时，一份来自台湾海关的最新数据正在传递更为精细的信号。2026年5月，台湾半导体设备进口额同比增长23\%，环比下降8\%。这个数字本身并不令人意外——台积电的资本开支计划早已为市场所熟知。但真正值得关注的，不是总量的增长，而是增长的结构：光刻设备、测试设备、清洗设备正在经历截然不同的需求节奏，而这种分化正在重塑全球半导体设备供应商的竞争格局。

Bernstein的这份研报，通过对台湾半导体设备进口数据的拆解，揭示了一个关键判断：市场对日本和欧洲半导体设备公司的估值，尚未充分反映其在先进制程扩张中的差异化受益程度。简单来说，不是所有设备商都能从这轮资本开支周期中同等受益，而那些在测试、光刻等关键环节拥有不可替代性的公司，其定价权正在被低估。

1. 台湾进口数据是观察全球半导体设备需求的领先指标，其信号意义远超台湾本土

台湾作为全球半导体制造的核心枢纽，其设备进口数据从来不只是反映本地需求。Bernstein的研报明确指出，这些数据对于覆盖日本和欧洲的设备公司具有“强相关性”。这背后的逻辑并不复杂：台积电的产能扩张直接决定了ASML、东京电子、爱德万测试等公司的订单节奏，而台湾海关每月公布的设备进口数据，本质上是对这些公司未来收入的提前映射。

从总量看，5月台湾半导体设备进口额约为44亿美元，虽然环比下降8\%，但同比仍保持23\%的增长。更重要的是，三个月移动平均环比增长达到13\%，表明增长的持续性并未被单月波动打断。这组数字支撑了Bernstein对日本和欧洲半导体设备板块的积极看法。

但真正有洞察力的分析，不在于看总量，而在于看结构。Bernstein的研报通过拆解不同设备类别的进口数据，揭示了三个截然不同的需求故事——这正是市场容易忽略的层次。

2. 光刻设备进口的月度波动掩盖了一个季度级的强力反弹，ASML的台湾收入可能远超市场预期

5月台湾光刻设备进口额为5.43亿欧元，环比下降38\%，同比上升21\%。单看这个数字，容易得出“光刻需求降温”的结论。但Bernstein的分析揭示了截然不同的图景：如果把4月和5月的数据合起来看，本季度前两个月的进口额环比增长约80\%。

这个数字的冲击力在于，4月的光刻设备进口额是历史第三高的月度水平，而通常每个季度的第一个月是季节性低点。这意味着，台积电的光刻设备部署正在加速，而不是放缓。Bernstein基于两个月进口数据的回归模型显示，ASML本季度在台湾的系统销售收入可能达到23.3亿欧元，环比增长61\%，同比增长19\%，台湾将占ASML总系统销售收入的37\%。

这个判断的含义是什么？市场对ASML的估值逻辑可能需要重新校准。当市场还在用季度环比波动来判断需求强度时，Bernstein的数据表明，台积电的先进制程扩张正在进入一个加速阶段，而这种加速对ASML的收入弹性可能被低估了。对于持有ASML头寸的投资者来说，这个信号值得认真对待。

3. 测试设备进口持续走强，爱德万测试的季度收入可能跑赢市场一致预期

与光刻设备不同，测试设备的需求节奏更为平稳，但趋势同样强劲。5月台湾从日本和马来西亚进口的测试设备总额为6.13亿美元，环比增长7\%，同比增长34\%，三个月移动平均环比增长4\%。

Bernstein特别指出了测试设备进口与爱德万测试台湾收入之间的高度相关性。基于回归模型，研报估算爱德万测试本季度台湾收入可能环比增长6\%，这意味着其整体收入可能环比增长5\%，而市场一致预期仅为3\%。换句话说，市场可能低估了爱德万测试在这轮周期中的受益程度。

这背后的产业逻辑值得展开。随着先进制程芯片的复杂度持续提升，测试环节的重要性正在被重新定义。每一代

[... middle omitted ...]


与光刻和测试设备的强势形成对比，东京电子在台湾的需求出现了明显的月度回落。5月，与东京电子产品线相关的台湾进口数据（包括CVD、干法刻蚀、清洗、涂胶显影、RTP等）环比下降19\%，同比上升7\%，三个月移动平均环比仅增长3\%。Bernstein的回归模型显示，东京电子本季度台湾收入可能环比下降24\%，远低于市场预期的持平。

这个数据点需要谨慎解读。一方面，单月数据波动在半导体设备行业并不罕见，尤其是当台积电的订单交付节奏出现季度间调整时。另一方面，东京电子在刻蚀和沉积设备领域的全球市场份额正在稳步提升，其技术壁垒并未因短期订单波动而减弱。

这里的核心问题是：市场是否会因为一个季度的台湾收入下滑，就低估东京电子在更长时间维度上的结构性增长？Bernstein给出的答案是否定的——其59200日元的目标价仍高于当前股价，表明长期观点并未改变。但这份数据确实提出了一个需要持续跟踪的问题：东京电子的订单节奏是否正在被台积电的采购策略变化所影响，还是仅仅是季度间的正常波动？

\subsection{[R032] Citi：中国2万亿AI基建计划真正重塑的不是需求，而是国内供应链的定价权}
{\small\color{refgray}归类：AI/算力：资本开支上行风险与供应链结构性重置}

Citi：中国2万亿AI基建计划真正重塑的不是需求，而是国内供应链的定价权

市场正在讨论一个巨大的数字：2万亿元人民币，五年，国家级数据中心网络。但这份Citi研报的核心判断，并非只是“需求会爆发”这么简单。它真正揭示的信号是：中国AI基础设施的竞争逻辑正在发生一次底层切换——从谁建得最快，转向谁能把国产化率80\%的硬约束转化为结构性优势。

国家发改委牵头的这一计划，要求至少80\%的关键技术依赖国内供应商。这意味着，过去市场对AI产业链的估值模型，很大程度上基于全球供应链的假设，如今需要重写。Citi研报在6月9日发布后迅速引发关注，因为它提供了一个量化锚点：以每兆瓦约2亿元的综合成本计算，2万亿投资相当于新增约10GW的AI数据中心容量，年均2GW。而中国目前全部IDC装机容量估计约为20GW。换句话说，仅这一计划就将使AI级数据中心容量在现有基础上扩张约50\%。

这个数字本身已经足够震撼，但真正值得深入推敲的，是它如何重塑产业链上每一家公司的竞争地位。以下是我们从这份研报中提炼出的四个关键层次，以及一个尚未被充分讨论的问题。

1. 运营商从“管道工”升级为AI经济的基础设施运营商，其估值逻辑需要重构

报告明确指出，根据媒体报道，国有电信运营商将运营这批互联数据中心的主体。这并非简单的“接单干活”，而是一次角色跃迁。

过去十年，三大运营商在资本市场上的定位长期被简化为“流量管道”，估值受制于ARPU值增长乏力、资本开支周期和竞争性降价。但Citi研报暗示，当运营商成为国家级AI算力网络的核心运营方，他们的收入结构将从“卖连接”转向“卖算力+连接”。这不仅意味着单位收入的提升，更意味着商业模式的本质变化——从流量批发商升级为平台型基础设施服务商。

这一转变的财务含义尚未被市场充分定价。如果运营商能够基于国产算力集群提供高附加值的算力调度、运维和安全服务，其利润率中枢可能显著上移。而当前市场对运营商的估值，仍然停留在传统的EV/EBITDA框架中，并未计入这一结构性溢价。

2. 中立IDC运营商正在经历一轮“订单质量”的跃迁，而非仅仅是订单数量

一个容易被忽视的细节是，Citi研报特别提到了GDS和VNET在2026年第一季度的创纪录批发订单：GDS约200MW，VNET超过500MW。但关键不在于订单规模本身，而在于这些订单的“质量”——它们来自国家级AI基建计划的前期需求。

传统上，IDC运营商面临的最大风险是客户集中度和需求波动性。但当需求方变为国家级战略项目，且建设周期明确、资金有保障时，订单的确定性和长期性大幅提升。这意味着，对于GDS和VNET这类公司，其收入可见性将从12-18个月延长至3-5年，这直接支撑其融资能力和扩张节奏。

不过，这里仍需注意一个关键假设：这些订单是否真正来自AI基建计划，还是来自其他商业客户？报告没有完全展开。但VNET超过500MW的单季订单量，放在历史数据中看，已经远超常规水平，这本身就是一个需要投资者密切跟踪的信号。

3. 国产AI服务器厂商的真正受益点，在于“供应约束”变成了“需求保障”

报告点名了IEIT、联想和中兴通讯作为AI服务器需求受益者。但更值得思考的是，国产化率80\%的硬约束，对这些厂商意味着什么。

过去，国内AI服务器厂商面临的最大困境是：需求存在，但高端芯片供应受限，导致产能无法充分释放，毛利率承压。而这份计划的核心变革在于，它通过政策手段创造了一个“确定性需求池”，并且这个需求池明确指向国产芯片方案。换句话说，原本的“供应约束”现在变成了“需求保障”——因为客户只能从国产方案中选择，而国产方案的主要供应商就是这几家。

这意味着，这些厂商的AI服务器业务，将从“努力争取订单”转向“分配产能”，其议价能力和盈利稳定性都将提升

[... middle omitted ...]

量价齐升。但Citi研报特别强调了“互联数据中心网络”这一概念。这意味着，国家级算力网络不仅仅是单个数据中心的建设，更是数据中心之间的高速互联。这将对高速光模块（如400G/800G）产生额外的需求增量，而不仅仅是传统的数据中心内部连接。Accelink作为国内领先的光模块厂商，在这一轮互联需求爆发中，其产品组合和价值量都将显著提升。

对于中软国际，受益逻辑更为间接但可能更具持续性。报告将其定位为“华为的IT服务和ISV合作伙伴”，提供算力服务和基于Token的AI服务。这实际上指向了一个更宏大的趋势：当硬件基础设施建成后，算力的“运营”和“调度”将成为新的价值高地。中软国际的角色类似于算力时代的“系统集成商+运维服务商”，其收入模式将从一次性项目制转向持续性的服务订阅。如果这一模式能够跑通，其估值逻辑将发生根本性变化。

报告尚未完全回答的关键问题：国产芯片的性能瓶颈何时会变成政策执行的障碍？

这是整份研报逻辑链条中最脆弱的一环，也是投资者必须自己判断的风险点。

\subsection{[R033] GS：市场对香港银行跨境监管风险的担忧，可能被放大了十倍}
{\small\color{refgray}归类：外汇与资本流动：日元对冲需求与韩元资本外流悖论}

自5月21日以来，HSBC控股和渣打银行的股价分别跑输欧洲银行板块指数5\%和7\%。在一个充斥着政策信号的月份里，这种下跌看起来顺理成章——中国证监会宣布对线上券商加强监管，香港金管局收紧开户流程，国务院发布了关于跨境投资的新规。投资者的直觉反应是：内地资金流入香港的通道正在收窄，依赖这些流量的银行将面临盈利压力。

但GS最新发布的这份报告提出了一个反直觉的判断：市场担忧的方向是对的，但对程度的估计可能错了。 即便在最极端的假设下——来自中国内地的净新资金完全归零——对HSBC和渣打税前利润的影响也仅为约1\%和2\%。这个数字与股价5\%至7\%的跌幅之间，存在一个明显的缺口。

这个缺口意味着什么？是市场过度反应，还是报告遗漏了某些更隐蔽的风险？GS的分析框架值得仔细拆解。

5月22日，中国证监会宣布了对线上券商的监管措施及相关罚款。同日，香港金管局发布通函，旨在加强开户程序和合规要求。6月1日，中国国务院发布了一项关于规范跨境投资的条例。

这组政策在时间上的密集出台，很容易被解读为对跨境资本流动的系统性收紧。但GS的经济学家提供了一个不同的视角：这些政策的核心关切是技术跨境转移的管理，而非限制资本外流。

具体而言，国务院的条例重点在于建立规范跨境技术转让的框架，支持中国企业走出去，并管理与海外业务和资产相关的地缘政治风险。条例中唯一与银行直接相关的条款是第33条，该条款规定“境外金融市场投资应遵守本条例，具体规则另行制定”。

关键短语是“另行制定”。这意味着，截至目前，针对银行跨境业务的实质性新规则尚未出台。GS坦承，他们尚无法预判这些规则的具体形式。但报告同时指出，对于HSBC和渣打来说，直接来自内地境内的净新资金占比“绝对是小部分”。

这里的逻辑链条是：如果政策的首要目标不是限制资本流动，且银行来自内地的直接资金敞口有限，那么市场对银行业务冲击的担忧可能被前置了。

GS的量化分析是这份报告最值得关注的部分。它提供了一个清晰的敏感性框架。

首先，报告拆解了香港银行内地客户存款的资金来源结构。根据HSBC和渣打的管理层指引，一个典型的持有香港账户的内地居民，其存款流中约90\%来自“离岸”来源（例如新加坡或香港本地），只有约10\%来自内地银行账户。

2025年，HSBC的集团净新资金约为860亿美元，渣打约为520亿美元。在近期的投资者会议上，渣打CFO指出，其净新资金中约30\%来自“全球华人”，而其中绝大部分资金已经存放在离岸。

基于这些数据，GS进行了压力测试：假设集团净新资金中约10\%来自内地境内来源，如果这部分未来净新资金完全停止流入，对HSBC和渣打分别意味着约90亿和50亿美元的资金缺口。

进一步，假设在这部分净新资金上赚取100个基点的利差，对税前利润的影响仅为HSBC的0.3\%和渣打的0.7\%。即便以300个基点的利差计算——当前HIBOR为2.7\%——对税前利润的影响也仅为每年约1\%和2\%。

*这个测算的隐含含义是：即便最坏情景发生，对银行基本面的冲击也远小于股价已经反映的幅度。** 市场可能混淆了“风险方向”和“风险程度”。

GS认为，这些产品目前要么不在新规的管辖范围内，要么已经符合现有规则。更重要的是，对这些产品采取更严格的限制，与政策制定者的几个核心目标并不一致：a）监管技术跨境转移的重点，b）支持人民币国际化的意愿，c）继续发展香港作为双向投资枢纽的定位。

这是一个被市场情绪忽略的结构性判断。如果政策制定者确实希望维持香港作为离岸人民币中心和财富管理枢纽的地位，那么对保险和投资产品的过度限制将与这一目标自相矛盾。

当然，报告没有完全展开的是：政策的不确定性本身就有价值。即便最终影响有限，在规则明确之前，银行的新客户获取成本和合规成本可能会上升。这部分“摩擦成本

[... middle omitted ...]

管理产品的偏好或信任度，那么财富收入的波动可能比存款利差的影响更大。

GS的测算仅聚焦于净新资金减少对利差收入的影响，但没有对财富管理手续费收入进行同样的压力测试。原因可能是报告所说“这些产品目前不在范围内”，但投资者需要追问的是：如果客户情绪变化导致财富管理产品的销售周期延长或转化率下降，对营收的冲击会是多少？

对于关注这两只股票的投资者，GS的分析提供了一个更精细的观察框架：不要只盯着净新资金的绝对规模变化，而要关注资金来源的结构。

报告的核心洞见之一是，HSBC和渣打的净新资金中，来自内地境内的直接流入占比很小。这意味着，即便未来跨境监管收紧，对银行整体资金池的冲击也是有限的。真正需要跟踪的指标应该是：

离岸华人客户的净新资金增速是否出现放缓 新开户客户中非香港居民的比例变化 合规成本上升是否影响了新客户的获取效率

如果这些指标保持稳定，那么股价的下跌可能提供了更好的入场窗口。反之，如果离岸华人客户的资金增速也出现系统性放缓，那才意味着更根本性的变化。

\subsection{[R034] DB：中国车市正在经历的，不是周期波动，而是竞争结构的根本性重构}
{\small\color{refgray}归类：中国新能源车：渗透率突破60\%后的竞争结构重构}

DB：中国车市正在经历的，不是周期波动，而是竞争结构的根本性重构

五月的数据看起来有些矛盾。乘用车批发总量同比下降4.9\%，但新能源车（NEV）批发量却逆势增长12\%，渗透率首次突破60\%大关。市场似乎同时在经历“总量收缩”和“结构加速”两个方向。

DB最新发布的这份月度批发量图册，提供了理解这一矛盾最直接的切口。它没有长篇大论地分析宏观政策或消费者信心，而是用一组清晰的数据，揭示了一个正在发生的底层变化：中国乘用车市场的竞争，已经从“谁能造出电动车”进入了“谁能把规模优势转化为定价权和成本护城河”的阶段。

这份报告最值得关注的判断，不是五月销量好坏，而是渗透率突破60\%这个数字背后所代表的竞争格局质变。当新能源车成为绝对主流，市场的游戏规则正在被重新定义——赢家通吃的逻辑更加明显，而落后者的生存空间正在被快速压缩。

1. 新能源渗透率突破60\%不是终点，而是竞争烈度升级的起点

五月新能源渗透率达到60.9\%，环比提升3个百分点，同比提升超过9个百分点。这个数字的意义，远不止于市场份额的里程碑。它意味着，在当前的月度销售中，传统燃油车已经退居次位，新能源车正式成为市场的主流选择。

从更长的时间线来看，这一趋势的斜率值得重视。2025年初，渗透率还在42\%左右，到年中攀升至50\%以上，并在2025年下半年稳定在52\%-56\%的区间。进入2026年，渗透率在经历了一季度季节性回落后，在四月和五月快速拉升，从48\%跳升至61\%。这种加速态势，通常不是由单一产品周期驱动的，而是反映了消费者认知、基础设施、以及产品性价比三者之间的正向循环已经形成。

对于产业决策者而言，这个数字的含义是清晰的：市场不再需要讨论“是否要电动化”，而是必须回答“在电动化已成定局的竞争环境下，我的规模、成本和品牌护城河在哪里”。那些仍然依赖燃油车利润输血的企业，时间窗口正在关闭。

2. 总量下滑掩盖了头部企业的结构性优势，比亚迪的份额反弹是信号

五月行业总量同比下降4.9\%，前五个月累计同比下降6.0\%。但在这个收缩的总量中，头部企业的表现出现了明显的分化。

比亚迪在经历了年初几个月的低迷后，五月批发量恢复至38.3万辆，同比增长0.3\%，环比增长19.4\%，市场份额从四月的14.0\%回升至17.0\%。这个数字尤其值得关注，因为比亚迪年初至今累计销量仍同比下降20.3\%，但五月的强势反弹表明，其产品更新周期和价格策略正在重新生效。

与此同时，特斯拉中国的表现同样稳健。五月批发量8.6万辆，同比增长39.4\%，市场份额稳定在3.9\%左右，接近历史高位。特斯拉的市场份额曲线在过去一年中呈现稳步上升趋势，从2025年初的3.0\%左右，逐步攀升至当前的3.9\%。

这些数据合在一起，揭示了一个关键判断：在总量收缩的环境中，头部企业正在通过规模效应和品牌溢价，进一步巩固自己的市场地位。对于比亚迪来说，五月的反弹可能意味着其年初以来的库存调整和新产品爬坡已接近完成；对于特斯拉，稳定的市场份额增长则体现了其在中国本土化生产后的成本优势持续释放。

新势力企业的五月数据，是这份报告中最能体现竞争结构变化的章节。

零跑汽车是最大亮点。五月批发量8.2万辆，同比增长81.0\%，环比增长14.3\%，市场份额从年初的1.6\%跃升至3.7\%，已经超越了理想汽车和蔚来。零跑的增长曲线在过去一年中几乎呈线性上升，从2025年初的月销2.5万辆，到当前突破8万辆，增长幅度在所有新势力中最为突出。

蔚来同样表现强劲。五月批发量3.8万辆，同比增长62.3\%，环比增长28.4\%，市场份额达到1.7\%，创下近期新高。蔚来的销量曲线显示，其在2025年下半年经历了一轮显著的增长，月销从2万辆左右攀升至4万辆以上，尽管2026年初有所回落，但整体趋势仍然向上。

[... middle omitted ...]

定位和产品定义，在10-20万元这一最主流的价格带实现了快速放量。而理想和小鹏的困境则说明，当市场渗透率超过60\%，单纯依靠品牌故事或单一产品亮点，已经不足以维持增长。

DB的这份图册提供了丰富的销量和市场份额数据，但它没有深入讨论一个核心问题：在渗透率突破60\%之后，行业何时能够迎来盈利能力的系统性改善？

当前的数据显示，尽管新能源车销量在增长，但多数新势力企业仍然处于亏损状态。零跑的增长速度令人瞩目，但其毛利率是否能随着规模扩大而同步改善，仍是一个需要验证的问题。蔚来的销量增长同样强劲，但高昂的研发和销售费用，使其盈利时间表仍存在不确定性。

另一个未被充分讨论的问题是：当渗透率达到60\%以上，传统车企的燃油车业务将面临怎样的压力？上汽集团五月销量同比下降8.4\%，长安汽车同比下降28.4\%，这些下滑不仅仅是个别企业的经营问题，更反映了整个燃油车产业链正在经历的收缩。这种收缩会如何影响传统车企向新能源转型的资金来源和转型节奏，是产业决策者需要密切关注的风险点。

\subsection{[R035] DB：中国新能源车市最值得关注的不是总量，而是订单结构的剧烈分化}
{\small\color{refgray}归类：中国新能源车：渗透率突破60\%后的竞争结构重构}

DB：中国新能源车市最值得关注的不是总量，而是订单结构的剧烈分化

DB最新发布的周度订单追踪显示，6月首周（W23）主要新能源车企合计订单环比下滑约5\%，同比更是下降了约15\%。如果只看这个数字，很容易得出“需求疲软”的结论。但这份报告的价值，恰恰在于它揭示了总量数字之下，各家车企正在经历的截然不同的命运。

真正重要的判断是：中国新能源车市场的竞争，已经从“谁能卖得更多”进入了“谁的订单结构更健康”的阶段。市场目前定价的核心变量，并非整体需求是否复苏，而是企业能否在订单波动中证明自己的定价能力和品牌黏性。

为什么这个时间点重要？因为5月刚经历了一轮由新车型发布和促销活动驱动的订单脉冲，6月首周的回落，恰好是检验哪些品牌拥有真实需求、哪些只是靠一次性刺激撑起数据的试金石。

6月首周，比亚迪订单量为4.77万辆，环比下滑2\%，同比骤降36\%。这是所有主要车企中同比降幅最大的。更值得警惕的是，这一数字不仅远低于去年同期的7.5万辆，也低于今年5月最后一周的4.87万辆。

表面上看，这是高基数效应——去年6月正值比亚迪多款冠军版车型密集上市。但问题不止于此。从更长的时间序列看，比亚迪的周度订单自3月第4周达到8.5万辆的峰值后，便进入了下行通道。尽管4月第1周仍维持在6.5万辆，但此后多数时间徘徊在4.5-5.5万辆区间。

这意味着什么？比亚迪的订单中枢正在下移。市场此前对比亚迪的定价，很大程度上基于其“规模效应持续放大”的叙事。但如果订单量无法维持在高位，规模效应的边际贡献就会递减，而固定成本分摊的压力反而会上升。

这里需要指出一个报告没有完全展开的隐含含义：比亚迪的订单下滑，可能是其产品矩阵内部竞争加剧的结果。当王朝网和海洋网的产品线越来越密集，消费者在不同车型之间的选择反而可能延迟决策，甚至出现“内部替代”效应。这个问题，目前尚未被市场充分定价。

2. 问界的“过山车”订单，揭示出华为生态的爆发力与脆弱性并存

HIMA（鸿蒙智行，主要以问界为代表）6月首周订单2.42万辆，环比大幅下滑41\%，但同比增长112\%。这个数据放在月度趋势中看尤为戏剧性：4月第4周，问界订单曾飙升至6.5万辆；5月最后一周又回升至4.1万辆；然后6月首周迅速回落。

这种脉冲式订单模式，在问界身上已经重复出现多次。每次新车型发布或重大OTA升级前后，订单都会出现一个尖峰，随后快速回落。这与特斯拉早期的模式有些相似——品牌关注度极高，但消费者的购买决策受新闻周期和营销事件驱动明显。

关键在于，这种模式的可持续性面临两个考验。第一，每次脉冲的峰值是否在衰减？从数据看，4月第4周的6.5万辆明显高于3月第4周的1.1万辆，但5月第4周的4.1万辆已经低于4月峰值。第二，脉冲之间的“谷底”是否在抬升？目前来看，问界的周度订单在非高峰期已从年初的2-3万辆提升到3-4万辆，这是一个积极信号。

报告没有直接给出结论，但数据暗示：问界正在从“纯事件驱动”向“品牌驱动”过渡，但过渡尚未完成。如果接下来两个月没有重大新品发布，问界的订单能否稳定在3万辆以上，将是检验其品牌黏性的关键窗口。

蔚来是6月首周最大的亮点：订单2.8万辆，同比增长419\%，环比仅下滑28\%。更惊人的是，就在5月最后一周，蔚来订单曾达到3.9万辆——这个数字甚至超过了比亚迪同期的4.8万辆。

从更长时间轴看，蔚来订单的爆发始于5月第3周，从之前的周均1万辆左右直接跳升到2.1万辆，随后在第5周达到3.9万辆的峰值。这种陡峭的增长曲线，通常对应着新车型上市或重大价格调整。报告没有明确指出具体原因，但从市场公开信息可以推断，蔚来子品牌乐道的首款车型L60的预售，以及蔚来主品牌在BaaS电池租赁方案上的调整，共同推动了这一波订单潮。

但这里有一个报告尚未回答的

[... middle omitted ...]

。这两家曾经的新势力头部玩家，如今在订单量上已经被蔚来、问界甚至小米甩开。

理想的困境更为典型。从月度趋势看，理想在5月第3周曾短暂冲高至1.35万辆，但随后迅速回落。这种“冲高回落”的形态，在理想L6上市后的表现中尤为明显。L6作为理想目前售价最低的车型，确实在上市初期带来了订单增量，但这种增量似乎正在快速消耗。

问题出在哪里？理想的“家庭定位”和“增程路线”曾经是其护城河，但如今问界M7、比亚迪唐等竞品在相同价位段提供了更丰富的智能化和舒适配置。当产品力差距缩小，品牌溢价就无法维持。小鹏的情况类似，其G6和X9虽然产品力不俗，但在价格战中没有建立起清晰的品牌认知。

这两个案例共同指向一个结论：在2026年的中国新能源车市场，“性价比”已经不是一个有效的竞争策略。因为所有车企都在卷价格，消费者对“便宜”的感知阈值越来越高。真正能带来订单黏性的，是品牌定位的独特性、技术体验的差异化，以及用户社群的真实活跃度。理想和小鹏在这三个维度上，目前都还没有找到新的突破口。

\subsection{[R036] GS：中国医疗设备招标的拐点正在到来，但真正的分化才刚刚开始}
{\small\color{refgray}归类：地产与消费：中国餐饮品牌分化加速与家电估值逻辑切换 / 医疗健康：药价谈判结构性变化与设备招标拐点}

GS：中国医疗设备招标的拐点正在到来，但真正的分化才刚刚开始

2026年5月的数据，让关注中国医疗设备行业的人稍微松了一口气。GS最新发布的招标跟踪数据显示，在经历了连续五个月的同比下滑之后，九大类医疗设备的总招标金额在5月录得了0.1\%的同比正增长。这个数字不大，甚至可以说刚刚转正，但它所承载的信号意义，远比数字本身重要。

这份报告的真正价值，不在于告诉你“5月数据转正了”，而在于它揭示了一个正在发生的结构性变化：中国医疗设备市场正在从“总量承压”转向“结构分化”。那些能在招标总量恢复的过程中，同时实现份额提升和海外扩张的公司，将获得远超行业平均的增长弹性。而那些仅仅依赖国内招标总量回暖的企业，可能会在接下来的竞争中感受到更大的压力。

为什么这个时间点值得关注？因为2026年5月的转正，是在2025年同期已经处于相对高基数的基础上实现的。GS的报告明确指出，这一变化符合其此前“行业将从温和同比下降转向温和同比增长”的预期，并且预计更为明显的复苏将在2026年下半年出现。这意味着，市场的底部可能已经确认，但接下来的上行路径，不会是普涨。

5月的数据表面上是总量层面的好消息，但拆开来看，结构性的分化才是核心。在九大类设备中，只有CT、MRI和直线加速器（LINAC）实现了同比正增长。这三类设备有一个共同特征：它们都属于大型、高单价、采购决策周期较长的设备。它们的恢复，可能意味着医院端对于大型资本开支的态度正在从“极度谨慎”转向“有条件放开”。

与之形成对比的是，此前连续保持正增长的患者监护仪，在5月未能延续强势。这看似是一个负面信号，但GS在报告中给出了一个值得深思的观察：疫情期间大量采购的监护仪，现在正进入替换周期。这意味着监护仪的需求逻辑正在从“应急采购”转向“存量替换”，这种需求的持续性可能比市场预期的更强。不过，报告并没有完全展开这个替换周期的具体节奏和规模，这里仍需验证。

超声、内窥镜、DR等品类的招标数据仍然疲软。超声在5月同比下降20\%，内窥镜同比下降25\%，DR同比下降26\%。这些品类的共同问题是：它们更多依赖基层医疗机构的采购和日常诊疗量驱动，而非大型医院的资本开支计划。它们的恢复，可能需要对基层医疗财政支持力度有更明确的信号。

所以，5月数据转正的真正含义是：大型设备率先企稳，但中小型设备的拐点尚未到来。对于投资者而言，这意味着不能简单用“行业复苏”来概括所有公司的前景。

2. CT和MRI的招标恢复：是需求回补，还是医院信心修复？

CT和MRI是5月招标数据中最亮眼的两类。CT招标金额同比增长12\%，MRI同比增长20\%。这两个数字，是在4月分别同比下降31\%和50\%的背景下实现的。这种大幅波动本身说明，大型设备招标的月度数据存在较大的噪音，单月数据不宜过度解读。

但如果我们把时间拉长，看今年前五个月的整体趋势，CT和MRI的招标金额仍然处于一个从低位缓慢爬升的过程中。GS的报告没有给出前五个月的累计同比数据，但从图表趋势来看，CT和MRI的招标金额在2025年下半年经历了一轮显著回升，进入2026年后有所回落，5月再次反弹。这种“W型”走势，更像是医院在财政预算约束下的“按需采购”行为，而非全面的需求爆发。

这里有一个关键问题：医院端的采购意愿是否已经真正修复？GS的报告没有直接回答这个问题，但我们可以从另一个角度来推断。报告提到，联影医疗在今年前五个月的招标金额同比增长了9\%，是主要竞争对手中唯一实现正增长的企业。如果整个行业的需求仍然疲软，联影很难单独实现增长。这说明，至少在CT和MRI这两个品类上，需求端的恢复是真实存在的，只是恢复的力度和节奏在不同企业之间出现了分化。

联影医疗的表现，是这份报告中最值得深入分析的部分。在GEHC、西门子医疗、飞利浦、东软等主

[... middle omitted ...]

相当激进的指引。GS目前对联影2026年的营收增长预测是19.3\%，略低于公司指引。考虑到国内招标需求仍处于恢复早期，这个预测是合理的。但值得关注的是，联影的海外业务正在成为新的增长引擎。报告显示，联影海外收入占比已经超过25\%，且2026年海外增速预计仍将高于国内增速。如果海外市场能够持续贡献增量，那么联影实现20\%以上增长的可能性就会显著提升。

这里有一个尚未完全回答的问题：联影的份额提升，在多大程度上是可持续的？是依靠产品力和服务能力的长期竞争优势，还是短期价格策略或政策导向的结果？GS的报告没有展开讨论，但这恰恰是判断联影长期估值中枢的关键。

4. 迈瑞医疗的国内去库存何时结束：一个正在被市场低估的拐点

迈瑞医疗的情况，与联影形成了有趣的对比。联影在招标端实现了增长，而迈瑞在国内市场仍处于去库存阶段。报告显示，迈瑞1Q26国内收入同比下降11\%，其中PMLS（病人监护、生命支持）业务同比下降44\%，医学影像业务同比下降37\%。这些数字看起来仍然令人担忧。

\subsection{[R037] HSBC：储能正在成为全球能源政策的新重心，光伏反而面临需求拐点}
{\small\color{refgray}归类：能源与大宗：欧洲钢铁定价权重塑与储能政策重心转移}

HSBC：储能正在成为全球能源政策的新重心，光伏反而面临需求拐点

2026年6月初的上海SNEC光伏展，传递出的信号并不像展馆内的热闹那样均衡。HSBC在其最新SNEC观展报告中给出了一个层次分明的判断：地面光伏正面临中国市场化定价改革带来的需求转弱，而储能——尤其是户用储能——正在成为全球能源政策切换的新重心。这一判断并非简单的“光伏不行了、储能行了”的二元叙事，而是建立在各国电网消纳瓶颈、政策补贴结构变化以及技术渗透阶段差异之上的结构性推演。

对于关注中国新能源产业链的投资者而言，这份报告的核心价值不在于复述展会上的产品陈列，而在于它揭示了一个正在发生的政策优先级转移：全球范围内，从“鼓励光伏装机”到“强制配套储能”的过渡，正在从早期市场向更多国家蔓延。这意味着，储能赛道的选股逻辑，需要从“跟着光伏走”切换到“独立于光伏、甚至受益于光伏瓶颈”的新框架。

HSBC在报告中明确指出了地面光伏面临的两个压力源。第一个是中国的政策切换：2025年6月，中国从固定上网电价转向市场化定价机制。这一变化的影响在一年后的SNEC上已经可以感知——参展商对地面项目需求的预期明显趋于谨慎。第二个压力来自竞争格局的恶化：非传统玩家如京东方和阳光电源都在SNEC上展示了光伏组件产品。这意味着组件环节的竞争正在从传统光伏制造商之间的“内卷”，升级为跨界巨头的“降维打击”。

这两点叠加在一起，指向一个结论：地面光伏的利润率压缩可能不是周期性的，而是结构性的。市场化定价意味着电站开发商不再享有稳定的电价保证，回报率的不确定性会抑制新项目开工；而跨界玩家的进入，则意味着即便需求恢复，价格竞争也难以缓解。

对于投资者而言，这意味着光伏制造端的选股逻辑需要从“规模扩张”转向“成本曲线和差异化”。HSBC对福斯特的评级为持有，隐含的正是对短期空间太阳能需求谨慎、对传统组件竞争格局担忧的综合判断。

HSBC报告中最重要的数据点之一，是中国锂电池对前十大市场的出口在2026年前四个月同比增长29\%。这一增长并非均匀分布——对澳大利亚的出口增长了1.9倍，对荷兰增长了1.48倍，对英国增长了23\%。这些数字背后，是各国政策正在从“补贴光伏”转向“补贴储能”的明显信号。

报告特别提到了两个关键政策节点。澳大利亚的“Cheaper Battery Program”虽然在2026年5月开始退坡，但英国的“Warm Home Plan”（150亿英镑用于屋顶光伏/储能/热泵）将在2026年底生效。更重要的是，HSBC在报告中绘制了一张包含20多个国家的“储能政策阶段图”，清晰地展示了全球从净计量（net metering）向净结算（net billing）乃至虚拟电厂（VPP）演进的路径。

这张图的含义是：当一个国家光伏加风电渗透率接近两位数时，电网瓶颈会成为制约可再生能源进一步发展的硬约束。解决这一约束的唯一办法，就是加速储能部署。因此，全球政策优先级从“促光伏”转向“促储能”不是偶然，而是电网物理规律的必然结果。

这可能是报告中最容易被忽视但最具杀伤力的信息点。2026年4月22日，欧盟委员会禁止欧盟资金支持使用中国逆变器的可再生能源项目。HSBC估算，这一禁令可能影响欧洲约20\%的地面光伏和储能市场，而对工商业项目的影响约为5\%。

这一影响的非对称性值得深思。地面项目通常依赖项目融资，融资成本的变化会直接改变项目的IRR。如果使用中国逆变器导致项目无法获得欧盟资金支持，那么开发商的理性选择要么是改用非中国逆变器（成本上升），要么是放弃项目。无论哪种选择，对阳光电源等中国逆变器龙头的欧洲地面项目业务都会形成压力。

而工商业项目对融资的依赖度较低，影响相对可控。这意味着，对于中国储能企业而言，欧洲市场的增长驱动力正在从“地面项目”转向

[... middle omitted ...]

小，其商业化速度反而可能更快。

对于空间钙钛矿，HSBC认为在地球上的应用（如BIPV、智能门锁电源、便携式电子设备）会比在太空中更快，因为太空应用需要足够多的在轨数据来验证可靠性，这一过程可能耗时数年。

这里的投资启示是：技术路线选择上，市场往往高估了短期内突破性技术的渗透速度，低估了渐进式改进的累积效应。HSBC对宏发股份的看好，部分逻辑正来自于此——800V高压直流继电器市场40\%的全球份额，意味着宏发是HVDC技术推广的直接受益者，而这一技术路线比固态变压器更可能在未来2-3年放量。

5. HSBC尚未完全回答的关键问题：储能电池产能过剩的消化路径

报告对亿纬锂能给出了101\%的上涨空间，核心逻辑是市场份额增长和产能扩张。但报告没有完全展开的一个问题是：全球储能电池产能过剩的消化路径。亿纬锂能计划新增约260GWh的LFP电池产能，而德业、阳光电源等系统集成商也在扩大产能。在需求端，虽然户用储能增长强劲，但地面项目面临欧盟禁令和融资环境收紧的双重压力。

\subsection{[R038] HSBC：SNEC 2026揭示的真正信号，不是光伏回暖，而是能源系统的权力交接}
{\small\color{refgray}归类：能源与大宗：欧洲钢铁定价权重塑与储能政策重心转移}

HSBC：SNEC 2026揭示的真正信号，不是光伏回暖，而是能源系统的权力交接

每年SNEC都是中国光伏行业的一次集体亮相，但今年的信号与往年截然不同。HSBC分析师团队在2026年6月初实地走访上海SNEC后，给出的核心判断并非关于组件价格何时触底，也不是产能出清到了哪个阶段——而是整个行业的叙事重心正在发生一次根本性转移：光伏不再只是发电源，而正在成为AI算力基础设施的底层能源架构。

这份报告的价值不在于告诉读者“行业还在调整”，而在于它捕捉到了一个正在发生的结构性变化：当传统光伏供应链的展台人流稀疏、讨论沉闷时，固态变压器、AI数据中心供电一体化方案、以及住宅储能软件的智能化，正在成为新的流量中心。HSBC的观察暗示，投资者真正需要关注的不再是组件出货量或电池效率突破，而是“谁能把光伏从发电设备重新定义为算力能源接口”。

这并非一个遥远的愿景。报告指出，在SNEC现场，“绿色电力驱动AIDC”已经成为大量展台的核心叙事，从渲染图到参考架构再到实物演示，生态系统的参与度远高于一年前。而固态变压器（SST）作为连接光伏直流发电与数据中心配电的关键环节，正在成为多个企业竞相布局的焦点。

以下是我们基于HSBC这份实地调研报告提炼出的五个核心洞察，以及一个尚未被充分回答的关键问题。

1. 光伏供应链的“冷静”并非坏事，它意味着资本正在从制造端向应用端迁移

HSBC团队的直观感受是，与储能和AIDC主题的火爆相比，传统光伏供应链展区“明显更安静”。这不是一个孤立的观察。当组件价格仍在RMB /W的主流区间承压，分布式和海外市场能清到RMB /W，而更低价的订单更多集中在年初时，整个制造环节的定价权正在被持续压缩。

但“安静”不等于“没有信号”。HSBC报告揭示了一个关键的结构性转变：设备投资周期正从新建产能转向升级改造。累计TOPCon产能已超过800GW，这意味着新线建设的高峰已经过去，取而代之的是每GW约RMB3000万的改造投资，用以将组件功率从620W提升至650W以上。这是一个典型的“存量优化”叙事——资本支出强度下降，但技术迭代仍在进行。

对于投资者而言，这意味着光伏制造环节的边际价值正在从“扩产”转向“提效”。铜基浆料替代银浆、降低金属化成本等工艺优化，每瓦可带来RMB 的成本下降，但规模爬坡需要时间。那些能够在这个“优化而非新建”的周期中持续降本的企业，将比单纯追求出货量的公司更具防御性。

但更值得关注的是，资本和注意力的转移方向。当传统供应链变得“安静”，而储能和AIDC展区人流密集时，行业正在用脚投票：下一轮价值创造的主战场，不在制造端，而在应用端和系统集成端。

2. 住宅储能正在经历“硬件同质化后的渠道与软件之战”，AI能源管理已成标配

HSBC报告指出，住宅储能是SNEC现场人流最密集的展区之一。但最值得关注的不是热度本身，而是竞争格局的变化方向。报告明确提出，硬件差异化正在收窄，竞争正在向渠道、认证、交付和软件服务转移。

“AI驱动的能源管理”已经不再是少数厂商的差异化卖点，而是被大量展商定位为“基线能力”。这意味着，能够实现电价套利、智能调度、预测控制和虚拟电厂参与等功能的系统，正在从“锦上添花”变成“入场券”。同时，一体化系统的普及正在简化安装流程和用户体验，进一步降低了新进入者的门槛。

这个趋势对行业格局的含义是双重的。一方面，硬件的商品化会加速价格竞争，压缩制造环节的利润空间。另一方面，软件和渠道能力的价值正在被重新定价——谁能通过AI算法更精准地管理用户侧储能，谁就能在电力市场化和虚拟电厂崛起的背景下，获得更高的客户粘性和服务溢价。

对于投资者而言，评估住宅储能企业的框架需要从“出货量增速”转向“软件订阅收入占比”和“渠道覆盖密度”。硬件制造的规模效

[... middle omitted ...]

口问题：直流耦合与系统级集成。传统的数据中心配电架构基于交流变压和UPS系统，而光伏发电天然是直流电，AI芯片的供电也越来越倾向于直流微网架构。SST作为连接这两端的核心器件，其效率和可靠性直接影响数据中心的PUE（电能使用效率）和整体运营成本。

然而，HSBC报告也给出了一个必要的审慎判断：尽管兴趣浓厚，但部署速度预计会较慢。原因在于，SST的生命周期经济性、效率表现、对PUE的实际影响、可靠性数据以及标准化进程，都还需要时间来验证和证明。

这意味着，SST目前处于一个典型的“预期先行、业绩滞后”的阶段。对于投资者来说，关键不是判断SST会不会成为主流——从技术逻辑看，这几乎是必然的——而是判断谁能在标准化和成本下降过程中占据先机。报告没有给出明确的胜出者，但指出“多家企业”正在竞逐，这意味着竞争格局尚未定型，技术路线和客户验证将成为未来12-18个月的核心观察变量。

4. 太空太阳能仍然“小众但旗舰”，商业化时间表指向2027年，SpaceX的目标被谨慎对待

\subsection{[R039] JPM：中国家电的估值逻辑正在从“周期股”转向“复合增长引擎”}
{\small\color{refgray}归类：地产与消费：中国餐饮品牌分化加速与家电估值逻辑切换}

JPM：中国家电的估值逻辑正在从“周期股”转向“复合增长引擎”

市场对家电板块的定价，可能还停留在旧世界里。10倍市盈率、6\%的股息率、约20\%的净资产收益率，这是当前市场给中国家电三巨头——美的、海尔、格力——的集体定价。这个定价隐含的假设是：这些公司本质上是国内耐用消费品周期股，补贴退坡后需求放缓，利润承压，估值没有扩张空间。

JPM最新发布的家电行业深度报告提出了一个根本性的问题：如果中国家电不再应该被当作一个国内家电周期来估值，而是作为“现金牛盈利 + 全球挑战者增长 + 工业科技期权价值”的组合，那会怎样？

这个问题的答案，决定了当前估值体系是否正在酝酿一次系统性重估。报告的核心判断是：市场低估的不是短期需求，而是供给侧的再定价——即龙头企业能否将中国市场的现金流，转化为全球份额扩张和B2B能力建设。而在这条逻辑线上，美的集团被选为最清晰的标的。

国内家电市场正在成熟，这是共识。补贴政策退坡后，2026年国内零售需求可能出现中个位数的下滑。JPM预计，二季度将是利润率压力测试的关键窗口：需求走弱已充分预期，但石化产品驱动的成本上升将真正考验龙头企业的定价能力和成本传导能力。基准情景下，三巨头在2026年二、三季度的毛利率将温和收缩0.5至1个百分点。

但这里的关键判断不是“利润会不会跌”，而是“龙头企业能否守住利润池”。报告指出，三巨头在国内市场合计控制了约60\%的份额，核心利润池相对稳固。如果它们能在二、三季度再次证明毛利率的相对稳定，“成熟品类叠加利润率压缩”的熊市叙事就会失去说服力。

这意味着，国内B2C业务的角色正在转变：它不再是增长的主要驱动力，而是变成了一个高现金回报的“燃料箱”。这个燃料箱能否被有效再配置到更高增长的领域，才是投资者真正应该关注的变量。

国内增长放缓，并不等于整个行业的增长曲线已经平坦。海外市场的结构性机会，是这份报告最有力的反驳论据。

数据很清楚：三巨头在国内的市场份额合计超过60\%，但在中国以外的全球市场，它们的合计份额只有约16\%。而海外可寻址市场大约是国内市场的三倍。这个差距本身就是增长空间。

更重要的是，增长的驱动力正在从“出口量”转向“OBM（自主品牌）规模、渠道控制和品牌经济性”。JPM预计，2026至2028财年，美的、海尔、格力的海外收入复合年增长率将分别达到11\%、9\%和7\%。这个增速意味着，即使国内需求趋于平稳，这些公司仍然能够实现持续的收入复合增长。

海外业务不再是“锦上添花”，而是正在成为新的增长引擎。这一点对于估值框架的转变至关重要——因为全球投资者对“中国出口型公司”的估值逻辑，与对“全球品牌运营商”的估值逻辑，是完全不同的两套体系。

如果说海外业务是“增长引擎”，那么B2B业务就是“估值乘数”。这是整份报告中最具洞察力的部分，也是市场定价与公司实际价值之间差距最大的地方。

商用暖通空调、数据中心液冷、欧洲热泵——这些业务看起来与传统家电无关，但它们恰好坐落在三巨头的核心能力之上：热管理、零部件深度、供应链控制、制造规模和成本下降执行力。全球商用暖通空调市场规模约1.2万亿元人民币，而三巨头合计份额仅约5\%。如果执行力能够维持，这是一个极长的增长跑道。

美的集团是最清晰的案例。到2026财年，其B2B业务预计将接近总收入的30\%。但市场仍然以“家电公司”的框架给它定价，几乎没有为这种业务结构的变化给予估值溢价。这个差距，就是重估的机会。

这里有一个值得深思的类比：报告专门用了一个章节的标题来提问——“美的到2030年：西门子，还是松下？”西门子代表了从工业集团转型为科技驱动型工业公司的成功案例，估值溢价显著；松下则代表了未能完成转型、被市场长期低估的教训。美的正站在这个十字路口，而市场目前给它定价的假设更接近后者。

[... middle omitted ...]

。

科沃斯是被明确回避的标的。报告认为，其面临三重利润率压力——物料成本上涨、竞争阻碍成本传导、以及湿干吸尘器市场份额流失——而市场共识尚未充分反映这一点。JPM的净利润率预测为8.0\%，远低于市场共识的9.6\%。

5. 报告尚未完全回答的关键问题：B2B转型的执行风险与时间窗口

这份报告的逻辑链条是清晰的，但并非没有未解的问题。其中最关键的两个，都指向B2B转型的执行层面。

第一，B2B业务的利润率曲线尚不明确。商用暖通空调和数据中心液冷虽然是高增长领域，但它们的利润率结构、客户集中度、以及资本回报率，与家电业务有本质区别。报告没有充分展开这些业务的单位经济模型，而这是评估“期权价值”能否转化为“实际价值”的核心。

第二，转型的时间窗口。即便B2B业务最终能够带来估值重估，这个过程需要多长时间？市场是否愿意为一个可能需要三到五年才能兑现的期权支付溢价？如果国内需求放缓超预期，现金流的再配置能力是否会受到制约？这些问题的答案，将决定当前估值是否已经足够安全。

\subsection{[R040] JPM：零售投资者在科技股上的踩踏式抛售，暴露的不仅是情绪，更是仓位结构的脆弱性}
{\small\color{refgray}归类：风险偏好：科技股信心结构性重置与亚洲动量泡沫}

JPM：零售投资者在科技股上的踩踏式抛售，暴露的不仅是情绪，更是仓位结构的脆弱性

这份发布于6月10日的JPMRetail Radar周报，捕捉到了一个关键信号：散户资金流在一周之内从极端买入科技股转向了创纪录的卖出。这并非普通的获利了结，而是一次由业绩回撤触发的、非选择性的仓位出清。

对于产业决策者和高净值投资者而言，真正值得关注的不是散户的情绪波动本身，而是这种极端行为背后所反映的市场结构——当最坚定的趋势追随者开始无差别地削减仓位时，此前由AI叙事驱动的拥挤交易正在经历一次真正的压力测试。报告的数据表明，市场正在经历的，可能不是简单的风格切换，而是对过去一年半“买科技、买AI”这一核心叙事的重新定价。

在这份报告中，JPM策略团队不仅记录了资金流向的剧变，还揭示了几个容易被忽视的结构性特征：散户的抛售并非为抄底其他板块筹集资金，而是纯粹的亏损止损；石油相关资产的仓位正在悄然累积；世界杯主题的交易热情尚未被点燃。这些信息为理解当前市场的底层逻辑提供了新的维度。

1. 散户的抛售不是“换仓”，而是“断腕”——这意味着风险偏好的系统性下降

报告中最引人注目的数据点是：在截至6月10日的一周内，零售投资者的美元净流入降至一年多以来的最低水平。ETF和个股的流入百分位分别仅为2.8\%和0.8\%。更关键的是，上周五创下了过去一年中最大的单日散户卖出记录。

这种行为模式与传统意义上的“获利了结”或“板块轮动”有本质区别。报告明确指出，散户卖出活动“主要是由业绩驱动的”——他们的投资组合遭受了回撤，尽管这些回撤尚未显著侵蚀此前积累的超额收益。换句话说，散户是在“砍仓”，而不是在“调仓”。

这意味着什么？当散户投资者开始因为浮亏而被迫卖出，而不是因为看到更好的机会而主动换仓时，市场整体的风险偏好正在经历一次系统性的下降。这种行为的背后逻辑是：投资者对科技股未来走势的信心出现了动摇，他们不再愿意承受短期波动以换取潜在回报，而是选择优先保护本金。这对依赖持续资金流入支撑估值的科技板块而言，是一个需要警惕的信号。

2. 科技ETF流出创历史纪录——拥挤交易的瓦解往往伴随超调

报告用“有史以来最大规模的卖出”来形容散户对科技ETF的抛售，其z-score达到了-3.7，而就在一周前，买入的z-score还高达3.1。这种从极热到极寒的急剧反转，是拥挤交易瓦解的典型特征。

具体来看，被抛售最严重的ETF包括SOXL、SMH和IGV，这些恰恰是过去一年散户最钟爱的半导体和软件ETF。值得注意的是，散户卖出科技ETF的同时，却买入了消费必需品ETF。这看起来像是一次防御性轮动，但报告的数据显示，这种买入规模远不足以抵消科技股的卖出规模。散户的资金整体是净流出的，而不是在板块间流动。

这种“只卖不买”的行为模式，意味着科技板块正在经历一次资金的净撤离，而不是简单的内部结构调整。对于持有科技仓位或关注科技估值的读者而言，这提示了一个风险：当趋势交易者开始集体离场时，价格的下行幅度往往会超出基本面所能解释的范围，因为流动性本身成为了驱动因素。

3. 散户没有“抄底”科技股——NVDA和TSLA的买入掩盖了板块内部的广泛疲软

一个常见的误解是，散户会在科技股下跌时“逢低买入”。但这份报告的数据清楚地表明，情况并非如此。虽然NVDA和TSLA重新回到了散户买入榜的前列，但这只是个别现象。

报告特别指出：“值得注意的是，科技板块没有出现广泛的逢低买入”。卖出名单上充斥着MU、MRVL、AMD、INTC、ARM、DELL等过去备受追捧的名字。其中MRVL的卖出z-score高达-6.9，ARM为-7.0，这些都是极端水平的卖出信号。

这揭示了一个重要的结构性分化：散户对AI龙头股（如NVDA）的信仰尚未完全崩塌，但对于更

[... middle omitted ...]

”。

这意味着，散户正在集体押注油价上涨。结合美国对伊朗发动新一波打击的背景，这种仓位积累有其基本面逻辑。但报告隐含的警示是：当一种交易的拥挤度持续处于高位且无法收敛时，一旦出现预期外的反转，其反向冲击的力度也会相应放大。

对于关注能源板块的读者而言，这里的关键问题不是油价本身会涨多少，而是当前散户在石油相关资产上的仓位是否已经隐含了过于乐观的预期。如果地缘政治风险溢价被充分定价，那么任何关于停火或谈判的进展都可能触发一轮快速的仓位出清。

5. 世界杯主题尚未被定价——这可能是市场情绪改善的潜在催化剂

报告提到，世界杯将于周四开幕，但“散户投资者并未积极交易这一主题”。JPM认为，考虑到宏观背景的不确定性、地缘政治风险以及对低端消费者的担忧，市场对潜在受益者的预期“相当低”。

这正是报告中一个值得关注的“预期差”。当一个明确的催化事件即将到来，而市场对其反应冷淡时，往往意味着该主题尚未被充分定价。JPM策略团队因此建议“战术性做多2026世界杯受益者篮子”。

\subsection{[R041] JEF：泡泡玛特的韧性来自股东结构变化，而非基本面改善}
{\small\color{refgray}归类：未被主题摘要引用}

这份JEF最新发布的IP玩具月度追踪报告，表面上是一份常规的数据更新，但如果我们用金字塔原理来拆解，会发现一个被市场情绪掩盖的核心判断：泡泡玛特股价在1Q26业绩发布后的韧性，更多来自股东结构的结构性变化，而非运营基本面的加速改善。段永平在5月底购入约6\%股份成为重要股东，这一事件对市场信心的支撑力度，可能超过了公司自身业绩所能提供的边际增量。

报告发布时点恰逢市场对泡泡玛特海外增长动能是否可持续产生分歧。4月电商GMV同比增长70\%、环比增长25\%的数据，以及Labubu二手价格环比企稳的信号，确实为多头提供了弹药。但与此同时，报告也坦率指出：公司App活跃用户数在5月出现下滑，全球及区域Google搜索趋势继续走弱，Instagram粉丝增速也在放缓。这些信号组合在一起，指向一个更复杂的图景——短期催化剂（Labubu 4.0发布、世界杯相关营销）能否转化为可持续的增长，仍需验证。

这份报告的价值不在于罗列数据，而在于它揭示了一个关键张力：市场正在用股东结构变化来定价，而非用运营数据的边际变化来定价。这种定价逻辑的切换，对投资者的观察框架提出了新要求。

报告开篇就给出了一个清晰的判断：在5月13日1Q26业绩发布后，泡泡玛特股价保持了韧性，但JEF并未看到基本面有任何实质性变化。真正改变市场情绪的事件，是段永平在5月25日至28日期间购入8100万股，约占公司总股本的6\%，使其成为重要股东。

这一判断的关键在于时间序列的对照。业绩发布到段永平入股之间有两个星期的窗口期，如果市场对业绩本身有强烈正面反应，股价韧性应该在这两周内就已经建立。但报告暗示，真正驱动股价走强的催化剂是后续的股东结构变化。段永平的入场不仅带来了资金层面的支撑，更重要的是向市场传递了一个信号：有长期视角的投资者愿意在这个价位建立头寸。

这意味着投资者需要区分两种不同类型的支撑力。运营数据改善带来的支撑是可持续的，因为它反映了公司竞争力的提升；而股东结构变化带来的支撑是情绪性的，其持续性取决于新股东是否会继续增持、是否会参与公司治理、以及市场是否会将这一事件解读为价值发现的信号。目前来看，这些维度都尚未得到充分验证。

2. 海外社交媒体的互动回暖与流量下滑并存，增长质量需要更细致的拆解

报告在海外运营数据部分呈现了一个看似矛盾的画面。一方面，TikTok上的互动量在5月出现明显回升，全球账号和泰国账号分别环比增长49\%和67\%；另一方面，Instagram粉丝净增量从4月的1.9万降至5月的1.3万，App活跃用户数在全球和美国市场均出现下滑，Google搜索趋势也继续下行。

这种“互动回暖但用户基础收缩”的组合，可能意味着泡泡玛特在海外市场的增长正在从“广度扩张”转向“深度运营”。互动量的提升说明现有粉丝的参与度在提高，这可能得益于世界杯相关营销活动的拉动；但粉丝增速放缓和搜索趋势下行，则暗示新用户的获取成本在上升，或者品牌在目标市场的渗透率已经接近一个阶段性平台。

对于投资者而言，这个信号的含义是：海外增长的下一个阶段，可能不再是由社交媒体病毒式传播驱动的自然增长，而是需要更系统性的渠道建设、本地化运营和产品创新来支撑。世界杯营销和Labubu 4.0能否带来新用户增量，还是仅仅激活了存量用户，将是判断海外增长质量的关键观察点。

报告提供的4月电商数据显示，泡泡玛特在主流电商平台的GMV同比增长70\%、环比增长25\%，增速显著领先于竞争对手。同期，Bloks的电商销售同比增长16\%，名创优品同比增长11\%。而52Toys和卡游的电商销售则出现同比下滑。

这一数据分化的背后，可能反映了三个结构性变化。第一，泡泡玛特凭借IP矩阵的深度和广度，正在将规模优势转化为电商渠道的流量获取效率优势。当消费者在电商平台搜索盲盒或潮

[... middle omitted ...]

4. 二手市场价格企稳但产品结构趋于分散，IP生命周期的管理能力面临考验

报告指出，Labubu毛绒玩具的二手价格在5月环比企稳。同时，在二手交易平台上最受欢迎的30款产品中，出现了更多样化的构成：其中有3款是非泡泡玛特产品，而泡泡玛特与aespa的联名款仍然是需求最旺盛的产品。

二手价格企稳是一个积极信号，它说明市场对Labubu系列的需求没有出现急剧萎缩。但产品结构的分散化趋势值得关注。当非泡泡玛特产品开始进入最受欢迎榜单的前列，意味着消费者的注意力正在被其他IP分散。联名款成为需求顶峰，也说明泡泡玛特自有IP的热度可能已经接近阶段性高点，未来的增长可能需要更多地依赖外部IP合作。

这引出了一个更深层的问题：泡泡玛特的IP生命周期管理能力是否足以支撑其估值溢价？如果Labubu的热度开始自然衰减，公司是否有足够的IP储备来填补空白？报告没有给出明确答案，但它在数据层面提供了一个观察框架：二手价格走势和热门产品构成的变化，将是判断IP生命周期阶段的前瞻性指标。

\subsection{[R042] 摩根斯坦利：市场真正低估的不是需求，而是供给纪律的长期重塑}
{\small\color{refgray}归类：半导体：欧洲设备需求周期尚未见顶与存储供给纪律重塑}

过去两个月，DRAM股票涨幅高达70\%至134\%，随后经历了一轮剧烈回调。市场主流叙事将这轮波动归结为“周期见顶的信号”——历史经验表明，当存储周期进入第六个季度，供给松动、价格增速放缓，股票往往提前四个月见顶。但摩根斯坦利在这份最新研报中提出了一个反直觉的核心判断：这轮回调不是周期的终结，而是周期得以延续的必要重置。

真正值得关注的，不是需求是否持续——AI带来的结构性需求已经是共识——而是供给端正在发生的历史性变化。过去四十年的存储周期，每一次崩溃的根源都不是需求不足，而是供给纪律从未真正建立。这一次，长期供应协议（LTA）正在改变游戏规则。

如果LTA能够覆盖未来三到五年70\%以上的总供给，那么DRAM股票的理论估值中枢将从当前的5倍PE，向8至10倍PE迁移。这一判断的隐含含义是：当前市场对存储公司的定价，仍然停留在“周期股”的旧框架里，而忽视了其向“准成长股”转型的可能性。

1. 这轮回调不是周期的终点，而是市场从“动量驱动”转向“基本面驱动”的必经之路

自2022年秋季生成式AI问世以来，存储板块已经经历了三次类似的回调。每一次都是抛物线式上涨后，杠杆ETF、对冲基金和散户的拥挤持仓与价格走势发生冲突，引发剧烈修正。摩根斯坦利在报告中明确指出，这种价格层面的重置并不意味着周期结束，恰恰相反，它是周期得以延续的必要条件。

关键在于理解这轮回调的性质。它不是基本面恶化引发的恐慌，而是持仓结构过载后的技术性修正。报告引用了韩国市场历史数据：在KOSPI熔断事件后，三星电子的下一个交易日上涨概率为87.5\%，三个月内平均回报率达到36.9\%。虽然历史不会简单重复，但这个统计至少说明：当波动由情绪和仓位而非基本面驱动时，风险回报比往往在下跌中改善。

更重要的信号在于估值本身。经过这轮回调，三星和SK海力士的远期PE分别回到5.2倍和4.7倍。在盈利修正仍然强劲、LTA框架正在重塑收入可见性的背景下，这个估值水平并没有反映供给端正在发生的结构性变化。

2. 供给纪律是存储周期的终极变量，而LTA正在改写这个变量的行为模式

过去每一次存储周期，需求从来不是问题——除了2008年全球金融危机和2021年疫情初期的需求冲击，每一轮周期都有真实且增长的需求。但周期总是以供给过剩、价格崩溃告终，原因只有一个：供给纪律从未真正建立。当价格上涨，厂商竞相扩产，最终供过于求，周期见顶。

摩根斯坦利认为，这一轮的不同之处在于LTA的普及。LTA的本质是下游客户（云厂商、GPU制造商）与存储厂商之间的长期锁量锁价协议。报告估算，如果LTA能够覆盖未来三到五年70\%以上的总供给，那么存储厂商的盈利可见性将发生质变——它们不再完全暴露于现货市场的价格波动。

这个判断的底层逻辑是：AI需求是基础设施级别的，而非消费电子级别的。GPU的每一代迭代都受限于内存带宽而非算力——从A100的80GB到Rubin GPU/Superchip的288至768GB，单AI单元的DRAM容量增长了4到7倍。当需求与长期基础设施投资绑定，而非与消费者换机周期绑定，存储厂商对未来需求的可见性大幅提升，这反过来抑制了盲目扩产的冲动。

报告引用管理层指引指出，新增产能将以更加可控的节奏释放，而非像过去那样一拥而上。如果这个判断成立，那么当前周期的“峰值”可能不是一个尖顶，而是一个更宽的平台——盈利在高位持续的时间可能远超市场预期。

3. 估值重构的逻辑：当70\%的盈利被LTA锁定，PE应该向市场均值靠拢

摩根斯坦利提出了一个简洁但有力的估值框架。如果LTA覆盖了2027年盈利的70\%以上，那么这部分盈利应该按照市场平均估值（滚动五年10至14倍PE）定价，而剩余30\%的非LTA盈利仍然按照历史周期峰值5倍PE定价。简单加权计算，DRA

[... middle omitted ...]

盈利修正，而盈利修正反过来支撑PE的稳定性。在5.2倍2027年预期市盈率下，即便股价已经大幅上涨，估值仍未反映LTA带来的结构性变化。

4. 报告尚未完全回答的关键问题：价格弹性如何改变AI需求的二阶效应

摩根斯坦利提出了一个值得深入探讨的命题：AI需求的价格弹性可能与消费电子需求完全不同。更低的DRAM价格会降低AI推理的运行成本，使AI部署更加便宜，从而创造新的需求、扩大可寻址市场，而不仅仅是降低固定消费电子设备的成本。

这个命题的推论是：即便2027年末新产能上线导致DRAM价格下跌，也不一定意味着周期终结。价格下跌可能刺激出新的AI应用场景，形成“价格下降-需求扩张-产能消化”的正反馈循环。

但报告没有完全回答的问题是：这个价格弹性到底有多大？在什么价格水平下，新的AI需求才会大规模涌现？如果价格下跌30\%，AI推理成本下降能否带来50\%以上的需求增量？这个问题至关重要，因为它决定了2027年后供给释放时，行业是否会重蹈“供给过剩-价格崩溃”的覆辙。

\subsection{[R043] 摩根斯坦利：中国电动车市场正进入“产品攻势间歇期”的定价博弈}
{\small\color{refgray}归类：中国新能源车：渗透率突破60\%后的竞争结构重构}

摩根斯坦利：中国电动车市场正进入“产品攻势间歇期”的定价博弈

中国电动车市场的周度订单数据，在经历了密集的新品投放季后，正进入一个微妙的“中场休息”阶段。摩根斯坦利最新发布的周度订单追踪报告显示，6月第一周（6月1日至7日），主要电动车厂商的周度订单普遍出现环比下滑。这份报告表面上看是一张常规的周度数据快照，但它揭示了一个更深层的结构性信号：市场对“产品周期”的线性外推正在失效，真正决定下一阶段竞争格局的，不是谁推出了新车，而是谁能在新品间歇期维持订单势能、并在下一轮密集投放前完成渠道与产能的校准。

对于产业决策者和投资者而言，这份报告的价值不在于单周订单的涨跌，而在于它提供了一个观察竞争格局演变的时间窗口。当市场注意力被比亚迪的“大唐”预售、华为与赛力斯的“尚界”系列、以及理想L8的即将推出所吸引时，摩根斯坦利的这份周报提醒我们：在电动车行业，产品攻势的节奏本身已经成为一种竞争变量。

6月第一周的订单数据，最显著的特征是“正常化”。蔚来汽车周订单从5月底的峰值骤降72\%，小鹏汽车下滑49\%，问界下滑69\%。这些数字如果孤立来看，会引发对需求疲软的担忧。但摩根斯坦利在报告中明确指出，这些下滑是“新品上市冲高后的正常回落”——蔚来的ES9、小鹏的GX、问界的M9，都在此前一周或数周内经历了订单爆发，而6月首周的数据只是回到了更可持续的水平。

这里的核心洞察不在于数字本身，而在于数字背后的预期管理。对于一家车企而言，新品上市当周的订单量往往带有“粉丝效应”和“早期尝鲜者”的溢价，这些订单的转化率、退单率与正常订单有本质区别。摩根斯坦利在报告中暗示，真正的考验是这些订单能否在6月交付中得到兑现。报告提到“强劲的订单积累应支持6月交付量回升”，这实际上是在提醒市场：不要因为单周订单下滑而过度悲观，交付数据才是更可靠的验证指标。

但反过来看，这种“新品效应”的衰减速度本身也是一个竞争信号。蔚来和小鹏的订单在上市后仅一周就出现如此剧烈的回落，说明消费者对电动车的购买决策正在变得更加理性、更加价格敏感。过去那种“新车上市即爆款”的确定性正在降低，取而代之的是“上市后需持续投入营销和渠道才能维持势能”的新常态。这意味着，车企的营销费用率和渠道效率，正在成为比产品本身更重要的竞争维度。

在所有主要厂商中，比亚迪的表现最为稳健。6月第一周订单 万辆，环比仅下降10\%，同比虽然下降16\%，但考虑到去年同期的基数效应，这个数字并不令人担忧。更重要的是，报告指出比亚迪后续订单的增长将取决于“大唐”预售订单的转化以及“大汉”的催化作用。

比亚迪的订单韧性，本质上反映的是其规模效应带来的两个竞争优势。第一是价格锚定能力。当其他品牌的新品上市后订单快速回落时，比亚迪凭借其从低端到中高端的完整产品矩阵，能够在不依赖单一爆款的情况下维持稳定的订单流。第二是渠道消化能力。比亚迪的经销商网络和产能布局，使其能够在新品上市后迅速将订单转化为交付，从而缩短订单到收入的周期。

但摩根斯坦利的报告也隐含了一个尚未回答的问题：比亚迪的同比下滑是否意味着其市场份额正在被侵蚀？报告显示比亚迪周订单同比下降16\%，而同期蔚来同比增长144\%、零跑增长66\%、小米增长60\%。这些高增长数字固然有基数效应，但也确实表明，在20-30万元价格带，比亚迪正面临来自新势力的更激烈竞争。比亚迪能否通过“大唐”和“大汉”的新品重新夺回这一价格带的话语权，将是下半年竞争格局的关键变量。

3. 华为系（HIMA）的订单波动暴露了“品牌叠加”的边际效应递减

华为系的订单表现是这份报告中最值得拆解的部分。HIMA整体周订单 万辆，环比下降58\%，但同比仍增长30\%。其中，问界周订单0.88-0.9万辆，环比骤降69\%，同比仅增长9\%。报告将这一下滑归因于M9上市后的正

[... middle omitted ...]

据上形成接力，那么“华为赋能”这一标签的边际价值可能正在下降。

理想汽车周订单 万辆，环比下降9\%，同比下滑20\%。小米汽车周订单0.8-0.82万辆，环比下降20\%，同比增长60\%。这两家公司的数据看似平淡，但它们的竞争逻辑与比亚迪和华为系截然不同。

理想的订单下滑幅度最小，这并非偶然。理想的产品策略一直以“稳健”著称——不追求单月订单爆发，而是通过L系列的持续迭代和用户口碑积累，维持稳定的订单流。报告提到L8即将在未来几周推出，这意味着理想的下一轮产品攻势正在酝酿中。理想的挑战在于，如何在L8上市后，避免与L6、L7形成内部竞争，同时保持毛利率的稳定。

小米的订单表现则反映了其“慢热型”的产品节奏。60\%的同比增长固然亮眼，但环比下降20\%也说明，小米SU7的订单积累正在趋于平稳。小米的优势在于其庞大的用户生态和营销能力，但劣势在于汽车业务仍处于产能爬坡期，订单转化率可能低于成熟厂商。摩根斯坦利的报告没有直接讨论小米的产能瓶颈，但这是一个值得关注的隐含变量。

\subsection{[R044] 摩根斯坦利：全球钢铁市场的真正拐点不在需求，而在欧洲供给侧的重新定价}
{\small\color{refgray}归类：能源与大宗：欧洲钢铁定价权重塑与储能政策重心转移}

摩根斯坦利：全球钢铁市场的真正拐点不在需求，而在欧洲供给侧的重新定价

全球钢铁市场正在经历一个被多数投资者低估的结构性转变。这份摩根斯坦利刚刚发布的研报，核心判断并非关于中国出口量的短期波动，而是指向一个更深层的趋势：欧洲正在通过政策工具和地缘政治现实，系统性地重塑其钢铁定价权。对于关注全球资产配置和产业趋势的决策者而言，理解这一转变，远比跟踪月度贸易数据更为关键。

报告揭示了一个看似矛盾的现象：中国5月成品钢净出口环比增长约10\%，达到989万吨，年化运行率约1.19亿吨，甚至高于摩根斯坦利中国材料分析师对2026年全年约1.05亿吨的预测。然而，出口的绝对数量并非重点。真正值得关注的是，这些增量正在被全球贸易格局的区域化趋势所抵消，而欧洲正成为这一趋势的最大受益者。

中国出口数据本身并不构成对欧洲市场的直接威胁。报告明确指出，中国仅占欧洲钢铁进口量的约13\%。但这组数字背后的结构性含义，才是投资者需要深挖的。当全球钢铁市场从“全球化套利”转向“区域化定价”，欧洲本土钢厂正在获得过去十年未曾有过的定价权。这并非一个短期现象，而是一个由政策、地缘和成本结构共同塑造的新均衡。

1. 欧洲正在通过政策工具建立一道隐形的结构性壁垒，而非仅仅针对中国的贸易摩擦

市场习惯将欧洲的钢铁政策视为针对中国过剩产能的临时性防御。但摩根斯坦利的分析表明，这套政策框架的设计逻辑远比“反倾销”复杂。从2026年1月1日起实施的CBAM（碳边境调节机制），以及从7月1日起收紧的保障措施（包括更高的配额外关税），正在共同构建一个多层级的保护体系。

报告特别提及了“熔化与浇注条款”（melted-and-poured clause）。这一条款的设计精妙之处在于，它不仅仅限制中国直接出口，还通过追溯钢材的原始冶炼地点，堵住了经由第三国（如越南、土耳其）转口的路径。这意味着，即使中国对欧洲的直接出口量有限，政策仍能通过截断中间环节来影响全球贸易流。

这一判断的核心含义是：欧洲不再依赖需求复苏来支撑钢价。即使终端需求持续疲软，政策本身就能创造供给缺口。报告估算，即使不考虑需求回升，现有政策框架也可能使欧洲结构性短缺 万吨。这相当于欧洲年消费量的约10\%。供给短缺意味着更高的出清价格，而欧洲热轧卷板（HRC）的吨钢毛利已升至404美元/吨，显著高于320美元/吨的长期均值。这不是一个周期性高点，而是一个结构性重定价的开始。

2. 地缘政治正在加速供应链的“近岸化”溢价，而不仅仅是成本比较

许多分析将欧洲钢价上涨归因于能源成本或碳成本。但这份研报提供了一个被低估的变量：供应链安全溢价。自中东冲突爆发以来，安特卫普CIF热轧卷板价格上涨了70美元/吨，而同期中国FOB价格仅上涨29美元/吨。这意味着欧洲的溢价在冲突期间额外扩大了41美元/吨。

这41美元/吨的溢价，本质上是市场对“可靠性”和“邻近性”的定价。当运费、保险和供应链中断风险上升时，买家愿意为更短的交货周期和更稳定的供应支付溢价。这种溢价在和平时期可能收缩，但地缘政治的不确定性已成为常态，而非例外。

对于投资者而言，这意味着欧洲钢厂的竞争优势不再仅仅来自成本端，而是来自其地理位置和市场准入。拥有本地一体化资产、能够灵活调整发货量、并受益于区域溢价扩大的企业，将获得持续的利润重估。报告明确将安赛乐米塔尔（ArcelorMittal）列为首选标的，原因正是其能够通过调整产量来改善固定成本吸收，并在政策框架将边际吨位重新导向国内供应时抢占市场份额。

3. 中国出口的结构性下降正在发生，但影响路径比表面数据更复杂

中国5月出口的环比增长，不应被误解为中国钢铁出口的再次扩张。报告数据显示，年初至今（YTD）中国净出口仍同比下降约8\%。与此同时，中国钢协（CISA）成员产量同比

[... middle omitted ...]

冲击，但也意味着全球贸易摩擦的“涟漪效应”将更加复杂。例如，当欧盟收紧政策时，中国出口可能转向东南亚，进而压低当地价格，再间接影响从东南亚向欧洲的转口贸易。

对于投资者而言，这意味着不能简单用“中国出口量”来预测全球钢价。真正的变量在于中国出口的“边际成本”和“最终目的地”。报告没有完全展开的一个关键问题是：当中国国内需求进一步恶化时，出口是否会突破当前的政策和成本约束，重新成为全球价格的下行压力来源？这个问题的答案，取决于中国钢厂能否在亏损状态下维持出口规模，以及全球其他市场能否吸收这些增量。

4. 报告尚未完全回答的关键问题：政策执行的“二阶效应”与需求端的脆弱性

摩根斯坦利的分析框架在供给端和政策端提供了清晰的逻辑，但需求端的假设仍存在不确定性。报告的核心论点——政策足以创造结构性短缺——建立在“即使没有需求复苏”的前提上。但这是一个强假设。如果欧洲经济出现超预期衰退，导致钢铁消费大幅下滑，那么政策创造的供给缺口可能被需求萎缩所抵消，定价权逻辑将面临挑战。

\subsection{[R045] NOM：中国汽车市场2026年可能首次出现国内需求双位数下滑}
{\small\color{refgray}归类：中国新能源车：渗透率突破60\%后的竞争结构重构}

NOM：中国汽车市场2026年可能首次出现国内需求双位数下滑

这份NOM研报的核心判断并不令人意外，但其警示的深度值得认真对待。5月数据公布后，市场普遍关注的是环比改善，而NOM选择将目光投向一个更根本的问题：国内汽车市场正在经历一个结构性转折点，2026年全年国内零售销量可能出现有史以来首次双位数同比下滑。这个判断如果成立，将从根本上改变对中国汽车产业链的定价逻辑。

报告发布时点恰逢国内车企密集推出新车型和技术升级的窗口期，但NOM的订单追踪显示，除了新车上市带来的短期脉冲，市场并未出现有意义的有机需求改善。这意味着，当前行业面临的问题不是供给端缺乏创新，而是需求端正在经历一个更深层次的收缩周期。对于投资者和产业决策者而言，理解这个收缩的性质和边界，远比关注月度环比数据更有价值。

1. 5月数据确认了国内需求的疲软并非短期扰动，而是趋势性下滑

5月中国乘用车零售销量（剔除微型客车）为150万辆，同比下降22.0\%。虽然环比增长9.2\%看起来是一个积极信号，但NOM明确指出，这种环比改善更多是季节性因素所致，而非需求回暖的标志。更值得关注的是，今年前5个月累计零售销量为710万辆，同比下降19.3\%，这是自2020年疫情冲击以来的最低水平。

将这个数据放在更长的历史背景下看，2020年国内零售全年约为1920万辆，而2026年前5个月的运行节奏如果延续，全年零售可能落在 万辆区间。这意味着国内市场规模正在退回五年前的水平。NOM在报告中提出一个大胆但逻辑自洽的假设：如果当前趋势延续且没有增量刺激政策，2026年国内汽车销量可能出现历史上首次双位数同比下滑。

这个判断的意义不在于预测的精确度，而在于它揭示了一个关键变化：过去几年市场一直习惯于将销量下滑归因于短期因素（如补贴退坡、库存调整、疫情扰动），但2026年的情况可能不同。即便在大量新车投放和技术升级的背景下，需求依然没有起色，这说明问题可能出在消费者购买力或消费意愿的结构性变化上。

5月新能源乘用车零售销量为95万辆，同比下降7.5\%，环比增长12.0\%。渗透率达到62.1\%，同比提升9.9个百分点，环比提升1.5个百分点，创下历史新高。与此同时，燃油车零售销量为56.1万辆，同比下降38.4\%。

渗透率的快速提升被市场广泛解读为利好，但NOM的数据揭示了一个更复杂的图景：新能源车销量本身也在下降，只是降幅远小于燃油车。这意味着渗透率的提升更多来自燃油车市场的加速萎缩，而非新能源市场本身的强劲扩张。这是一个重要的区别——如果新能源市场能够维持正增长，渗透率提升是健康的；但如果新能源市场也开始收缩，渗透率的继续提升可能只是分母效应。

从车型结构看，A00级微型电动车市场份额在4月已降至6\%，NOM预计2026年全年这一比例将继续维持在个位数。国家补贴政策对大众市场车型越来越不利，这意味着新能源市场的增长重心正在向中高端转移。对于主打低价走量策略的车企而言，这个趋势意味着它们需要重新审视自己的产品定位和盈利模型。

NOM对主要车企的零售数据拆解，提供了一个观察竞争格局变化的窗口。比亚迪5月零售20.74万辆，同比下降29.2\%，市场份额21.8\%，同比下降6.7个百分点。虽然环比有所恢复，但同比的大幅下滑说明，即便是行业龙头，在整体需求萎缩的环境下也难以独善其身。

吉利5月零售10.92万辆，同比下降16.3\%，市场份额11.5\%，同比下降1.2个百分点。零跑汽车是少数实现同比增长的车企之一，5月零售6.14万辆，同比增长48.3\%，显示其在特定细分市场的产品定位获得了消费者认可。

值得注意的是，特斯拉中国5月零售4.73万辆，同比增长22.5\%，环比大幅增长82.2\%。这个环比跳跃可能与新车型交付节奏有关，但NOM并未展开分析。蔚来5

[... middle omitted ...]

对可持续性提出了关键疑问

5月中国乘用车出口80.9万辆，同比增长73.0\%，前5个月累计出口350万辆，同比增长70\%。这个增速在任何一个行业都是惊人的。NOM指出，中国车企在欧盟、巴西、澳大利亚等关键市场取得了实质性进展。

然而，NOM在这份报告中提出了两个需要认真思考的问题。第一是目标市场的容量问题。随着现有目标市场逐渐饱和，车企需要不断开拓新市场，而这会带来更多的不确定性和成本。第二是贸易摩擦风险。报告引用了彭博关于中欧贸易关系面临进一步压力的报道，并指出俄罗斯、巴西、墨西哥等多个国家已经提高了汽车相关关税。

NOM的判断是：中国电动车企在2026年至少仍将在海外市场享受强劲增长，但成为真正的全球玩家并不容易。这个判断的潜台词是，出口高速增长的可持续性需要打一个问号。如果贸易壁垒继续升级，出口增速可能会在2027年前后出现拐点。对于高度依赖出口来消化国内过剩产能的车企而言，这个风险不容忽视。

5. 报告没有完全回答的关键问题：内需收缩的根源到底是什么？

\subsection{[R046] JPM：日本股市真正的博弈不在AI，而在价值与动量的交叉点}
{\small\color{refgray}归类：风险偏好：科技股信心结构性重置与亚洲动量泡沫}

JPM最新发布的日本权益风格投资报告，在看似延续“AI半导体引领市场”的共识背后，藏着一个更值得推敲的判断：当前日本股市最被低估的交易机会，并非追高动量股，而是那些同时具备价值特征和适度动量的股票。这个判断的逻辑链条并不复杂，但它的含义却可能被大多数投资者忽视。

这份报告发布的时间点值得注意。日经225指数在6月初逼近68000点，全球AI半导体热潮持续推高相关标的，动量因子在过去12个月录得0.44的夏普比率，Beta因子表现同样强劲。市场看起来气势如虹。但JPM的量化团队却在此时刻意提醒两个风险：市场对日本央行可能更为鹰派的加息路径定价不足，以及赢家与输家之间极端分化可能出现的部分修正。

换句话说，这份报告的核心主张是：在动量过热、价值被压抑的当下，投资者不应简单地在两个阵营之间做二选一，而应该寻找那个被忽略的交集地带。

JPM在报告中提供了一个关键数据点：动量因子在5月录得6.4\%的月度回报，4月更是高达14.1\%。但3月的数据是-12.9\%。这种剧烈波动本身就说明，动量交易已经不再是低波动率的“躺赢”策略。当动量因子在三个月内从-12.9\%跳到14.1\%再回到6.4\%，它本质上已经变成了一个高Beta、高换手率的博弈工具。

更值得关注的是持仓集中度。报告明确指出，动量因子的持仓已经高度集中在AI相关标的上。这不是一个分散化的动量组合，而是一个主题化的动量赌注。当市场出现任何对AI主题不利的信号——无论是地缘政治风险、利率预期变化，还是行业自身的周期性回调——这个拥挤的动量仓位都可能面临急剧的去杠杆压力。

JPM也承认，在动量出现“迷你崩盘”时，基本面投资者往往会逢低买入，这在一定程度上限制了下行空间。但问题在于，这种“基本面兜底”的逻辑依赖于一个关键假设：AI半导体的盈利增长预期不会出现系统性下修。如果宏观环境发生变化，这个假设就可能被动摇。

所以，动量交易的风险回报比已经不再像几个月前那样有吸引力。这不是说AI主题会消失，而是说在这个价位上，继续押注动量所承担的风险，可能已经超过了潜在的回报。

价值因子在5月的表现可以用“惨淡”来形容。根据报告中的数据，Price to Book因子5月仅上涨0.8\%，Earnings Yield下跌2.3\%，Div Yield下跌5.9\%，ROE下跌8.9\%，JPM Quality更是下跌10.3\%。与此同时，动量因子上涨6.4\%，Size因子上涨6.3\%，Beta因子上涨6.2\%。

这种极端的风格分化，很容易让人得出一个结论：价值投资在日本市场已经失效。但JPM提供了一个截然不同的视角。报告中的关键图表（Figure 3）展示了价值因子表现与JPM编制的日本央行讲话鹰鸽指数之间的关系。数据显示，当前价值因子的低迷表现，与日本央行鹰派言论的升温高度相关。

这意味着，价值因子的跑输不是因为它所代表的公司基本面出了问题，而是因为市场正在为日本央行加息定价。当利率上升预期增强时，那些对利率敏感的价值股——银行、保险、能源、公用事业——自然会承压。但这恰恰是问题的关键：如果市场对加息的定价已经过度，或者加息落地后利空出尽，那么当前被压制的价值股就可能迎来修复行情。

这里需要指出的是，JPM并没有断言加息预期已经见顶。相反，报告提到“一些市场参与者担心日本央行可能落后于曲线”，这意味着加息风险可能尚未完全释放。但这也正是价值股可能存在的机会窗口——如果市场最终发现加息幅度低于预期，或者企业盈利能够消化利率上升的影响，那么当前的价值股价格就提供了安全边际。

JPM报告中最具操作性的洞察，并非对单一风格的推荐，而是对“Value with Momentum”这个交叉策略的系统性分析。报告专门用一张表格（Figure 4）列出了30只同时具备价值和动量特征的日

[... middle omitted ...]

这意味着它们还没有被动量资金充分挖掘。

JPM特别指出，在AI动量股出现短期轮动时，“Value with Momentum”特征的股票可能成为资金的避风港。这个判断隐含了一个重要假设：资金不会离开日本股市，只是从拥挤的动量交易转向相对低估的标的。如果这个假设成立，那么这种轮动对日本股市整体而言是健康的，但对那些押注纯动量策略的投资者来说，则意味着需要重新调整仓位。

报告中的日本QMI（量化宏观指标）在5月继续处于扩张区间，这是支撑当前风险偏好和动量交易的重要宏观背景。但JPM在报告中埋下了一个关键的伏笔：QMI改善的速度在一定程度上反映了“一次性利好”，比如内存领域的特殊需求。从统计上看，维持当前改善速度“具有挑战性”。

这个表述的潜台词是：当前的扩张可能已经接近周期顶部。报告中的Figure 5和Figure 6展示了QMI的历史走势和状态位置，虽然数据截至5月仍处于扩张区间，但JPM明确提醒投资者“应关注潜在的周期性减速，并对Beta效应的过热保持警惕”。

\subsection{[R047] Bernstein：印度350cc摩托车市场真正的较量，不是产品，是分销}
{\small\color{refgray}归类：区域市场：印度制造业竞争力缺口与韩元贬值悖论}

Bernstein：印度350cc摩托车市场真正的较量，不是产品，是分销

印度摩托车市场正在经历一场罕见的正面交锋。Bajaj Auto与Triumph的合资品牌，在2026年4月完成了一次关键的产品线迁移——将旗下主力车型从398cc调整至349cc，正式进入Royal Enfield统治了数十年的350cc排量区间。表面上看，这是一场关于产品力的竞争：马力、保修、服务间隔。但Bernstein这份研报通过深入的经销商调研，给出了一个更克制的判断：市场尚未真正测试Royal Enfield的护城河，因为Triumph还没有让足够多的消费者做出“二选一”的决定。

这份报告的核心洞察在于：Triumph的挑战不是产品够不够好，而是它能否在Royal Enfield的“文化领地”上，建立起足以改变消费者决策路径的分销网络。这不是一个关于发动机排量的故事，而是一个关于“消费者从哪里获取信息、在哪里做决定、信任谁的服务网络”的故事。

以下是我们从Bernstein研报中提炼出的五个关键层次，它们共同指向一个尚未终结的竞争格局。

2025年8月，印度政府对摩托车消费税进行了结构性调整：350cc及以下排量的税率从28\%降至18\%，而350cc以上排量的税率则从28-31\%提升至40\%。这意味着，一个22个百分点的税收悬崖，突然出现在Bajaj和Triumph的398cc产品线面前。

Bernstein的分析指出，这一变化被市场普遍解读为“Bajaj获得了进入350cc市场的战略机遇”，但更精确的表述是：它创造了一个无法回避的“义务”。在税改之前，Bajaj选择在398cc排量上竞争，是一个自由的商业决策；税改之后，维持这一排量意味着承受结构性价格劣势，而这不是任何品牌能够长期承担的。

工程上的解决方案并不复杂——将398cc单缸发动机的缸径缩小至349.13cc，保留冲程，性能几乎不受影响。2026年4月，Triumph的4-5款车型以349cc排量重新上市。Bajaj旗下的Pulsar NS400Z、Dominar 400和KTM 390 Duke也随后跟进。

但关键问题在于：价格差距并没有因为排量下调而显著缩小。Bernstein的对比数据显示，Speed 400与Classic 350的价格差仅从约3.4万卢比缩小至2.5万卢比；Speed T4与Hunter 350的价差从3.5万卢比收窄至3万卢比。换句话说，Triumph进入了相同的排量类别，但并未进入相同的价格区间。

这意味着什么？Triumph现在确实“在对话中”——消费者在考虑Royal Enfield时，会开始提及Triumph的名字。但“被提及”是获胜的必要条件，远非充分条件。

Bernstein在孟买进行的经销商调研，是这份报告最具有现场感的章节。它呈现了两幅截然不同的画面。

在Royal Enfield的展厅，销售人员首先提到的不是产品，而是“等待期”。标准车型：3个月。定制化订单：45-60天，还需支付 卢比的溢价。销售人员提及等待期时，语气中没有任何歉意——这是一种“凭证”，证明品牌处于满负荷运转状态。展厅的设计语言是“目的地”而非“卖场”：暗色墙壁、琥珀色灯光、每款Classic 350之间有足够的间距。每一个表面都在传达同一个信息：这不是一笔交易，而是一场入会仪式。

在Triumph的展厅，黑色的内饰和琥珀色的灯光几乎如出一辙。但销售人员明显接受过更好的培训——对规格参数更熟悉，对比数据更流畅，应对Royal Enfield问题的技巧更娴熟。他的核心话术是：Triumph的产品在性能和工程上太强了，不应该直接与Royal Enfield的350cc车型比较。他提到Hunter的潜在客户也在询问Speed T4和Spe

[... middle omitted ...]

：品牌忠诚度高，看重“突突”声、二手残值、社群认同。即将推出的Bonneville 350是第一个真正的测试，但价格溢价（2.5-4.5万卢比）和有限的分销网络，限制了实质性的影响。威胁：低至中等。

Meteor 350买家（占12\%）：周末旅行者，注重巡航舒适性和长距离骑行。没有Triumph车型直接对应其使用场景。威胁：低。

Hunter 350买家（占23\%）：首次购买高端摩托车的年轻人，年龄18-28岁，看重品牌、轻量化、现代感和价格。Speed T4与之直接竞争，价差已经缩小到0-1万卢比。威胁：高。

Hunter的23\%销量占比，可能低估了它的战略重要性。它是Royal Enfield最年轻的车型，也是增长最快的车型之一，旨在将25岁以下的买家引入品牌。目前，Royal Enfield约30\%的客户年龄在25岁以下，而五年前这一比例约为18-20\%。这个人口结构的演进，是Royal Enfield最重要的长期资产——而它恰恰是Speed T4的靶心。

\subsection{[R048] Bernstein：Visa与OpenAI的合作标志着卡片支付正式进入代理时代}
{\small\color{refgray}归类：新兴科技：AI音乐变现与量子计算速度精度取舍}

Bernstein：Visa与OpenAI的合作标志着卡片支付正式进入代理时代

市场一直在讨论AI代理将如何颠覆支付行业，但大多数人把方向搞反了。他们担心Visa和Mastercard会被新技术边缘化，而Bernstein这份最新研报给出了一个截然不同的判断：卡片网络非但不会被替代，反而正在成为代理经济最核心的信任基础设施。

这不是一个遥远的假设。就在今天，Visa宣布与OpenAI建立合作伙伴关系，共同构建代理商业的基础设施。几乎同一时间，Mastercard也推出了面向机器支付的Agent Pay for Machines服务。两件事发生在同一天，绝非巧合。

如果你还停留在“AI代理需要新的支付方式”这个思维框架里，这份报告值得你重新审视整个逻辑链条。

1. 代理商业真正考验的不是技术，而是信任——这正是卡片网络的护城河

很多人把支付理解成“把钱从A挪到B”，所以他们认为区块链、稳定币或者某种新协议就能解决代理支付问题。但Bernstein的观点更深刻：代理商业的核心矛盾不是资金转移的效率，而是信任的建立与验证。

当AI代理代替人类去完成购物、订阅、订票甚至B2B采购时，谁来证明这个代理是经过授权的？谁来确保交易不会被篡改？谁来处理代理发起的欺诈行为？这些问题没有简单的技术答案。

Visa和Mastercard本质上不是支付处理公司，它们是信任网络。7亿以上的凭证、超过1亿的商户覆盖、几十年来积累的风险评估和欺诈检测能力——这些不是一夜之间可以复制的。OpenAI在合作声明中说得非常清楚：“通过与Visa智能商业的整合，我们正在为安全、透明、用户可控的代理交易构建基础设施。”

这不是一个技术集成，这是OpenAI在为代理经济选择信任层。

Bernstein报告指出，代理商业不仅不会侵蚀Visa和Mastercard的现有收入，反而会打开一个全新的增值服务市场。这一点值得深入拆解。

Visa今天同时发布了三项新服务：Agent Score（评估商户网站对代理商业的友好度）、Agentic Directory（验证商户和代理的身份）、Large Transaction Model（针对大额交易的欺诈检测）。这些服务本质上是在为代理时代建立一套新的认证和风险管理标准。

请注意这里面的商业逻辑。传统支付中，Visa的主要收入来自交易手续费。但在代理商业中，每一笔交易可能都伴随多个增值服务：身份验证、风险评估、代理授权、交易上下文分析。这意味着每笔交易的服务费可能更高，而不是更低。

Bernstein也提到，Visa正在增强其token的能力，加入更多数据维度——交易类型、使用场景、支付方身份、token历史信用信号。这些改进不仅提升了安全性，更创造了新的货币化机会。这个逻辑链条目前尚未被市场充分定价。

报告还提到了一个被很多人忽略的变量：监管环境的变化。Bernstein明确指出，几个关键的不确定性正在消退：CCCA法案的可能性已经大幅降低，MDL和解以及伊利诺伊州的卡片费用诉讼也在推进中。

对于Visa和Mastercard这样长期面临反垄断和费率监管压力的公司来说，监管风险的边际改善往往比业绩增长更能推动估值修复。当市场开始意识到“最坏的情况可能不会发生”，估值折价就会快速收窄。

更关键的是，Bernstein认为卡片网络是优秀的通胀对冲工具。在通胀环境下，名义消费支出上升，Visa和Mastercard按bps收费的模式意味着收入会自然增长。这个特性在当前宏观环境下尤为重要。

Bernstein在研报中提到了一个极具想象力的方向：B2B代理支付。报告写道：“代理商业在B2B领域拥有最强的产品市场契合度。”但坦率地说，这个判断目前还缺乏足够的数据支撑。

B2B支付本身就是一个比

[... middle omitted ...]

. 投资者的观察框架：从“AI受害者”到“AI赢家”的叙事转换

Bernstein提出了一个非常有意思的观察角度：市场对Visa和Mastercard的定价，可能正在从“AI贸易的融资来源”转向“AI的赢家”。这个叙事转换一旦确立，估值逻辑就会发生根本变化。

过去一年，市场把Visa和Mastercard视为“AI基建的买单者”——投资者卖出这些股票，买入英伟达和微软。但如果代理商业真的需要卡片网络作为信任层，那么Visa和Mastercard就变成了AI经济的直接受益者。

Bernstein给出了一个5年的时间框架。这不是一个短期催化剂驱动的故事，而是结构性变化的早期信号。对于长线投资者来说，当前的市场定价可能还没有充分反映代理商业对卡片网络的增量价值。

当然，这个判断需要持续验证。代理商业的渗透速度、商户和消费者的接受度、技术落地的实际效果，都会影响最终的结果。但至少，Bernstein这份研报提供了一个清晰的逻辑框架，而不是简单的“买入”或“卖出”建议。

\subsection{[R049] GS：游戏行业真正的拐点不在2026年的游戏阵容，而在成本结构与IP运营能力的分化}
{\small\color{refgray}归类：未被主题摘要引用}

GS：游戏行业真正的拐点不在2026年的游戏阵容，而在成本结构与IP运营能力的分化

6月初的一周内，State of Play、Summer Game Fest、Xbox Games Showcase、Nintendo Direct四场发布会密集召开，2026年及之后的游戏管线几乎全部揭晓。市场目光自然落在《GTA VI》定档11月、任天堂Switch 2的首年阵容、以及各大发行商的新作排期上。但在这些热闹背后，GS这份研报给出了一个更冷静的判断：真正决定公司未来两年股价走势的，不是管线数量，而是每家公司在“成本上升+硬件周期后半段”这一组合压力下的应对能力。

这份报告覆盖了五家日本游戏公司，给出了两个Buy（索尼、Capcom）、两个Sell（Square Enix、万代南梦宫）以及一个Buy但存在短期不确定性的评级（任天堂）。评级的分化并非基于管线强弱——事实上，Square Enix的管线进展甚至超出预期——而是基于估值、成本结构、以及IP变现效率的深层差异。

1. 索尼是平台商中唯一同时具备“管线强度”与“成本管控话语权”的公司

在四家主要平台/发行商中，GS对索尼的评级最为坚定。理由可以拆解为两层。

第一层是管线。索尼的PS5在2026年下半年拥有《GTA VI》这一行业级别的“核弹”，同时第一方与第三方的软件阵容在研报中被评价为“strong”。从Exhibit 1的排期来看，PS5独占或同步发行的重磅作品包括《Marvel‘s Wolverine》（9月）、《ONIMUSHA: Way of the Sword》（9月）以及多款第三方大作。在2026年秋季这个“神仙打架”的窗口，索尼的独占内容密度是足够的。

但真正让GS给出Buy评级的原因在第二层：成本管控。研报明确指出，在内存成本为首的硬件成本上升背景下，索尼已“明确表达了在业务运营中平衡增长与盈利改善的政策”。更重要的是，PS5处于生命周期后半段，索尼可以通过调整出货量与售价来管理硬件盈利。这一点在与任天堂的对比中尤为关键——任天堂Switch 2刚进入生命周期，硬件成本压力尚未完全消化，且近期已宣布在日本、美国、欧洲涨价。索尼在成本端拥有更大的操作空间。

这里的含义是：在通胀与供应链成本上升的行业背景下，平台商的“成本转嫁能力”正在成为比“独占内容数量”更重要的估值变量。

2. 任天堂的短期股价压力来自第一方管线不足，但硬件渗透率逻辑尚未被破坏

任天堂在本次Nintendo Direct上并未带来让市场兴奋的消息。2026年唯一确定的重大第一方作品是《塞尔达传说：时之笛》重制版，这与市场预期存在落差。研报用“somewhat underwhelming”来描述这一局面。

但GS并未因此下调评级，仍然维持Buy。其核心逻辑在于：Switch 2的硬件渗透逻辑并不完全依赖首发阵容的强度。研报引用了PS5的先例——PS5在经历多次涨价后，其渗透速度仍与PS4相当。这意味着，只要管线在中长期能够持续扩充，硬件销量不会因首发阵容相对平淡而崩塌。

不过，这里存在一个尚未回答的问题：Switch 2的价格上涨是否会影响其在新兴市场的竞争力？研报承认“it may take time”，股价需要等到实际销售数据确认Switch 2 momentum后才能逐步恢复。对于投资者而言，这意味着任天堂的短期催化剂是“sell-through数据”，而非“发布会预期”。

3. Capcom是发行商中管线深度与估值性价比最突出的标的

在发行商板块，GS对Capcom的推荐最为明确。理由来自三个维度。

第一，管线深度超预期。除了市场已有预期的《Monster Hunter Wilds》大型DLC，Capcom还宣布了《Resident Evi

[... middle omitted ...]

的组合下，Capcom的性价比明显优于同行。

第三，IP运营效率。这一点虽然研报没有直接展开，但从管线安排可以看出：Capcom没有过度依赖单一IP，而是形成了《怪物猎人》《生化危机》《鬼武者》《洛克人》等多IP的轮动节奏。这种“IP组合管理”能力，在行业成本上升、单款产品开发风险加大的背景下，是一种结构性优势。

4. Square Enix的中长期叙事改善，但估值与短期执行风险仍不支持买入

Square Enix是这份报告中评级与叙事之间“偏离感”最明显的公司。从基本面看，Square Enix的管线进展相当积极：《最终幻想VII 重制版》第三部定档2027年春季；《勇者斗恶龙XII》虽然未定发售日，但发布了游戏概念大幅调整的实机画面；多款新作也在稳步推进。研报甚至明确表示“中期印象是正面的”。

估值方面，Square Enix的FY2 P/E为26倍，远高于全球游戏板块中位数的16倍。在当前的利率环境和行业增长放缓背景下，这一估值溢价缺乏足够的基本面支撑。

\subsection{[R050] JPM：MLCC材料端正在经历一场被低估的定价权迁移}
{\small\color{refgray}归类：半导体：欧洲设备需求周期尚未见顶与存储供给纪律重塑}

市场对MLCC（多层陶瓷电容器）的关注，长期集中在下游被动元件厂商的库存周期和产能利用率上。但一份来自JPM与国瓷材料（Sinocera）管理层的最新交流纪要揭示了一个更关键的结构性变化：本轮周期的驱动力不再仅仅是需求的回补，而是高端产品对材料端定价权的系统性重塑。

这份报告的核心判断是：MLCC 粉体行业已在2026年上半年进入一个强于市场预期的上升周期，其可持续性并非来自传统消费电子的温和复苏，而是来自AI服务器和汽车电子对高端粉体需求的刚性增长，以及中国对稀土出口管控带来的长期竞争壁垒。这意味着，材料端正在从“被动跟随”转向“主动定价”，而市场对这一转变的定价可能尚不充分。

国瓷材料作为占据国内MLCC粉体市场超过80\%份额、全球约20\%份额的绝对龙头，其产能规划、产品结构变化和客户认证进展，为理解这一产业趋势提供了清晰的微观注脚。以下是我们从这场对话中提炼出的五个关键洞察。

1. 高端产品线正在制造一个利润“断层”，而非简单的量价齐升

传统MLCC粉体与高端产品之间的利润鸿沟，远比市场认知的更为陡峭。报告数据显示，常规MLCC粉体价格在每吨6-7万元人民币，毛利率在33\%以上。而汽车级产品价格跃升至每吨8-9万元，AI服务器级粉体更是突破每吨10万元。高端产品的毛利率区间高达45\%至50\%。

这不仅仅是数字上的差异。它意味着，当一家企业能够将产能从常规产品向高端产品转移时，其整体盈利能力的提升是倍数级的，而非线性增长。国瓷材料在2025年底投产的2000吨高端生产线，以及计划在2026年底前完成的另外3000吨扩产，正是对这一利润断层的直接回应。公司预计2026年高端产品销量将超过1000吨，这将成为其全年营收增长超过20\%目标的核心支撑。

对于投资者而言，需要关注的不是MLCC粉体的总出货量，而是高端出货占比的变化。每一吨从常规向高端的切换，都意味着利润池的显著扩大。

2. 三星的认证通过是国产材料进入全球高端供应链的“压力测试”节点

报告中最具标志性的信号之一，是国瓷材料的AI和汽车级产品已通过三星的认证，并将在2026年第三季度进入大规模量产。这一认证周期耗时近一年，恰恰说明了高端MLCC粉体的技术门槛之高。

技术细节揭示了认证的苛刻之处：高端MLCC粉体对粒径的要求极为严格，国内AI客户要求粒径范围在100-150纳米，三星的要求为200纳米，同时对粉体的均匀性和分散性提出了更高要求。这直接导致了高端粉体的良率可能仅为80\%，远低于常规产品超过95\%的良率。

这意味着，能够通过此类认证的供应商，已经建立起了相当宽的技术护城河。对于已经通过认证的产品，后续客户认证周期将大幅缩短。这不仅为国瓷材料打开了三星这一全球核心客户的大门，更重要的是，它向全球其他MLCC制造商传递了一个信号：中国的高端MLCC粉体供应商已经具备了与国际对手同台竞技的能力。这可能会加速全球供应链的本地化和多元化布局。

3. 中国稀土出口管控正在创造一种“不对称竞争优势”，其影响尚未被完全定价

报告揭示了一个容易被忽视的长期变量：中国供应了全球80\%至95\%的MLCC粉体生产所必需的稀土资源。尽管日本厂商目前持有1-2年的战略性库存，短期内不会出现供应压力，但中国持续的严格稀土出口管制，将在长期内对海外产能形成结构性约束。

这是一项“不对称竞争优势”。海外竞争对手不仅面临原材料获取的不确定性，还面临更高的成本结构。而国内供应商，如国瓷材料，则处于一个相对稳定且成本可控的原材料环境中。这种优势不会在季度财报中立即体现，但它会随着时间的推移，逐渐转化为持续的订单转移和市场份额提升。

对于产业决策者来说，这意味着在评估MLCC粉体供应链的长期稳定性时，不能仅看当前的产能和价格，必须将地缘政治因素纳入核心考量

[... middle omitted ...]

。通过选择那些技术壁垒高、海外垄断强、国内替代需求迫切的方向进行布局，公司能够以较低的市场教育成本，获取较高的成长确定性。报告提到，其球形硅微粉、低轨卫星陶瓷封装外壳、陶瓷基板等新材料项目均在有序推进，部分将在2026年下半年或2027年贡献收入。

对于读者而言，这提供了一个观察框架：不要孤立地看待国瓷材料的每一个新业务，而应将其视为一个基于“国产替代”逻辑、由技术平台驱动的系统性扩张。其研发投入占营收约7\%，且关键项目不设预算上限，这种持续投入是支撑其多线并进的基础。

5. 报告尚未完全回答的关键问题：高端良率爬坡的节奏与AI需求的可持续性

尽管报告提供了大量乐观信号，但仍有几个关键问题悬而未决，这些也正是我们后续需要持续跟踪的焦点。

第一，高端粉体80\%的良率，是目前的技术瓶颈还是常态？如果是瓶颈，随着产能爬坡，良率提升的速度将直接决定实际出货量能否达到1000吨的目标，以及毛利率能否稳定在45\%-50\%的区间。报告并未详细展开良率提升的具体路径和时间表。

\section{Disclaimer}
{\scriptsize\color{refgray}
This community is a paid learning and information-sharing community created for general educational purposes only. The membership fee is charged solely for access to community discussions, educational content, organization, and operational costs. It is not an advisory fee, management fee, performance fee, brokerage fee, or compensation for personalized financial, investment, legal, tax, or professional advice.\par
I participate and share information solely as an individual and for educational discussion among community members. I am not acting as a financial adviser, investment adviser, broker, dealer, portfolio manager, legal adviser, tax adviser, accountant, fiduciary, or any other licensed professional adviser. Nothing shared in this community constitutes, and should not be interpreted as, financial advice, investment advice, legal advice, tax advice, a recommendation, solicitation, offer, endorsement, instruction, or guarantee to buy, sell, hold, short, trade, or invest in any security, cryptocurrency, fund, derivative, strategy, or other financial product.\par
All content shared in this community, including messages, discussions, links, files, charts, examples, market commentary, watchlists, AI-generated or AI-polished summaries, and third-party materials, is general, impersonal, and educational in nature. It does not take into account any individual's financial situation, investment objectives, risk tolerance, investment horizon, liquidity needs, tax status, legal restrictions, or suitability requirements. No content should be relied upon as suitable for any specific person, account, portfolio, or financial decision.\par
Information shared in this community may be incomplete, inaccurate, outdated, speculative, AI-generated, AI-edited, or based on third-party internet sources that have not been independently verified. I make no representation or warranty as to the accuracy, completeness, timeliness, reliability, or suitability of any information shared.\par
Financial markets involve substantial risk, including the possible loss of principal. Past performance, historical data, hypothetical examples, simulated results, backtests, personal opinions, or educational examples do not guarantee future results. Each member is solely responsible for conducting their own research, exercising independent judgment, and making their own decisions. Before making any financial, investment, legal, or tax decision, members should consult qualified and licensed professionals.\par
Participation in this community, payment of any membership fee, or communication with me or other members does not create any adviser-client, fiduciary, professional, contractual, agency, partnership, employment, or other advisory relationship. By joining or participating in this community, you acknowledge and agree that you are solely responsible for your own decisions, actions, transactions, profits, losses, risks, and consequences, and that I shall not be liable for any direct, indirect, incidental, consequential, financial, legal, tax, or other loss or damage arising from or related to any content, discussion, information, or material shared in this community.\par
}
\end{document}
